%============================================================
%  Bd. 08 - Axiomatischer Aufbau der endlichen Mengen
%============================================================

\documentclass[main.tex]{subfiles}

\ifSubfilesClassLoaded{
  \BandSetupStandalone{B08}
}{
}

\title{Bd. 08 - Axiomatischer Aufbau der endlichen Mengen}
\author{Martin Kunze}
\date{}
\setcounter{file}{8}

\begin{document}

\title{Bd. 08 - Axiomatischer Aufbau der endlichen Mengen}
\author{Martin Kunze}
\date{}
\hypersetup{pageanchor=false}
\maketitle
\hypersetup{pageanchor=true}
\setcounter{tocdepth}{3}
\tableofcontents
\setcounter{file}{8}
\renewcommand{\FormulaBandID}{8}

\chapter{Einleitung}

Band~03 führt das primitive Endlichkeitsprädikat und seine drei Axiome ein.
Dieser Band entwickelt deren Konsequenzen, verbindet sie mit den in Band~07
aufgebauten natürlichen Zahlen und untersucht die verschiedenen
Charakterisierungen endlicher Mengen.

\chapter{Endlichkeitsaxiome}

Band~03 führt das Prädikat \(\Finite{\cdot}\) als primitives Symbol ein.
Die folgenden drei dort registrierten Endlichkeitsaxiome bilden den
Ausgangspunkt dieses Bandes. Zuerst stehen sie gesammelt; anschließend
werden sie einzeln aufgeführt. Die Registrierungen aus Band~03 bleiben dabei
unverändert, damit sämtliche Verweise im Folgenden auf dieselben Axiome
zeigen.

\FormulaDefDeltaK[Endlichkeitsaxiome]{%
  \Finite{\,\cdot\,}%
}{Endlichkeitsaxiome}{%
  \DeltaRow{\textbf{Axiome}}{}
  \DeltaRow{Leere Menge}
           {\Finite{\varnothing}}
           [\FormulaRefAuto{FiniteEmptyAxiom}[ax]]
  \DeltaRow{\makecell[l]{Adjunktion eines\\ neuen Elements}}
           {a\notin A\dsep \Finite A\vdash \Finite{A\cup\{a\}}}
           [\FormulaRefAuto{FiniteAdjunctionAxiom}[ax]]
  \DeltaRow{\makecell[l]{Induktionsregel f\"ur\\ endliche Mengen}}
           {\begin{aligned}[t]
              &P(\varnothing)\dsep
              [\Finite A,P(A),a\notin A]\vdots{}\\
              &P(A\cup\{a\})\dsep \Finite M\vdash P(M)
            \end{aligned}}
           [\FormulaRefAuto{FiniteInductionRule}[ax]]
  \DeltaRow{\textbf{Symbol aus Band~03}}{}
  \DeltaRow{Endlichkeit}{\Finite{\,\cdot\,}}
}

\begin{samepage}
  \begin{formulaAx}[Die leere Menge ist endlich]
    \[
      \Finite{\varnothing}
    \]
  \end{formulaAx}
  \par\noindent\emph{Registriert in Band~03:}
  \FormulaRefAuto{FiniteEmptyAxiom}[ax].\par
\end{samepage}

\begin{samepage}
  \begin{formulaAx}[Adjunktion eines neuen Elements]
    \[
      a\notin A\dsep \Finite A\vdash \Finite{A\cup\{a\}}
    \]
  \end{formulaAx}
  \par\noindent\emph{Registriert in Band~03:}
  \FormulaRefAuto{FiniteAdjunctionAxiom}[ax].\par
\end{samepage}

\begin{samepage}
  \begin{formulaAx}[Induktionsregel f\"ur endliche Mengen]
    \[
      P(\varnothing)\dsep
      [\Finite A,P(A),a\notin A]\vdots P(A\cup\{a\})\dsep
      \Finite M\vdash P(M)
    \]
  \end{formulaAx}
  \par\noindent\emph{Registriert in Band~03:}
  \FormulaRefAuto{FiniteInductionRule}[ax].\par
\end{samepage}

Neue Endlichkeitsaxiome werden in diesem Band nicht deklariert. Insbesondere
sind Transport, Vereinigung, Teilmenge, Bild, Potenzmenge sowie die
Dedekind- und Tarski-Charakterisierungen hier zu beweisende Folgerungen.

Die Bijektionscharakterisierung wird in diesem Band neu und ausschließlich mit
den in Band~07 entwickelten Peano-Anfangsabschnitten \(\NatSeg{n}\)
hergeleitet. Aus Band~03 werden dafür keine endlichkeitstheoretischen
Folgerungen übernommen.
\chapter{Endliche Mengen}

\section{Unmittelbare Folgerungen der Axiome}

Zunächst halten wir die beiden elementaren Axiome als Sätze fest. Dadurch
können die folgenden Beweise einheitlich auf registrierte Theorem-IDs
verweisen, ohne zusätzliche Annahmen einzuführen.

\FormulaThmDeltaK[Die leere Menge ist endlich]{%
  \Finite{\varnothing}%
}{FiniteEmpty}{%
  \DeltaRow{Mengen}{}%
}
\begin{tabproof}
  \proofstep{}{\Finite{\varnothing}}{\FormulaRefAuto{FiniteEmptyAxiom}[ax]}
\end{tabproof}

\FormulaThmDeltaK[Adjunktion eines neuen Elements erhält Endlichkeit]{%
  a\notin A\dsep \Finite A\vdash \Finite{A\cup\{a\}}%
}{FiniteAdjunction}{%
  \DeltaRow{Mengen}{A\dsep a}%
}
\begin{tabproof}
  \proofstep{1}{a\notin A}{\rA}
  \proofstep{2}{\Finite A}{\rA}
  \proofstep{1,2}{\Finite{A\cup\{a\}}}{%
    \FormulaRefAuto{FiniteAdjunctionAxiom}[ax]{1,2}}
\end{tabproof}

\FormulaThmDeltaK[Adjunktion eines beliebigen Elements]{%
  \Finite A\vdash \Finite{A\cup\{a\}}%
}{FiniteInsert}{%
  \DeltaRow{Mengen}{A\dsep a}%
}
\begin{tabproof}
  \proofstep{1}{\Finite A}{\rA}
  \proofstep{}{a\in A\lor a\notin A}{\FormulaRefAuto{P\lor\neg P}}
  \proofstep{3}{a\in A}{\rA}
  \proofstep{3}{A\cup\{a\}=A}{%
    \FormulaRefAuto{A\cup\{a\}=A\eqvdash a\in A}{3}}
  \proofstep{1,3}{\Finite{A\cup\{a\}}}{\rIE{\FormulaRefAuto{a=b\vdash b=a}{4},1}}
  \proofstep{6}{a\notin A}{\rA}
  \proofstep{1,6}{\Finite{A\cup\{a\}}}{%
    \FormulaRefAuto{FiniteAdjunction}[thm]{6,1}}
  \proofstep{1}{\Finite{A\cup\{a\}}}{\rOE{2,3,5,6,7}}
\end{tabproof}

\FormulaThmDeltaK[Bilder endlicher Teilmengen sind endlich]{%
  C\subseteq A\dsep F\colon A\to B\dsep \Finite C
  \vdash \Finite{F[C]}%
}{FiniteImage}{%
  \DeltaRow{Mengen}{A\dsep B\dsep C\dsep c}%
  \DeltaRow{Funktionen}{F}%
}
Für festes \(F\colon A\to B\) wenden wir die Endlichkeitsinduktion auf
\[
  P(C)\;\coloneqq\;\bigl(C\subseteq A\rightarrow\Finite{F[C]}\bigr)
\]
an. Die Mengen \(A,B\) und die Funktion \(F\) sind dabei lediglich feste
Parameter der Eigenschaft \(P\).
\begin{tabproofsplitwide}
  \proofpartwideindR[Induktionsanfang]{%
    F\colon A\to B\vdash
    \bigl(\varnothing\subseteq A\rightarrow\Finite{F[\varnothing]}\bigr)%
  }{%
    F\colon A\to B\vdash
    \bigl(\varnothing\subseteq A\rightarrow\Finite{F[\varnothing]}\bigr)%
  }[FiniteImageBase]
    \proofstepwidestar[1]{F\colon A\to B}{\rA}
    \proofstepwidestar[2]{\varnothing\subseteq A}{\rA}
    \proofstepwidestar[1]{F[\varnothing]=\varnothing}{%
      \FormulaRefAuto{F\colon A\to B\vdash F[\varnothing]=\varnothing}{1}}
    \proofstepwidestar[]{\Finite{\varnothing}}{\FormulaRefAuto{FiniteEmpty}[thm]}
    \proofstepwidestar[1]{\Finite{F[\varnothing]}}{%
      \rIE{\FormulaRefAuto{a=b\vdash b=a}{3},4}}
    \proofstepwidestar[1]{\varnothing\subseteq A\rightarrow\Finite{F[\varnothing]}}{%
      \rRI{2,5}}
  \closeproofpartwideindR

  \proofpartwideindR[Induktionsschritt]{%
      \begin{aligned}[t]
      &F\colon A\to B\dsep \Finite C\dsep
       (C\subseteq A\rightarrow\Finite{F[C]})\dsep c\notin C\vdash{}\\
      &(C\cup\{c\}\subseteq A\rightarrow\Finite{F[C\cup\{c\}]})
    \end{aligned}%
  }{%
    F\colon A\to B\dsep \Finite C\dsep
    (C\subseteq A\rightarrow\Finite{F[C]})\dsep c\notin C
    \vdash (C\cup\{c\}\subseteq A\rightarrow\Finite{F[C\cup\{c\}]})%
  }[FiniteImageStep]
    \proofstepwidestar[1]{F\colon A\to B}{\rA}
    \proofstepwidestar[2]{\Finite C}{\rA}
    \proofstepwidestar[3]{C\subseteq A\rightarrow\Finite{F[C]}}{\rA}
    \proofstepwidestar[4]{c\notin C}{\rA}
    \proofstepwidestar[5]{C\cup\{c\}\subseteq A}{\rA}
    \proofstepwidestar[5]{C\subseteq A}{%
      \FormulaRefAuto{A\subseteq B\dsep B\subseteq C\vdash A\subseteq C}{%
        \FormulaRefAuto{A\subseteq A\cup B},5}}
    \proofstepwidestar[5]{c\in A}{%
      \FormulaRefAuto{A\subseteq B,\,x\in A\vdash x\in B}{5,\FormulaRefAuto{a\in A\cup\{a\}}}}
    \proofstepwidestar[3,5]{\Finite{F[C]}}{\rRE{3,6}}
    \proofstepwidestar[3,5]{\Finite{F[C]\cup\{F(c)\}}}{%
      \FormulaRefAuto{FiniteInsert}[thm]{8}}
    \proofstepwidestar[1,5]{F[C\cup\{c\}]=F[C]\cup\{F(c)\}}{%
      \rIE{%
        \FormulaRefAuto{x\in A\dsep F\colon A\to B\vdash F[\{x\}]=\{F(x)\}}{7,1},
        \FormulaRefAuto{F\colon M\to N\dsep A\subseteq M\dsep B\subseteq M
          \vdash F[A\cup B]=F[A]\cup F[B]}{%
          1,6,\FormulaRefAuto{a\in A\vdash\{a\}\subseteq A}{7}}}}
    \proofstepwidestar[1,3,5]{\Finite{F[C\cup\{c\}]}}{%
      \rIE{\FormulaRefAuto{a=b\vdash b=a}{10},9}}
    \proofstepwidestar[1,3]{%
      C\cup\{c\}\subseteq A\rightarrow\Finite{F[C\cup\{c\}]}}{\rRI{5,11}}
  \closeproofpartwideindR

  \proofpartwide{Induktionsschluss}
    \proofstepwidestar[1]{C\subseteq A}{\rA}
    \proofstepwidestar[2]{F\colon A\to B}{\rA}
    \proofstepwidestar[3]{\Finite C}{\rA}
    \proofstepwidestar[2,3]{C\subseteq A\rightarrow\Finite{F[C]}}{%
      \FormulaRefAuto{FiniteInductionRule}[ax]{IA,IS,3}}
    \proofstepwidestar[1,2,3]{\Finite{F[C]}}{\rRE{4,1}}
  \closeproofpartwide
\end{tabproofsplitwide}

\FormulaThmDeltaK[Transport der Endlichkeit mittels Bijektion]{%
  \Finite M\dsep F\colon M\bij N\vdash \Finite N%
}{FiniteBijectionTransport}{%
  \DeltaRow{Mengen}{M\dsep N}%
  \DeltaRow{Funktionen}{F}%
}
\begin{tabproof}
  \proofstep{1}{\Finite M}{\rA}
  \proofstep{2}{F\colon M\bij N}{\rA}
  \proofstep{2}{F\colon M\to N}{\FormulaRefAuto{Bijektive Funktion}{2}}
  \proofstep{1,2}{\Finite{F[M]}}{%
    \FormulaRefAuto{FiniteImage}[thm]{\FormulaRefAuto{A\subseteq A},3,1}}
  \proofstep{2}{F[M]=N}{\FormulaRefAuto{F\colon A\bij B\vdash F[A]=B}{2}}
  \proofstep{1,2}{\Finite N}{\rIE{5,4}}
\end{tabproof}

\FormulaThmDeltaK[Transport der Endlichkeit entlang der Gleichmächtigkeit]{%
  \Finite M\dsep M\EqCard N\vdash \Finite N%
}{FiniteEqCardTransport}{%
  \DeltaRow{Mengen}{M\dsep N}%
  \DeltaRow{Funktionen}{F}%
}
\begin{tabproof}
  \proofstep{1}{\Finite M}{\rA}
  \proofstep{2}{M\EqCard N}{\rA}
  \proofstep{2}{\exists F\colon M\bij N}{%
    \FormulaRefAuto{A\EqCard B\coloneqq \exists F\colon A\bij B}{2}}
  \proofstep{4}{F\colon M\bij N}{\rA}
  \proofstep{1,4}{\Finite N}{\FormulaRefAuto{FiniteBijectionTransport}[thm]{1,4}}
  \proofstep{1,2}{\Finite N}{\rEE{3,4,5}}
\end{tabproof}


% Die reinen NatSeg-Grundlagen stehen nun fachlich in B07.
% Der bisherige B05-Block bleibt bis zur abschliessenden Abnahme als
% auskommentierte Ausgangsfassung erhalten.
\iffalse
\section{Peano-Anfangsabschnitte}

Die endlichen Anfangsbereiche werden in diesem Band nicht mit natuerlichen
Zahlen als Mengen identifiziert. Fuer eine natuerliche Zahl \(n\) bezeichnet
\(\NatSeg{n}\) den durch die Peano-Ordnung bestimmten Anfangsabschnitt.

\FormulaDefDeltaK[Peano-Anfangsabschnitt]{%
  n\in\mathbb{N}\vdash
  \NatSeg{n}\coloneqq \{i\in\mathbb{N}\mid i<n\}%
}{PeanoNatSeg}{%
  \DeltaRow{Mengen}{i\dsep n}%
  \DeltaRow{Definitionsprinzip}{Peano-Ordnung}%
  \DeltaRow{Neue Symbole}{\NatSeg{n}}%
}

\FormulaThmDeltaK[Charakterisierung des Peano-Anfangsabschnitts]{%
  n\in\mathbb{N}\dsep i\in\mathbb{N}\vdash
  i\in\NatSeg{n}\leftrightarrow i<n%
}{PeanoNatSegMembership}{%
  \DeltaRow{Mengen}{i\dsep n}%
}
\begin{tabproof}
  \proofstep{1}{n\in\mathbb{N}}{\rA}
  \proofstep{2}{i\in\mathbb{N}}{\rA}
  \proofstep{1}{\NatSeg{n}=\{i\in\mathbb{N}\mid i<n\}}{%
    \FormulaRefAuto{PeanoNatSeg}{1}}
  \proofstep{1,2}{i\in\NatSeg{n}\leftrightarrow i<n}{%
    \FormulaRefAuto{x\in\{u\in A\mid P(u)\}\eqvdash x\in A\land P(x)}{3,2}}
\end{tabproof}

\FormulaThmDeltaK[Anfangsabschnitte sind Teilmengen von \(\mathbb{N}\)]{%
  n\in\mathbb{N}\vdash \NatSeg{n}\subseteq\mathbb{N}%
}{PeanoNatSegSubsetN}{%
  \DeltaRow{Mengen}{n}%
}
\begin{tabproof}
  \proofstep{1}{n\in\mathbb{N}}{\rA}
  \proofstep{1}{\NatSeg{n}=\{i\in\mathbb{N}\mid i<n\}}{\FormulaRefAuto{PeanoNatSeg}{1}}
  \proofstep{1}{\NatSeg{n}\subseteq\mathbb{N}}{%
    \FormulaRefAuto{A=\{x\in B\mid P(x)\}\vdash A\subseteq B}{2}}
\end{tabproof}

\FormulaThmDeltaK[Leere Separation bei leerem Praedikat]{%
  A=\{x\in B\mid P(x)\}\dsep
  \forall x\in B\,\neg P(x)\vdash A=\varnothing%
}{EmptySeparationByNegPredicate}{%
  \DeltaRow{Mengen}{A\dsep B\dsep x}%
  \DeltaRow{Praedikat}{P}%
}
\begin{tabproof}
  \proofstep{1}{A=\{x\in B\mid P(x)\}}{\rA}
  \proofstep{2}{\forall x\in B\,\neg P(x)}{\rA}
  \proofstep{3}{x\in A}{\rA}
  \proofstep{1,3}{x\in \{x\in B\mid P(x)\}}{\rIE{1,3}}
  \proofstep{1,3}{x\in B\land P(x)}{%
    \FormulaRefAuto{x\in\{u\in A\mid P(u)\}\eqvdash x\in A\land P(x)}{4}}
  \proofstep{1,3}{x\in B}{\rAEa{5}}
  \proofstep{1,3}{P(x)}{\rAEb{5}}
  \proofstep{1,2,3}{\neg P(x)}{\rRE{\rUE{2},6}}
  \proofstep{1,2,3}{\bot}{\rBI{7,8}}
  \proofstep{1,2}{x\notin A}{\rCI{3,9}}
  \proofstep{1,2}{\forall x\,x\notin A}{\rUI{\rRI{3,10}}}
  \proofstep{1,2}{A=\varnothing}{%
    \FormulaRefAuto{\forall x\,x\notin A\vdash A=\varnothing}{11}}
\end{tabproof}

\FormulaThmDeltaK[Null-Anfangsabschnitt]{%
  \NatSeg{0}=\varnothing%
}{PeanoNatSegZero}{%
  \DeltaRow{Mengen}{}%
}
\begin{tabproof}
  \proofstep{}{0\in\mathbb{N}}{\FormulaRefAuto{PeanoZero}[ax]}
  \proofstep{}{\NatSeg{0}=\{i\in\mathbb{N}\mid i<0\}}{%
    \FormulaRefAuto{PeanoNatSeg}{1}}
  \proofstep{}{\NatSeg{0}=\varnothing}{%
    \FormulaRefAuto{EmptySeparationByNegPredicate}{2}}
\end{tabproof}

\FormulaThmDeltaK[Nachfolgerfall der strikten Ordnung]{%
  n\in\mathbb{N}\dsep i\in\mathbb{N}\vdash
  i<\Succ(n)\leftrightarrow i<n\lor i=n%
}{PeanoStrictSuccSplit}{%
  \DeltaRow{Mengen}{i\dsep n}%
  \DeltaRow{Relationen}{<_{\mathbb{N}}}%
}
\begin{tabproof}
  \proofstep{1}{n\in\mathbb{N}}{\rA}
  \proofstep{2}{i\in\mathbb{N}}{\rA}
  \proofstep{1}{\Succ(n)\in\mathbb{N}}{\FormulaRefAuto{PeanoSuccClosure}{1}}
  \proofstep{1,2}{i<\Succ(n)\leftrightarrow i\in\Succ(n)}{%
    \FormulaRefAuto{VonNeumannStrictMembership}[thm]{2,3}}
  \proofstep{1,2}{i\in\Succ(n)\leftrightarrow i\leq n}{%
    \FormulaRefAuto{VonNeumannSuccMembershipLeqEquiv}[thm]{2,1}}
  \proofstep{1,2}{i<\Succ(n)\leftrightarrow i\leq n}{%
    \FormulaRefAuto{P\leftrightarrow Q\dsep Q\leftrightarrow R\vdash P\leftrightarrow R}{4,5}}
  \proofstep{7}{i\leq n}{\rA}
  \proofstep{1,2,7}{i<n\lor i=n}{%
    \FormulaRefAuto{VonNeumannLeqSplit}[thm]{2,1,7}}
  \proofstep{1,2}{i\leq n\rightarrow i<n\lor i=n}{\rRI{7,8}}
  \proofstep{10}{i<n\lor i=n}{\rA}
  \proofstep{11}{i<n}{\rA}
  \proofstep{1,2,11}{i\leq n}{%
    \FormulaRefAuto{VonNeumannStrictImpliesLeq}[thm]{2,1,11}}
  \proofstep{13}{i=n}{\rA}
  \proofstep{1}{n\leq n}{\FormulaRefAuto{VonNeumannLeqReflexive}[thm]{1}}
  \proofstep{1,13}{i\leq n}{\rIE{13,14}}
  \proofstep{1,2,10}{i\leq n}{\rOE{10,11,12,13,15}}
  \proofstep{1,2}{(i<n\lor i=n)\rightarrow i\leq n}{\rRI{10,16}}
  \proofstep{1,2}{i\leq n\leftrightarrow i<n\lor i=n}{\rLRI{9,17}}
  \proofstep{1,2}{i<\Succ(n)\leftrightarrow i<n\lor i=n}{%
    \FormulaRefAuto{P\leftrightarrow Q\dsep Q\leftrightarrow R\vdash P\leftrightarrow R}{6,18}}
\end{tabproof}

\FormulaThmDeltaK[Nachfolger-Anfangsabschnitt]{%
  n\in\mathbb{N}\vdash
  \NatSeg{\Succ(n)}=\NatSeg{n}\cup\{n\}%
}{PeanoNatSegSucc}{%
  \DeltaRow{Mengen}{n}%
}
\begin{tabproof}
  \proofstep{1}{n\in\mathbb{N}}{\rA}
  \proofstep{1}{\Succ(n)\in\mathbb{N}}{\FormulaRefAuto{PeanoSuccClosure}{1}}
  \proofstep{1}{\NatSeg{\Succ(n)}=\{i\in\mathbb{N}\mid i<\Succ(n)\}}{%
    \FormulaRefAuto{PeanoNatSeg}{2}}
  \proofstep{1}{\NatSeg{n}=\{i\in\mathbb{N}\mid i<n\}}{%
    \FormulaRefAuto{PeanoNatSeg}{1}}
  \proofstep{1}{\NatSeg{\Succ(n)}=\NatSeg{n}\cup\{n\}}{%
    \FormulaRefAuto{n\in\mathbb{N}\dsep i\in\mathbb{N}\vdash i<\Succ(n)\leftrightarrow i<n\lor i=n}{1}}
\end{tabproof}
\FormulaThmDeltaK[Elemente eines Nachfolger-Anfangsabschnitts]{%
  n\in\mathbb{N}\dsep x\in\NatSeg{\Succ(n)}\vdash
  x\in\NatSeg{n}\lor x=n%
}{PeanoNatSegSuccElim}{%
  \DeltaRow{Mengen}{n\dsep x}%
}
\begin{tabproof}
  \proofstep{1}{n\in\mathbb{N}}{\rA}
  \proofstep{2}{x\in\NatSeg{\Succ(n)}}{\rA}
  \proofstep{1}{\NatSeg{\Succ(n)}=\NatSeg{n}\cup\{n\}}{%
    \FormulaRefAuto{PeanoNatSegSucc}{1}}
  \proofstep{1,2}{x\in\NatSeg{n}\cup\{n\}}{\rIE{3,2}}
  \proofstep{1,2}{x\in\NatSeg{n}\lor x=n}{%
    \FormulaRefAuto{x\in A\cup\{a\} \vdash x\in A\lor x=a}{4}}
\end{tabproof}

\FormulaThmDeltaK[Einbettung des Anfangsabschnitts in den Nachfolger]{%
  n\in\mathbb{N}\vdash \NatSeg{n}\subseteq\NatSeg{\Succ(n)}%
}{PeanoNatSegSubsetSucc}{%
  \DeltaRow{Mengen}{n}%
}
\begin{tabproof}
  \proofstep{1}{n\in\mathbb{N}}{\rA}
  \proofstep{1}{\NatSeg{\Succ(n)}=\NatSeg{n}\cup\{n\}}{%
    \FormulaRefAuto{PeanoNatSegSucc}{1}}
  \proofstep{}{\NatSeg{n}\subseteq\NatSeg{n}\cup\{n\}}{%
    \FormulaRefAuto{A\subseteq A\cup B}}
  \proofstep{1}{\NatSeg{n}\subseteq\NatSeg{\Succ(n)}}{\rIE{2,3}}
\end{tabproof}

\FormulaThmDeltaK[Strikte Ordnung als Mitgliedschaft im Peano-Anfangsabschnitt]{%
  m\in\mathbb{N}\dsep n\in\mathbb{N}\vdash
  m<n \leftrightarrow m\in \NatSeg{n}%
}{PeanoNatSegStrictMembership}{%
  \DeltaRow{Mengen}{m\dsep n}%
  \DeltaRow{Relationen}{<_{\mathbb{N}}}%
}
\begin{tabproof}
  \proofstep{1}{m\in\mathbb{N}}{\rA}
  \proofstep{2}{n\in\mathbb{N}}{\rA}
  \proofstep{2,1}{m\in\NatSeg{n}\leftrightarrow m<n}{%
    \FormulaRefAuto{PeanoNatSegMembership}{2,1}}
  \proofstep{1,2}{m<n \leftrightarrow m\in\NatSeg{n}}{%
    \FormulaRefAuto{P \leftrightarrow Q \vdash Q \leftrightarrow P}{3}}
\end{tabproof}
\FormulaThmDeltaK[Kein Endpunkt im eigenen Anfangsabschnitt]{%
  n\in\mathbb{N}\vdash n\notin \NatSeg{n}%
}{PeanoNatSegEndpointNotIn}{%
  \DeltaRow{Mengen}{n}%
}
\begin{tabproof}
  \proofstep{1}{n\in\mathbb{N}}{\rA}
  \proofstep{2}{n\in\NatSeg{n}}{\rA}
  \proofstep{1,2}{n<n}{%
    \FormulaRefAuto{PeanoNatSegStrictMembership}{1,1,2}}
  \proofstep{1}{n\leq n}{%
    \FormulaRefAuto{VonNeumannLeqReflexive}[thm]{1}}
  \proofstep{1,2}{\bot}{%
    \FormulaRefAuto{PeanoOrderNoCycle}{1,1,3,4}}
  \proofstep{1}{n\notin\NatSeg{n}}{\rCI{2,5}}
\end{tabproof}

\FormulaThmDeltaK[Endpunkt im Nachfolger-Anfangsabschnitt]{%
  n\in\mathbb{N}\vdash n\in \NatSeg{\Succ(n)}%
}{PeanoNatSegEndpointInSucc}{%
  \DeltaRow{Mengen}{n}%
}
\begin{tabproof}
  \proofstep{1}{n\in\mathbb{N}}{\rA}
  \proofstep{1}{n<\Succ(n)}{%
    \FormulaRefAuto{n\in\mathbb{N}\vdash n<\Succ(n)}{1}}
  \proofstep{1}{\Succ(n)\in\mathbb{N}}{%
    \FormulaRefAuto{PeanoSuccClosure}{1}}
  \proofstep{1}{n\in\NatSeg{\Succ(n)}}{%
    \FormulaRefAuto{PeanoNatSegStrictMembership}{1,3,2}}
\end{tabproof}

\FormulaThmDeltaK[Nichtstrikte Ordnung als strikte Nachfolgerordnung]{%
  m\in\mathbb{N}\dsep n\in\mathbb{N}\vdash
  m\leq n\leftrightarrow m<\Succ(n)%
}{PeanoLeqSuccStrict}{%
  \DeltaRow{Mengen}{m\dsep n}%
  \DeltaRow{Relationen}{\leq_{\mathbb{N}}\dsep <_{\mathbb{N}}}%
}
\begin{tabproof}
  \proofstep{1}{m\in\mathbb{N}}{\rA}
  \proofstep{2}{n\in\mathbb{N}}{\rA}
  \proofstep{2}{\Succ(n)\in\mathbb{N}}{\FormulaRefAuto{PeanoSuccClosure}{2}}
  \proofstep{1,2}{m\in\Succ(n)\leftrightarrow m\leq n}{%
    \FormulaRefAuto{VonNeumannSuccMembershipLeqEquiv}[thm]{1,2}}
  \proofstep{1,2}{m\leq n\leftrightarrow m\in\Succ(n)}{%
    \FormulaRefAuto{P\leftrightarrow Q\vdash Q\leftrightarrow P}{4}}
  \proofstep{1,2}{m<\Succ(n)\leftrightarrow m\in\Succ(n)}{%
    \FormulaRefAuto{VonNeumannStrictMembership}[thm]{1,3}}
  \proofstep{1,2}{m\in\Succ(n)\leftrightarrow m<\Succ(n)}{%
    \FormulaRefAuto{P\leftrightarrow Q\vdash Q\leftrightarrow P}{6}}
  \proofstep{1,2}{m\leq n\leftrightarrow m<\Succ(n)}{%
    \FormulaRefAuto{P\leftrightarrow Q\dsep Q\leftrightarrow R\vdash P\leftrightarrow R}{5,7}}
\end{tabproof}

\FormulaThmDeltaK[Nichtstrikte Ordnung als Nachfolger-Anfangsabschnitt]{%
  m\in\mathbb{N}\dsep n\in\mathbb{N}\vdash
  m\leq n \leftrightarrow m\in \NatSeg{\Succ(n)}%
}{PeanoNatSegLeqSuccMembership}{%
  \DeltaRow{Mengen}{m\dsep n}%
  \DeltaRow{Relationen}{\leq_{\mathbb{N}}}%
}
\begin{tabproof}
  \proofstep{1}{m\in\mathbb{N}}{\rA}
  \proofstep{2}{n\in\mathbb{N}}{\rA}
  \proofstep{2}{\Succ(n)\in\mathbb{N}}{%
    \FormulaRefAuto{PeanoSuccClosure}{2}}
  \proofstep{1,2}{m\leq n\leftrightarrow m<\Succ(n)}{%
    \FormulaRefAuto{PeanoLeqSuccStrict}{1,2}}
  \proofstep{1,2}{m<\Succ(n)\leftrightarrow m\in\NatSeg{\Succ(n)}}{%
    \FormulaRefAuto{PeanoNatSegStrictMembership}{1,3}}
  \proofstep{1,2}{m\leq n \leftrightarrow m\in \NatSeg{\Succ(n)}}{%
    \FormulaRefAuto{P\leftrightarrow Q\dsep Q\leftrightarrow R\vdash P\leftrightarrow R}{4,5}}
\end{tabproof}

\FormulaThmDeltaK[Ordnung und Anfangsabschnittsinklusion]{%
  m\in\mathbb{N}\dsep n\in\mathbb{N}\vdash
  m\leq n \leftrightarrow \NatSeg{m}\subseteq \NatSeg{n}%
}{PeanoNatSegLeqSubset}{%
  \DeltaRow{Mengen}{m\dsep n}%
  \DeltaRow{Relationen}{\leq_{\mathbb{N}}}%
}
\begin{tabproof}
  \proofstep{1}{m\in\mathbb{N}}{\rA}
  \proofstep{2}{n\in\mathbb{N}}{\rA}
  \proofstep{1,2}{m\leq n \leftrightarrow \NatSeg{m}\subseteq \NatSeg{n}}{%
    \FormulaRefAuto{PeanoInductionRule}[ax], \FormulaRefAuto{PeanoNatSegSucc}}
\end{tabproof}

\FormulaThmDeltaK[Strikte Ordnung liefert Anfangsabschnittsinklusion]{%
  m\in\mathbb{N}\dsep n\in\mathbb{N}\dsep m\in \NatSeg{n}
  \vdash \NatSeg{m}\subseteq \NatSeg{n}%
}{PeanoNatSegMembershipSubset}{%
  \DeltaRow{Mengen}{m\dsep n}%
}
\begin{tabproof}
  \proofstep{1}{m\in\mathbb{N}}{\rA}
  \proofstep{2}{n\in\mathbb{N}}{\rA}
  \proofstep{3}{m\in\NatSeg{n}}{\rA}
  \proofstep{1,2,3}{m<n}{%
    \FormulaRefAuto{PeanoNatSegStrictMembership}{1,2,3}}
  \proofstep{1,2,3}{m\leq n}{%
    \FormulaRefAuto{m\in\mathbb{N}\dsep n\in\mathbb{N}\dsep m<n\vdash m\leq n}{1,2,4}}
  \proofstep{1,2,3}{\NatSeg{m}\subseteq\NatSeg{n}}{%
    \FormulaRefAuto{PeanoNatSegLeqSubset}{1,2,5}}
\end{tabproof}

\FormulaThmDeltaK[Teilmengen eines Nachfolger-Anfangsabschnitts ohne Endpunkt]{%
  n\in\mathbb{N}\dsep B\subseteq \NatSeg{\Succ(n)}\dsep n\notin B
  \vdash B\subseteq \NatSeg{n}%
}{PeanoNatSegSubsetSuccWithoutEndpoint}{%
  \DeltaRow{Mengen}{B\dsep n}%
}
\begin{tabproof}
  \proofstep{1}{n\in\mathbb{N}}{\rA}
  \proofstep{2}{B\subseteq \NatSeg{\Succ(n)}}{\rA}
  \proofstep{3}{n\notin B}{\rA}
  \proofstep{4}{x\in B}{\rA}
  \proofstep{2,4}{x\in\NatSeg{\Succ(n)}}{%
    \FormulaRefAuto{A\subseteq B,\, x\in A\vdash x\in B}{2,4}}
  \proofstep{1,2,4}{x\in\NatSeg{n}\lor x=n}{%
    \FormulaRefAuto{PeanoNatSegSuccElim}{1,5}}
  \proofstep{7}{x\in\NatSeg{n}}{\rA}
  \proofstep{8}{x=n}{\rA}
  \proofstep{3,4,8}{\bot}{\rBI{\rIE{8,4},3}}
  \proofstep{3,4,8}{x\in\NatSeg{n}}{\FormulaRefAuto{P,\,\neg P \vdash Q}{4,9}}
  \proofstep{1,2,3,4}{x\in\NatSeg{n}}{\rOE{6,7,7,8,10}}
  \proofstep{1,2,3}{B\subseteq\NatSeg{n}}{%
    \FormulaRefAuto{A \subseteq B \coloneq \forall x\,(x\in A \rightarrow x\in B)}{\rUI{\rRI{4,11}}}}
\end{tabproof}

\FormulaThmDeltaK[Teilmengen eines Nachfolger-Anfangsabschnitts mit Endpunkt]{%
  n\in\mathbb{N}\dsep B\subseteq \NatSeg{\Succ(n)}\dsep n\in B
  \vdash B=(B\cap \NatSeg{n})\cup\{n\}%
}{PeanoNatSegSubsetSuccWithEndpoint}{%
  \DeltaRow{Mengen}{B\dsep n}%
}
\begin{tabproof}
  \proofstep{1}{n\in\mathbb{N}}{\rA}
  \proofstep{2}{B\subseteq \NatSeg{\Succ(n)}}{\rA}
  \proofstep{3}{n\in B}{\rA}
  \proofstep{1,2,3}{B=(B\cap \NatSeg{n})\cup\{n\}}{%
    \FormulaRefAuto{PeanoNatSegSucc}{1}, \FormulaRefAuto{A\subseteq B\dsep B\subseteq A \vdash A=B}}
\end{tabproof}


\fi

\section{Endlichkeit abstrakter Peano-Anfangsabschnitte}

Die folgenden beiden Sätze betreffen ausdrücklich die abstrakten
Peano-Anfangsabschnitte. Sie werden direkt aus den drei Endlichkeitsaxiomen
gewonnen und dienen nicht als allgemeine Endlichkeitszeugen.

\FormulaThmDeltaK[Peano-Anfangsabschnitte sind endlich]{%
  n\in\mathbb{N}\vdash \Finite{\NatSeg{n}}%
}{PeanoNatSegFinite}{%
  \DeltaRow{Mengen}{n}%
}
\begin{tabproofsplitwide}
  \proofpartwideindR[Induktionsanfang]{%
    \Finite{\NatSeg{0}}%
  }{%
    \Finite{\NatSeg{0}}%
  }[PeanoNatSegFiniteBase]
    \proofstepwidestar[]{\NatSeg{0}=\varnothing}{%
      \FormulaRefAuto{PeanoNatSegZero}[thm]}
    \proofstepwidestar[]{\Finite{\varnothing}}{%
      \FormulaRefAuto{FiniteEmpty}[thm]}
    \proofstepwidestar[]{\Finite{\NatSeg{0}}}{%
      \rIE{\FormulaRefAuto{a=b\vdash b=a}{1},2}}
  \closeproofpartwideindR

  \proofpartwideindR[Induktionsschritt]{%
    n\in\mathbb{N}\dsep \Finite{\NatSeg{n}}
    \vdash \Finite{\NatSeg{\Succ(n)}}%
  }{%
    n\in\mathbb{N}\dsep \Finite{\NatSeg{n}}
    \vdash \Finite{\NatSeg{\Succ(n)}}%
  }[PeanoNatSegFiniteStep]
    \proofstepwidestar[1]{n\in\mathbb{N}}{\rA}
    \proofstepwidestar[2]{\Finite{\NatSeg{n}}}{\rA}
    \proofstepwidestar[1]{n\notin\NatSeg{n}}{%
      \FormulaRefAuto{PeanoNatSegEndpointNotIn}[thm]{1}}
    \proofstepwidestar[1,2]{\Finite{\NatSeg{n}\cup\{n\}}}{%
      \FormulaRefAuto{FiniteAdjunction}[thm]{3,2}}
    \proofstepwidestar[1]{%
      \NatSeg{\Succ(n)}=\NatSeg{n}\cup\{n\}}{%
      \FormulaRefAuto{PeanoNatSegSucc}[thm]{1}}
    \proofstepwidestar[1,2]{\Finite{\NatSeg{\Succ(n)}}}{%
      \rIE{\FormulaRefAuto{a=b\vdash b=a}{5},4}}
  \closeproofpartwideindR

  \proofpartwide{Induktionsschluss}
    \proofstepwidestar[1]{n\in\mathbb{N}}{\rA}
    \proofstepwidestar[1]{\Finite{\NatSeg{n}}}{%
      \FormulaRefAuto{PeanoInductionRule}{IA,IS,1}}
  \closeproofpartwide
\end{tabproofsplitwide}

\FormulaThmDeltaK[Einführung der Endlichkeit aus einem Anfangsabschnittszeugen]{%
  n\in\mathbb{N}\dsep \NatSeg{n}\EqCard M\vdash \Finite{M}%
}{FiniteNatSegEqCardIntro}{%
  \DeltaRow{Mengen}{M\dsep n}%
}
\begin{tabproof}
  \proofstep{1}{n\in\mathbb{N}}{\rA}
  \proofstep{2}{\NatSeg{n}\EqCard M}{\rA}
  \proofstep{1}{\Finite{\NatSeg{n}}}{%
    \FormulaRefAuto{PeanoNatSegFinite}[thm]{1}}
  \proofstep{1,2}{\Finite{M}}{%
    \FormulaRefAuto{FiniteEqCardTransport}[thm]{3,2}}
\end{tabproof}

\section{Bijektionscharakterisierung der Endlichkeit}

Die Charakterisierung wird lokal aus den drei Endlichkeitsaxiomen aufgebaut.
Ihre Zählmengen sind durchgehend die Peano-Anfangsabschnitte aus Band~07.

\FormulaThmDeltaK[Anfangsabschnittszeuge für endliche Mengen]{%
  \Finite M\vdash
  \exists n\in\mathbb{N}\,\bigl(\NatSeg{n}\EqCard M\bigr)%
}{FiniteNatSegEqCardWitness}{%
  \DeltaRow{Mengen}{A\dsep M\dsep a\dsep n}%
}
Hier verwenden wir die Endlichkeitsinduktion mit
\[
  P(A)\;\coloneqq\;
  \exists n\in\mathbb{N}\,\bigl(\NatSeg{n}\EqCard A\bigr).
\]
\begin{tabproofsplitwide}
  \proofpartwideindR[Induktionsanfang]{%
    \exists n\in\mathbb{N}\,\bigl(\NatSeg{n}\EqCard\varnothing\bigr)%
  }{%
    \exists n\in\mathbb{N}\,\bigl(\NatSeg{n}\EqCard\varnothing\bigr)%
  }[FiniteNatSegEqCardWitnessBase]
    \proofstepwidestar[]{0\in\mathbb{N}}{\FormulaRefAuto{PeanoZero}[ax]}
    \proofstepwidestar[]{\NatSeg{0}=\varnothing}{%
      \FormulaRefAuto{PeanoNatSegZero}[thm]}
    \proofstepwidestar[]{\varnothing\EqCard\varnothing}{%
      \FormulaRefAuto{A\EqCard B\coloneqq \exists F\colon A\bij B}{%
        \FormulaRefAuto{\varnothing\colon\varnothing\bij\varnothing}}}
    \proofstepwidestar[]{\NatSeg{0}\EqCard\varnothing}{\rIE{2,3}}
    \proofstepwidestar[]{%
      \exists n\in\mathbb{N}\,\bigl(\NatSeg{n}\EqCard\varnothing\bigr)}{%
      \rEI{\rAI{1,4}}}
  \closeproofpartwideindR

  \proofpartwideindR[Induktionsschritt]{%
    \begin{aligned}[t]
      &\Finite A\dsep
       \exists n\in\mathbb{N}\,\bigl(\NatSeg{n}\EqCard A\bigr)
       \dsep a\notin A\vdash{}\\
      &\exists m\in\mathbb{N}\,\bigl(\NatSeg{m}\EqCard A\cup\{a\}\bigr)
    \end{aligned}%
  }{%
    \Finite A\dsep
    \exists n\in\mathbb{N}\,\bigl(\NatSeg{n}\EqCard A\bigr)
    \dsep a\notin A\vdash
    \exists m\in\mathbb{N}\,\bigl(\NatSeg{m}\EqCard A\cup\{a\}\bigr)%
  }[FiniteNatSegEqCardWitnessStep]
    \proofstepwidestar[1]{\Finite A}{\rA}
    \proofstepwidestar[2]{%
      \exists n\in\mathbb{N}\,\bigl(\NatSeg{n}\EqCard A\bigr)}{\rA}
    \proofstepwidestar[3]{a\notin A}{\rA}
    \proofstepwidestar[4]{n\in\mathbb{N}\land\NatSeg{n}\EqCard A}{\rA}
    \proofstepwidestar[4]{n\in\mathbb{N}}{\rAEa{4}}
    \proofstepwidestar[4]{\NatSeg{n}\EqCard A}{\rAEb{4}}
    \proofstepwidestar[4]{n\notin\NatSeg{n}}{%
      \FormulaRefAuto{PeanoNatSegEndpointNotIn}[thm]{5}}
    \proofstepwidestar[3,4]{%
      \NatSeg{n}\cup\{n\}\EqCard A\cup\{a\}}{%
      \FormulaRefAuto{EqCardFreshAdjunction}[thm]{6,7,3}}
    \proofstepwidestar[4]{\NatSeg{\Succ(n)}=\NatSeg{n}\cup\{n\}}{%
      \FormulaRefAuto{PeanoNatSegSucc}[thm]{5}}
    \proofstepwidestar[3,4]{%
      \NatSeg{\Succ(n)}\EqCard A\cup\{a\}}{\rIE{9,8}}
    \proofstepwidestar[4]{\Succ(n)\in\mathbb{N}}{%
      \FormulaRefAuto{PeanoSuccClosure}{5}}
    \proofstepwidestar[3,4]{%
      \exists m\in\mathbb{N}\,\bigl(\NatSeg{m}\EqCard A\cup\{a\}\bigr)}{%
      \rEI{\rAI{11,10}}}
    \proofstepwidestar[2,3]{%
      \exists m\in\mathbb{N}\,\bigl(\NatSeg{m}\EqCard A\cup\{a\}\bigr)}{%
      \rEE{2,4,12}}
  \closeproofpartwideindR

  \proofpartwide{Induktionsschluss}
    \proofstepwidestar[1]{\Finite M}{\rA}
    \proofstepwidestar[1]{%
      \exists n\in\mathbb{N}\,\bigl(\NatSeg{n}\EqCard M\bigr)}{%
      \FormulaRefAuto{FiniteInductionRule}[ax]{IA,IS,1}}
  \closeproofpartwide
\end{tabproofsplitwide}

\FormulaThmDeltaK[Charakterisierung durch Peano-Anfangsabschnitte]{%
  \Finite M\eqvdash
  \exists n\in\mathbb{N}\,\bigl(\NatSeg{n}\EqCard M\bigr)%
}{FiniteEqCardEquiv}{%
  \DeltaRow{Mengen}{M\dsep n}%
}
\begin{tabproofsplit}
  \proofpart{\(\vdash\)}
    \proofstep{1}{\Finite M}{\rA}
    \proofstep{1}{\exists n\in\mathbb{N}\,(\NatSeg{n}\EqCard M)}{%
      \FormulaRefAuto{FiniteNatSegEqCardWitness}[thm]{1}}
  \closeproofpart

  \proofpart{\(\dashv\)}
    \proofstep{1}{\exists n\in\mathbb{N}\,(\NatSeg{n}\EqCard M)}{\rA}
    \proofstep{2}{n\in\mathbb{N}\land\NatSeg{n}\EqCard M}{\rA}
    \proofstep{2}{n\in\mathbb{N}}{\rAEa{2}}
    \proofstep{2}{\NatSeg{n}\EqCard M}{\rAEb{2}}
    \proofstep{2}{\Finite M}{%
      \FormulaRefAuto{FiniteNatSegEqCardIntro}[thm]{3,4}}
    \proofstep{1}{\Finite M}{\rEE{1,2,5}}
  \closeproofpart
\end{tabproofsplit}

\section{Elementare Eigenschaften endlicher Mengen}

Die grundlegenden Bild- und Transportaussagen wurden oben direkt aus den drei
Endlichkeitsaxiomen hergeleitet; die Bijektionscharakterisierung wird dafür
nicht vorausgesetzt.

\FormulaThmDeltaK[Vereinigung endlicher Mengen ist endlich]{%
  \Finite A\dsep \Finite B\vdash \Finite{A\cup B}%
}{FiniteUnion}{%
  \DeltaRow{Mengen}{A\dsep B\dsep b}%
}
Für festes endliches \(A\) induzieren wir ausschließlich über \(B\) mit
\[
  P(B)\;\coloneqq\;\Finite{A\cup B}.
\]
Im Induktionsschritt ist \(b\notin B\) die vom Axiom verlangte
Frischebedingung. Falls dennoch \(b\in A\) gilt, deckt
\FormulaRefAuto{FiniteInsert}[thm] diesen Fall ab.
\begin{tabproofsplitwide}
  \proofpartwideindR[Induktionsanfang]{%
    \Finite A\vdash \Finite{A\cup\varnothing}%
  }{%
    \Finite A\vdash \Finite{A\cup\varnothing}%
  }[FiniteUnionBase]
    \proofstepwidestar[1]{\Finite A}{\rA}
    \proofstepwidestar[]{A\cup\varnothing=\varnothing\cup A}{%
      \FormulaRefAuto{A\cup B=B\cup A}}
    \proofstepwidestar[]{A=\varnothing\cup A}{\FormulaRefAuto{A=\varnothing\cup A}}
    \proofstepwidestar[]{\varnothing\cup A=A}{\FormulaRefAuto{a=b\vdash b=a}{3}}
    \proofstepwidestar[]{A\cup\varnothing=A}{\rIE{2,4}}
    \proofstepwidestar[]{A=A\cup\varnothing}{\FormulaRefAuto{a=b\vdash b=a}{5}}
    \proofstepwidestar[1]{\Finite{A\cup\varnothing}}{\rIE{6,1}}
  \closeproofpartwideindR

  \proofpartwideindR[Induktionsschritt]{%
    \begin{aligned}[t]
      &\Finite B\dsep \Finite{A\cup B}\dsep b\notin B\\
      &\vdash \Finite{A\cup(B\cup\{b\})}
    \end{aligned}%
  }{%
    \Finite B\dsep \Finite{A\cup B}\dsep b\notin B
    \vdash \Finite{A\cup(B\cup\{b\})}%
  }[FiniteUnionStep]
    \proofstepwidestar[1]{\Finite B}{\rA}
    \proofstepwidestar[2]{\Finite{A\cup B}}{\rA}
    \proofstepwidestar[3]{b\notin B}{\rA}
    \proofstepwidestar[2]{\Finite{(A\cup B)\cup\{b\}}}{%
      \FormulaRefAuto{FiniteInsert}[thm]{2}}
    \proofstepwidestar[]{A\cup(B\cup\{b\})=(A\cup B)\cup\{b\}}{%
      \FormulaRefAuto{A\cup(B\cup C)=(A\cup B)\cup C}}
    \proofstepwidestar[]{(A\cup B)\cup\{b\}=A\cup(B\cup\{b\})}{%
      \FormulaRefAuto{a=b\vdash b=a}{5}}
    \proofstepwidestar[2]{\Finite{A\cup(B\cup\{b\})}}{\rIE{6,4}}
  \closeproofpartwideindR

  \proofpartwide{Induktionsschluss}
    \proofstepwidestar[1]{\Finite A}{\rA}
    \proofstepwidestar[2]{\Finite B}{\rA}
    \proofstepwidestar[1,2]{\Finite{A\cup B}}{%
      \FormulaRefAuto{FiniteInductionRule}[ax]{IA,IS,2}}
  \closeproofpartwide
\end{tabproofsplitwide}

\section{Beispiele endlicher Mengen}

\subsection{Adjunktion eines neuen Elements}

Die Erhaltung der Endlichkeit unter Adjunktion ist hier die unmittelbare
Theoremform von \FormulaRefAuto{FiniteAdjunctionAxiom}[ax]. Die nachfolgende
ausführliche Konstruktion der erweiterten Bijektion bleibt als technische
Peano-Variante deaktiviert.

\iffalse

\FormulaThmDelta[Typisierung der vorangestellten Inklusionsabbildung]{%
  F\colon \NatSeg{n}\bij A \vdash \iota_{A,A\cup\{a\}}\circ F\colon \NatSeg{n}\to A\cup\{a\}
}{%
  \DeltaRow{Mengen}{A\dsep a\dsep n}
  \DeltaRow{Funktionen}{F}
}
\begin{tabproof}
  \proofstep{1}{F\colon \NatSeg{n}\bij A}{\rA}
  \proofstep{1}{F\colon \NatSeg{n}\to A}{\FormulaRefAuto{Bijektive Funktion}{1}}
  \proofstep{}{A\subseteq A\cup\{a\}}{\FormulaRefAuto{A\subseteq A\cup B}}
  \proofstep{1}{\iota_{A,A\cup\{a\}}\circ F\colon \NatSeg{n}\to A\cup\{a\}}{%
    \FormulaRefAuto{B\subseteq C \dsep F\colon A\to B \vdash \iota_{B,C}\circ F\colon A\to C}{3,2}}
\end{tabproof}

\FormulaThmDelta[Werte der vorangestellten Inklusionsabbildung]{%
  F\colon \NatSeg{n}\bij A \dsep u\in \NatSeg{n} \vdash
  \iota_{A,A\cup\{a\}}\circ F(u)=F(u)
}{%
  \DeltaRow{Mengen}{A\dsep a\dsep n\dsep u}
  \DeltaRow{Funktionen}{F}
}
\begin{tabproof}
  \proofstep{1}{F\colon \NatSeg{n}\bij A}{\rA}
  \proofstep{2}{u\in \NatSeg{n}}{\rA}
  \proofstep{1}{F\colon \NatSeg{n}\to A}{\FormulaRefAuto{Bijektive Funktion}{1}}
  \proofstep{2}{\iota_{A,A\cup\{a\}}\circ F(u)=\iota_{A,A\cup\{a\}}(F(u))}{%
    \FormulaRefAuto{x\in A \vdash (F\circ G)(x)=F(G(x))}{2}}
  \proofstep{1,2}{F(u)\in A}{\FormulaRefAuto{F\colon A\to B\dsep x\in A\vdash F(x)\in B}{3,2}}
  \proofstep{}{A\subseteq A\cup\{a\}}{\FormulaRefAuto{A\subseteq A\cup B}}
  \proofstep{5}{\iota_{A,A\cup\{a\}}(F(u))=F(u)}{%
    \FormulaRefAuto{A\subseteq B \dsep x\in A \vdash \iota_{A,B}(x)=x}{6,5}}
  \proofstep{1,2}{\iota_{A,A\cup\{a\}}\circ F(u)=F(u)}{\rIE{4,7}}
\end{tabproof}

\FormulaThmDelta[Werte der vorangestellten Inklusionsabbildung auf kleineren Punkten]{%
  n\in\mathbb{N}\dsep u\in\mathbb{N}\dsep F\colon \NatSeg{n}\bij A\dsep u<n \vdash
  \iota_{A,A\cup\{a\}}\circ F(u)=F(u)
}{%
  \DeltaRow{Mengen}{A\dsep a\dsep n\dsep u}
  \DeltaRow{Funktionen}{F}
  \DeltaRow{Relationen}{<_{\mathbb{N}}}
}
\begin{tabproof}
  \proofstep{1}{n\in\mathbb{N}}{\rA}
  \proofstep{2}{u\in\mathbb{N}}{\rA}
  \proofstep{3}{F\colon \NatSeg{n}\bij A}{\rA}
  \proofstep{4}{u<n}{\rA}
  \proofstep{1,2,4}{u\in \NatSeg{n}}{%
    \FormulaRefAuto{PeanoNatSegStrictMembership}{2,1,4}}
  \proofstep{1,2,3,4}{\iota_{A,A\cup\{a\}}\circ F(u)=F(u)}{%
    \FormulaRefAuto{F\colon \NatSeg{n}\bij A \dsep u\in \NatSeg{n} \vdash \iota_{A,A\cup\{a\}}\circ F(u)=F(u)}{3,5}}
\end{tabproof}

\FormulaThmDelta[Werte der adjungierten Abbildung auf alten Punkten]{%
  n\in\mathbb{N}\dsep F\colon \NatSeg{n}\bij A\dsep u\in \NatSeg{n} \vdash
  \ExtMap{\iota_{A,A\cup\{a\}}\circ F}{n}{a}(u)=F(u)
}{%
  \DeltaRow{Mengen}{A\dsep a\dsep n\dsep u}
  \DeltaRow{Funktionen}{F}
}
\begin{tabproof}
  \proofstep{1}{n\in\mathbb{N}}{\rA}
  \proofstep{2}{F\colon \NatSeg{n}\bij A}{\rA}
  \proofstep{3}{u\in \NatSeg{n}}{\rA}
  \proofstep{2}{\iota_{A,A\cup\{a\}}\circ F\colon \NatSeg{n}\to A\cup\{a\}}{%
    \FormulaRefAuto{F\colon \NatSeg{n}\bij A \vdash \iota_{A,A\cup\{a\}}\circ F\colon \NatSeg{n}\to A\cup\{a\}}{2}}
  \proofstep{}{a\in A\cup\{a\}}{\FormulaRefAuto{a\in A\cup\{a\}}}
  \proofstep{1}{n\notin \NatSeg{n}}{\FormulaRefAuto{n\in\mathbb{N}\vdash n\notin \NatSeg{n}}{1}}
  \proofstep{1,2,3}{\ExtMap{\iota_{A,A\cup\{a\}}\circ F}{n}{a}(u)=\iota_{A,A\cup\{a\}}\circ F(u)}{%
    \FormulaRefAuto{a\notin A \dsep f\colon A\to M \dsep b\in M \vdash \forall x\in A\,\ExtMap{f}{a}{b}(x)=f(x)}{6,4,5}}
  \proofstep{2,3}{\iota_{A,A\cup\{a\}}\circ F(u)=F(u)}{%
    \FormulaRefAuto{F\colon \NatSeg{n}\bij A \dsep u\in \NatSeg{n} \vdash \iota_{A,A\cup\{a\}}\circ F(u)=F(u)}{2,3}}
  \proofstep{1,2,3}{\ExtMap{\iota_{A,A\cup\{a\}}\circ F}{n}{a}(u)=F(u)}{\rIE{7,8}}
\end{tabproof}

\FormulaThmDelta[Werte der adjungierten Abbildung auf kleineren Punkten]{%
  n\in\mathbb{N}\dsep u\in\mathbb{N}\dsep F\colon \NatSeg{n}\bij A\dsep u < n \vdash
  \ExtMap{\iota_{A,A\cup\{a\}}\circ F}{n}{a}(u)=F(u)
}{%
  \DeltaRow{Mengen}{A\dsep a\dsep n\dsep u}
  \DeltaRow{Funktionen}{F}
  \DeltaRow{Relationen}{<_{\mathbb{N}}}
}
\begin{tabproof}
  \proofstep{1}{n\in\mathbb{N}}{\rA}
  \proofstep{2}{u\in\mathbb{N}}{\rA}
  \proofstep{3}{F\colon \NatSeg{n}\bij A}{\rA}
  \proofstep{4}{u < n}{\rA}
  \proofstep{1,2,4}{u\in \NatSeg{n}}{%
    \FormulaRefAuto{PeanoNatSegStrictMembership}{2,1,4}}
  \proofstep{1,2,3,4}{\ExtMap{\iota_{A,A\cup\{a\}}\circ F}{n}{a}(u)=F(u)}{%
    \FormulaRefAuto{n\in\mathbb{N}\dsep F\colon \NatSeg{n}\bij A\dsep u\in \NatSeg{n} \vdash \ExtMap{\iota_{A,A\cup\{a\}}\circ F}{n}{a}(u)=F(u)}{1,3,5}}
\end{tabproof}

\FormulaThmDelta[Wert der adjungierten Abbildung am neuen Punkt]{%
  n\in\mathbb{N}\dsep F\colon \NatSeg{n}\bij A \vdash
  \ExtMap{\iota_{A,A\cup\{a\}}\circ F}{n}{a}(n)=a
}{%
  \DeltaRow{Mengen}{A\dsep a\dsep n}
  \DeltaRow{Funktionen}{F}
}
\begin{tabproof}
  \proofstep{1}{n\in\mathbb{N}}{\rA}
  \proofstep{2}{F\colon \NatSeg{n}\bij A}{\rA}
  \proofstep{2}{\iota_{A,A\cup\{a\}}\circ F\colon \NatSeg{n}\to A\cup\{a\}}{%
    \FormulaRefAuto{F\colon \NatSeg{n}\bij A \vdash \iota_{A,A\cup\{a\}}\circ F\colon \NatSeg{n}\to A\cup\{a\}}{2}}
  \proofstep{}{a\in A\cup\{a\}}{\FormulaRefAuto{a\in A\cup\{a\}}}
  \proofstep{1}{n\notin \NatSeg{n}}{\FormulaRefAuto{n\in\mathbb{N}\vdash n\notin \NatSeg{n}}{1}}
  \proofstep{1,2}{\ExtMap{\iota_{A,A\cup\{a\}}\circ F}{n}{a}(n)=a}{%
    \FormulaRefAuto{a\notin A \dsep f\colon A\to M \dsep b\in M \vdash \ExtMap{f}{a}{b}(a)=b}{5,3,4}}
\end{tabproof}

\FormulaThmDelta[Typisierung der adjungierten Abbildung]{%
  n\in\mathbb{N}\dsep F\colon \NatSeg{n}\bij A \vdash
  \ExtMap{\iota_{A,A\cup\{a\}}\circ F}{n}{a}\colon (\NatSeg{n}\cup\{n\})\to A\cup\{a\}
}{%
  \DeltaRow{Mengen}{A\dsep a\dsep n}
  \DeltaRow{Funktionen}{F}
}
\begin{tabproof}
  \proofstep{1}{n\in\mathbb{N}}{\rA}
  \proofstep{2}{F\colon \NatSeg{n}\bij A}{\rA}
  \proofstep{2}{\iota_{A,A\cup\{a\}}\circ F\colon \NatSeg{n}\to A\cup\{a\}}{%
    \FormulaRefAuto{F\colon \NatSeg{n}\bij A \vdash \iota_{A,A\cup\{a\}}\circ F\colon \NatSeg{n}\to A\cup\{a\}}{2}}
  \proofstep{}{a\in A\cup\{a\}}{\FormulaRefAuto{a\in A\cup\{a\}}}
  \proofstep{1}{n\notin \NatSeg{n}}{\FormulaRefAuto{n\in\mathbb{N}\vdash n\notin \NatSeg{n}}{1}}
  \proofstep{1,2}{\ExtMap{\iota_{A,A\cup\{a\}}\circ F}{n}{a}\colon (\NatSeg{n}\cup\{n\})\to A\cup\{a\}}{%
    \FormulaRefAuto{a\notin A \dsep f\colon A\to M \dsep b\in M \vdash \ExtMap{f}{a}{b}\colon (A\cup\{a\})\to M}{5,3,4}}
\end{tabproof}

\FormulaThmDeltaR[Fall $x,y\in \NatSeg{n}$ der adjungierten Abbildung]{%
  \begin{aligned}[t]
    &n\in\mathbb{N}\dsep F\colon \NatSeg{n}\bij A\dsep x\in \NatSeg{n}\dsep y\in \NatSeg{n}\dsep{}\\
    &\ExtMap{\iota_{A,A\cup\{a\}}\circ F}{n}{a}(x)= {}\\
    &\qquad \ExtMap{\iota_{A,A\cup\{a\}}\circ F}{n}{a}(y)\vdash x=y
  \end{aligned}
}{%
  n\in\mathbb{N}\dsep F\colon \NatSeg{n}\bij A\dsep x\in \NatSeg{n}\dsep y\in \NatSeg{n}\dsep
  \ExtMap{\iota_{A,A\cup\{a\}}\circ F}{n}{a}(x)=
  \ExtMap{\iota_{A,A\cup\{a\}}\circ F}{n}{a}(y)\vdash x=y
}{%
  \DeltaRow{Mengen}{A\dsep a\dsep n\dsep x\dsep y}
  \DeltaRow{Funktionen}{F}
}
\begin{tabproof}
  \proofstep{1}{n\in\mathbb{N}}{\rA}
  \proofstep{2}{F\colon \NatSeg{n}\bij A}{\rA}
  \proofstep{3}{x\in \NatSeg{n}}{\rA}
  \proofstep{4}{y\in \NatSeg{n}}{\rA}
  \proofstep{5}{%
    \begin{aligned}[t]
      &\ExtMap{\iota_{A,A\cup\{a\}}\circ F}{n}{a}(x)= {}\\
      &\qquad \ExtMap{\iota_{A,A\cup\{a\}}\circ F}{n}{a}(y)
    \end{aligned}
  }{\rA}

  \proofstep{2}{F\colon \NatSeg{n}\inj A}{\FormulaRefAuto{F\colon A\bij B \vdash F\colon A\inj B}{2}}
  \proofstep{1,2,3}{\ExtMap{\iota_{A,A\cup\{a\}}\circ F}{n}{a}(x)=F(x)}{%
    \FormulaRefAuto{n\in\mathbb{N}\dsep F\colon \NatSeg{n}\bij A\dsep u\in \NatSeg{n} \vdash \ExtMap{\iota_{A,A\cup\{a\}}\circ F}{n}{a}(u)=F(u)}{1,2,3}}
  \proofstep{1,2,4}{\ExtMap{\iota_{A,A\cup\{a\}}\circ F}{n}{a}(y)=F(y)}{%
    \FormulaRefAuto{n\in\mathbb{N}\dsep F\colon \NatSeg{n}\bij A\dsep u\in \NatSeg{n} \vdash \ExtMap{\iota_{A,A\cup\{a\}}\circ F}{n}{a}(u)=F(u)}{1,2,4}}
  \proofstep{1,2,3,4}{F(x)=F(y)}{\rIE{7,\rIE{5,8}}}
  \proofstep{1,2,3,4}{x=y}{%
    \FormulaRefAuto{Injektivität}{6,3,4,9}}
\end{tabproof}

\FormulaThmDeltaR[Fall $x<n,\; y<n$ der adjungierten Abbildung]{%
  \begin{aligned}[t]
    &n\in\mathbb{N}\dsep x\in\mathbb{N}\dsep y\in\mathbb{N}\dsep{}\\
    &F\colon \NatSeg{n}\bij A\dsep x<n\dsep y<n\dsep{}\\
    &\ExtMap{\iota_{A,A\cup\{a\}}\circ F}{n}{a}(x)= {}\\
    &\qquad \ExtMap{\iota_{A,A\cup\{a\}}\circ F}{n}{a}(y)\vdash x=y
  \end{aligned}
}{%
  n\in\mathbb{N}\dsep x\in\mathbb{N}\dsep y\in\mathbb{N}\dsep
  F\colon \NatSeg{n}\bij A\dsep x<n\dsep y<n\dsep
  \ExtMap{\iota_{A,A\cup\{a\}}\circ F}{n}{a}(x)=
  \ExtMap{\iota_{A,A\cup\{a\}}\circ F}{n}{a}(y)\vdash x=y
}{%
  \DeltaRow{Mengen}{A\dsep a\dsep n\dsep x\dsep y}
  \DeltaRow{Funktionen}{F}
  \DeltaRow{Relationen}{<_{\mathbb{N}}}
}
\begin{tabproof}
  \proofstep{1}{n\in\mathbb{N}}{\rA}
  \proofstep{2}{x\in\mathbb{N}}{\rA}
  \proofstep{3}{y\in\mathbb{N}}{\rA}
  \proofstep{4}{F\colon \NatSeg{n}\bij A}{\rA}
  \proofstep{5}{x<n}{\rA}
  \proofstep{6}{y<n}{\rA}
  \proofstep{7}{%
    \begin{aligned}[t]
      &\ExtMap{\iota_{A,A\cup\{a\}}\circ F}{n}{a}(x)= {}\\
      &\qquad \ExtMap{\iota_{A,A\cup\{a\}}\circ F}{n}{a}(y)
    \end{aligned}
  }{\rA}

  \proofstep{1,2,5}{x\in \NatSeg{n}}{%
    \FormulaRefAuto{PeanoNatSegStrictMembership}{2,1,5}}
  \proofstep{1,3,6}{y\in \NatSeg{n}}{%
    \FormulaRefAuto{PeanoNatSegStrictMembership}{3,1,6}}
  \proofstep{1,2,3,4,5,6,7}{x=y}{%
    \FormulaRefAuto{%
      n\in\mathbb{N}\dsep F\colon \NatSeg{n}\bij A\dsep x\in \NatSeg{n}\dsep y\in \NatSeg{n}\dsep
      \ExtMap{\iota_{A,A\cup\{a\}}\circ F}{n}{a}(x)=
      \ExtMap{\iota_{A,A\cup\{a\}}\circ F}{n}{a}(y)\vdash x=y}{1,4,8,9,7}}
\end{tabproof}

\FormulaThmDeltaR[Fall $x\in \NatSeg{n},\; y=n$ der adjungierten Abbildung]{%
  \begin{aligned}[t]
    &a\notin A\dsep n\in\mathbb{N}\dsep F\colon \NatSeg{n}\bij A\dsep{}\\
    &x\in \NatSeg{n}\dsep y=n\dsep{}\\
    &\ExtMap{\iota_{A,A\cup\{a\}}\circ F}{n}{a}(x)= {}\\
    &\qquad \ExtMap{\iota_{A,A\cup\{a\}}\circ F}{n}{a}(y)\vdash x=y
  \end{aligned}
}{%
  a\notin A\dsep n\in\mathbb{N}\dsep F\colon \NatSeg{n}\bij A\dsep x\in \NatSeg{n}\dsep y=n\dsep
  \ExtMap{\iota_{A,A\cup\{a\}}\circ F}{n}{a}(x)=
  \ExtMap{\iota_{A,A\cup\{a\}}\circ F}{n}{a}(y)\vdash x=y
}{%
  \DeltaRow{Mengen}{A\dsep a\dsep n\dsep x\dsep y}
  \DeltaRow{Funktionen}{F}
}
\begin{tabproof}
  \proofstep{1}{a\notin A}{\rA}
  \proofstep{2}{n\in\mathbb{N}}{\rA}
  \proofstep{3}{F\colon \NatSeg{n}\bij A}{\rA}
  \proofstep{4}{x\in \NatSeg{n}}{\rA}
  \proofstep{5}{y=n}{\rA}
  \proofstep{6}{%
    \begin{aligned}[t]
      &\ExtMap{\iota_{A,A\cup\{a\}}\circ F}{n}{a}(x)= {}\\
      &\qquad \ExtMap{\iota_{A,A\cup\{a\}}\circ F}{n}{a}(y)
    \end{aligned}
  }{\rA}

  \proofstep{3}{F\colon \NatSeg{n}\to A}{\FormulaRefAuto{Bijektive Funktion}{3}}
  \proofstep{2,3,4}{\ExtMap{\iota_{A,A\cup\{a\}}\circ F}{n}{a}(x)=F(x)}{%
    \FormulaRefAuto{n\in\mathbb{N}\dsep F\colon \NatSeg{n}\bij A\dsep u\in \NatSeg{n} \vdash \ExtMap{\iota_{A,A\cup\{a\}}\circ F}{n}{a}(u)=F(u)}{2,3,4}}
  \proofstep{2,3}{\ExtMap{\iota_{A,A\cup\{a\}}\circ F}{n}{a}(n)=a}{%
    \FormulaRefAuto{n\in\mathbb{N}\dsep F\colon \NatSeg{n}\bij A \vdash \ExtMap{\iota_{A,A\cup\{a\}}\circ F}{n}{a}(n)=a}{2,3}}
  \proofstep{1,2,3,4,5}{F(x)=a}{\rIE{8,\rIE{6,\rIE{5,9}}}}
  \proofstep{3,4}{F(x)\in A}{\FormulaRefAuto{F\colon A\to B\dsep x\in A\vdash F(x)\in B}{7,4}}
  \proofstep{1,2,3,4,5}{a\in A}{\rIE{10,11}}
  \proofstep{1,2,3,4,5}{\bot}{\rBI{12,1}}
  \proofstep{1,2,3,4,5}{x=y}{\rCE{13}}
\end{tabproof}

\FormulaThmDeltaR[Fall $x<n,\; y=n$ der adjungierten Abbildung]{%
  \begin{aligned}[t]
    &a\notin A\dsep n\in\mathbb{N}\dsep x\in\mathbb{N}\dsep{}\\
    &F\colon \NatSeg{n}\bij A\dsep x<n\dsep y=n\dsep{}\\
    &\ExtMap{\iota_{A,A\cup\{a\}}\circ F}{n}{a}(x)= {}\\
    &\qquad \ExtMap{\iota_{A,A\cup\{a\}}\circ F}{n}{a}(y)\vdash x=y
  \end{aligned}
}{%
  a\notin A\dsep n\in\mathbb{N}\dsep x\in\mathbb{N}\dsep
  F\colon \NatSeg{n}\bij A\dsep x<n\dsep y=n\dsep
  \ExtMap{\iota_{A,A\cup\{a\}}\circ F}{n}{a}(x)=
  \ExtMap{\iota_{A,A\cup\{a\}}\circ F}{n}{a}(y)\vdash x=y
}{%
  \DeltaRow{Mengen}{A\dsep a\dsep n\dsep x\dsep y}
  \DeltaRow{Funktionen}{F}
  \DeltaRow{Relationen}{<_{\mathbb{N}}}
}
\begin{tabproof}
  \proofstep{1}{a\notin A}{\rA}
  \proofstep{2}{n\in\mathbb{N}}{\rA}
  \proofstep{3}{x\in\mathbb{N}}{\rA}
  \proofstep{4}{F\colon \NatSeg{n}\bij A}{\rA}
  \proofstep{5}{x<n}{\rA}
  \proofstep{6}{y=n}{\rA}
  \proofstep{7}{%
    \begin{aligned}[t]
      &\ExtMap{\iota_{A,A\cup\{a\}}\circ F}{n}{a}(x)= {}\\
      &\qquad \ExtMap{\iota_{A,A\cup\{a\}}\circ F}{n}{a}(y)
    \end{aligned}
  }{\rA}

  \proofstep{2,3,5}{x\in \NatSeg{n}}{%
    \FormulaRefAuto{PeanoNatSegStrictMembership}{3,2,5}}
  \proofstep{1,2,3,4,5,6,7}{x=y}{%
    \FormulaRefAuto{%
      a\notin A\dsep n\in\mathbb{N}\dsep F\colon \NatSeg{n}\bij A\dsep x\in \NatSeg{n}\dsep y=n\dsep
      \ExtMap{\iota_{A,A\cup\{a\}}\circ F}{n}{a}(x)=
      \ExtMap{\iota_{A,A\cup\{a\}}\circ F}{n}{a}(y)\vdash x=y}{1,2,4,8,6,7}}
\end{tabproof}

\FormulaThmDeltaR[Fall $x=n,\; y=n$ der adjungierten Abbildung]{%
  \begin{aligned}[t]
    &n\in\mathbb{N}\dsep x=n\dsep y=n\dsep{}\\
    &\ExtMap{\iota_{A,A\cup\{a\}}\circ F}{n}{a}(x)= {}\\
    &\qquad \ExtMap{\iota_{A,A\cup\{a\}}\circ F}{n}{a}(y)\vdash x=y
  \end{aligned}
}{%
  n\in\mathbb{N}\dsep x=n\dsep y=n\dsep
  \ExtMap{\iota_{A,A\cup\{a\}}\circ F}{n}{a}(x)=
  \ExtMap{\iota_{A,A\cup\{a\}}\circ F}{n}{a}(y)\vdash x=y
}{%
  \DeltaRow{Mengen}{A\dsep a\dsep n\dsep x\dsep y}
  \DeltaRow{Funktionen}{F}
}
\begin{tabproof}
  \proofstep{1}{n\in\mathbb{N}}{\rA}
  \proofstep{2}{x=n}{\rA}
  \proofstep{3}{y=n}{\rA}
  \proofstep{4}{%
    \begin{aligned}[t]
      &\ExtMap{\iota_{A,A\cup\{a\}}\circ F}{n}{a}(x)= {}\\
      &\qquad \ExtMap{\iota_{A,A\cup\{a\}}\circ F}{n}{a}(y)
    \end{aligned}
  }{\rA}
  \proofstep{3}{n=y}{\FormulaRefAuto{a=b\vdash b=a}{3}}
  \proofstep{2,3}{x=y}{\rIE{2,5}}
\end{tabproof}

\FormulaThmDelta[Injektivität der adjungierten Abbildung]{%
  \begin{aligned}[t]
    &a\notin A \dsep n\in\mathbb{N}\dsep F\colon \NatSeg{n}\bij A\dsep{}\\
    &x\in \NatSeg{n}\cup\{n\}\dsep y\in \NatSeg{n}\cup\{n\}\dsep{}\\
    &\ExtMap{\iota_{A,A\cup\{a\}}\circ F}{n}{a}(x)= {}\\
    &\qquad \ExtMap{\iota_{A,A\cup\{a\}}\circ F}{n}{a}(y)\vdash x=y
  \end{aligned}
}{%
  \DeltaRow{Mengen}{A\dsep a\dsep n\dsep x\dsep y}
  \DeltaRow{Funktionen}{F}
}
\begin{tabproofsplitwide}
  \proofpartwideindR[Fall $x\in \NatSeg{n}$]{%
    \begin{aligned}[t]
      &a\notin A\dsep n\in\mathbb{N}\dsep F\colon \NatSeg{n}\bij A\dsep{}\\
      &x\in \NatSeg{n}\cup\{n\}\dsep y\in \NatSeg{n}\cup\{n\}\dsep x\in \NatSeg{n}\dsep{}\\
      &\ExtMap{\iota_{A,A\cup\{a\}}\circ F}{n}{a}(x)= {}\\
      &\qquad \ExtMap{\iota_{A,A\cup\{a\}}\circ F}{n}{a}(y)\vdash x=y
    \end{aligned}
  }{%
    a\notin A\dsep n\in\mathbb{N}\dsep F\colon \NatSeg{n}\bij A\dsep
    x\in \NatSeg{n}\cup\{n\}\dsep y\in \NatSeg{n}\cup\{n\}\dsep x\in \NatSeg{n}\dsep
    \ExtMap{\iota_{A,A\cup\{a\}}\circ F}{n}{a}(x)=
    \ExtMap{\iota_{A,A\cup\{a\}}\circ F}{n}{a}(y)\vdash x=y
  }
    \proofstepwidestar[1]{a\notin A}{\rA}
    \proofstepwidestar[2]{n\in\mathbb{N}}{\rA}
    \proofstepwidestar[3]{F\colon \NatSeg{n}\bij A}{\rA}
    \proofstepwidestar[4]{x\in \NatSeg{n}\cup\{n\}}{\rA}
    \proofstepwidestar[5]{y\in \NatSeg{n}\cup\{n\}}{\rA}
    \proofstepwidestar[6]{x\in \NatSeg{n}}{\rA}
    \proofstepwidestar[7]{%
      \begin{aligned}[t]
        &\ExtMap{\iota_{A,A\cup\{a\}}\circ F}{n}{a}(x)= {}\\
        &\qquad \ExtMap{\iota_{A,A\cup\{a\}}\circ F}{n}{a}(y)
      \end{aligned}
    }{\rA}

    \proofstepwidestar[5]{y\in \NatSeg{n}\lor y=n}{%
      \FormulaRefAuto{x\in A\cup\{a\} \vdash x\in A\lor x=a}{5}}

    \proofstepwidestar[9]{y\in \NatSeg{n}}{\rA}
    \proofstepwidestar[2,3,6,9]{x=y}{%
      \FormulaRefAuto{n\in\mathbb{N}\dsep F\colon \NatSeg{n}\bij A\dsep x\in \NatSeg{n}\dsep y\in \NatSeg{n}\dsep \ExtMap{\iota_{A,A\cup\{a\}}\circ F}{n}{a}(x)=\ExtMap{\iota_{A,A\cup\{a\}}\circ F}{n}{a}(y)\vdash x=y}{2,3,6,9,7}}

    \proofstepwidestar[11]{y=n}{\rA}
    \proofstepwidestar[1,2,3,6,11]{x=y}{%
      \FormulaRefAuto{a\notin A\dsep n\in\mathbb{N}\dsep F\colon \NatSeg{n}\bij A\dsep x\in \NatSeg{n}\dsep y=n\dsep \ExtMap{\iota_{A,A\cup\{a\}}\circ F}{n}{a}(x)=\ExtMap{\iota_{A,A\cup\{a\}}\circ F}{n}{a}(y)\vdash x=y}{1,2,3,6,11,7}}

    \proofstepwidestar[1,2,3,5,6]{x=y}{\rOE{8,9,10,11,12}}
  \closeproofpartwideindR

  \proofpartwideindR[Fall $x=n$]{%
    \begin{aligned}[t]
      &a\notin A\dsep n\in\mathbb{N}\dsep F\colon \NatSeg{n}\bij A\dsep{}\\
      &x\in \NatSeg{n}\cup\{n\}\dsep y\in \NatSeg{n}\cup\{n\}\dsep x=n\dsep{}\\
      &\ExtMap{\iota_{A,A\cup\{a\}}\circ F}{n}{a}(x)= {}\\
      &\qquad \ExtMap{\iota_{A,A\cup\{a\}}\circ F}{n}{a}(y)\vdash x=y
    \end{aligned}
  }{%
    a\notin A\dsep n\in\mathbb{N}\dsep F\colon \NatSeg{n}\bij A\dsep
    x\in \NatSeg{n}\cup\{n\}\dsep y\in \NatSeg{n}\cup\{n\}\dsep x=n\dsep
    \ExtMap{\iota_{A,A\cup\{a\}}\circ F}{n}{a}(x)=
    \ExtMap{\iota_{A,A\cup\{a\}}\circ F}{n}{a}(y)\vdash x=y
  }
    \proofstepwidestar[1]{a\notin A}{\rA}
    \proofstepwidestar[2]{n\in\mathbb{N}}{\rA}
    \proofstepwidestar[3]{F\colon \NatSeg{n}\bij A}{\rA}
    \proofstepwidestar[4]{x\in \NatSeg{n}\cup\{n\}}{\rA}
    \proofstepwidestar[5]{y\in \NatSeg{n}\cup\{n\}}{\rA}
    \proofstepwidestar[6]{x=n}{\rA}
    \proofstepwidestar[7]{%
      \begin{aligned}[t]
        &\ExtMap{\iota_{A,A\cup\{a\}}\circ F}{n}{a}(x)= {}\\
        &\qquad \ExtMap{\iota_{A,A\cup\{a\}}\circ F}{n}{a}(y)
      \end{aligned}
    }{\rA}

    \proofstepwidestar[5]{y\in \NatSeg{n}\lor y=n}{%
      \FormulaRefAuto{x\in A\cup\{a\} \vdash x\in A\lor x=a}{5}}

    \proofstepwidestar[9]{y\in \NatSeg{n}}{\rA}
    \proofstepwidestar[7]{%
      \begin{aligned}[t]
        &\ExtMap{\iota_{A,A\cup\{a\}}\circ F}{n}{a}(y)= {}\\
        &\qquad \ExtMap{\iota_{A,A\cup\{a\}}\circ F}{n}{a}(x)
      \end{aligned}
    }{\FormulaRefAuto{a=b\vdash b=a}{7}}
    \proofstepwidestar[1,2,3,6,9]{y=x}{%
      \FormulaRefAuto{a\notin A\dsep n\in\mathbb{N}\dsep F\colon \NatSeg{n}\bij A\dsep x\in \NatSeg{n}\dsep y=n\dsep \ExtMap{\iota_{A,A\cup\{a\}}\circ F}{n}{a}(x)=\ExtMap{\iota_{A,A\cup\{a\}}\circ F}{n}{a}(y)\vdash x=y}{1,2,3,9,6,10}}
    \proofstepwidestar[1,2,3,6,9]{x=y}{%
      \FormulaRefAuto{a=b\vdash b=a}{11}}

    \proofstepwidestar[13]{y=n}{\rA}
    \proofstepwidestar[2,6,13]{x=y}{%
      \FormulaRefAuto{n\in\mathbb{N}\dsep x=n\dsep y=n\dsep \ExtMap{\iota_{A,A\cup\{a\}}\circ F}{n}{a}(x)=\ExtMap{\iota_{A,A\cup\{a\}}\circ F}{n}{a}(y)\vdash x=y}{2,6,13,7}}

    \proofstepwidestar[1,2,3,5,6]{x=y}{\rOE{8,9,12,13,14}}
  \closeproofpartwideindR

  \proofpartwide{Schluss}
    \proofstepwidestar[1]{a\notin A}{\rA}
    \proofstepwidestar[2]{n\in\mathbb{N}}{\rA}
    \proofstepwidestar[3]{F\colon \NatSeg{n}\bij A}{\rA}
    \proofstepwidestar[4]{x\in \NatSeg{n}\cup\{n\}}{\rA}
    \proofstepwidestar[5]{y\in \NatSeg{n}\cup\{n\}}{\rA}
    \proofstepwidestar[6]{%
      \begin{aligned}[t]
        &\ExtMap{\iota_{A,A\cup\{a\}}\circ F}{n}{a}(x)= {}\\
        &\qquad \ExtMap{\iota_{A,A\cup\{a\}}\circ F}{n}{a}(y)
      \end{aligned}
    }{\rA}

    \proofstepwidestar[4]{x\in \NatSeg{n}\lor x=n}{%
      \FormulaRefAuto{x\in A\cup\{a\} \vdash x\in A\lor x=a}{4}}

    \proofstepwidestar[8]{x\in \NatSeg{n}}{\rA}
    \proofstepwidestar[1,2,3,4,5,8]{x=y}{%
      \FormulaRefAuto{a\notin A\dsep n\in\mathbb{N}\dsep F\colon \NatSeg{n}\bij A\dsep x\in \NatSeg{n}\cup\{n\}\dsep y\in \NatSeg{n}\cup\{n\}\dsep x\in \NatSeg{n}\dsep \ExtMap{\iota_{A,A\cup\{a\}}\circ F}{n}{a}(x)=\ExtMap{\iota_{A,A\cup\{a\}}\circ F}{n}{a}(y)\vdash x=y}{1,2,3,4,5,8,6}}

    \proofstepwidestar[10]{x=n}{\rA}
    \proofstepwidestar[1,2,3,4,5,10]{x=y}{%
      \FormulaRefAuto{a\notin A\dsep n\in\mathbb{N}\dsep F\colon \NatSeg{n}\bij A\dsep x\in \NatSeg{n}\cup\{n\}\dsep y\in \NatSeg{n}\cup\{n\}\dsep x=n\dsep \ExtMap{\iota_{A,A\cup\{a\}}\circ F}{n}{a}(x)=\ExtMap{\iota_{A,A\cup\{a\}}\circ F}{n}{a}(y)\vdash x=y}{1,2,3,4,5,10,6}}

    \proofstepwidestar[1,2,3,4,5]{x=y}{\rOE{7,8,9,10,11}}
  \closeproofpartwide
\end{tabproofsplitwide}

\FormulaThmDeltaR[Injektivität der adjungierten Abbildung in Ordnungssprache]{%
  \begin{aligned}[t]
    &a\notin A \dsep n\in\mathbb{N}\dsep x\in\mathbb{N}\dsep y\in\mathbb{N}\dsep{}\\
    &F\colon \NatSeg{n}\bij A\dsep x\leq n\dsep y\leq n\dsep{}\\
    &\ExtMap{\iota_{A,A\cup\{a\}}\circ F}{n}{a}(x)= {}\\
    &\qquad \ExtMap{\iota_{A,A\cup\{a\}}\circ F}{n}{a}(y)\vdash x=y
  \end{aligned}
}{%
  a\notin A \dsep n\in\mathbb{N}\dsep x\in\mathbb{N}\dsep y\in\mathbb{N}\dsep
  F\colon \NatSeg{n}\bij A\dsep x\leq n\dsep y\leq n\dsep
  \ExtMap{\iota_{A,A\cup\{a\}}\circ F}{n}{a}(x)=
  \ExtMap{\iota_{A,A\cup\{a\}}\circ F}{n}{a}(y)\vdash x=y
}{%
  \DeltaRow{Mengen}{A\dsep a\dsep n\dsep x\dsep y}
  \DeltaRow{Funktionen}{F}
  \DeltaRow{Relationen}{\leq_{\mathbb{N}}}
}
\begin{tabproof}
  \proofstep{1}{a\notin A}{\rA}
  \proofstep{2}{n\in\mathbb{N}}{\rA}
  \proofstep{3}{x\in\mathbb{N}}{\rA}
  \proofstep{4}{y\in\mathbb{N}}{\rA}
  \proofstep{5}{F\colon \NatSeg{n}\bij A}{\rA}
  \proofstep{6}{x\leq n}{\rA}
  \proofstep{7}{y\leq n}{\rA}
  \proofstep{8}{%
    \begin{aligned}[t]
      &\ExtMap{\iota_{A,A\cup\{a\}}\circ F}{n}{a}(x)= {}\\
      &\qquad \ExtMap{\iota_{A,A\cup\{a\}}\circ F}{n}{a}(y)
    \end{aligned}
  }{\rA}

  \proofstep{2,3,6}{x\in\NatSeg{\Succ(n)}}{%
    \FormulaRefAuto{PeanoNatSegLeqSuccMembership}{3,2,6}}
  \proofstep{2,4,7}{y\in\NatSeg{\Succ(n)}}{%
    \FormulaRefAuto{PeanoNatSegLeqSuccMembership}{4,2,7}}
  \proofstep{2}{\NatSeg{\Succ(n)}=\NatSeg{n}\cup\{n\}}{%
    \FormulaRefAuto{PeanoNatSegSucc}{2}}
  \proofstep{2,3,6}{x\in \NatSeg{n}\cup\{n\}}{%
    \FormulaRefAuto{A = B,\, x \in A \vdash x \in B}{11,9}}
  \proofstep{2,4,7}{y\in \NatSeg{n}\cup\{n\}}{%
    \FormulaRefAuto{A = B,\, x \in A \vdash x \in B}{11,10}}
  \proofstep{1,2,3,4,5,6,7,8}{x=y}{%
    \FormulaRefAuto{%
      \begin{aligned}[t]
        &a\notin A \dsep n\in\mathbb{N}\dsep F\colon \NatSeg{n}\bij A\dsep{}\\
        &x\in \NatSeg{n}\cup\{n\}\dsep y\in \NatSeg{n}\cup\{n\}\dsep{}\\
        &\ExtMap{\iota_{A,A\cup\{a\}}\circ F}{n}{a}(x)= {}\\
        &\qquad \ExtMap{\iota_{A,A\cup\{a\}}\circ F}{n}{a}(y)\vdash x=y
      \end{aligned}
    }{1,2,5,12,13,8}}
\end{tabproof}

\FormulaThmDelta[Fall $z\in A$ der Surjektivität der adjungierten Abbildung]{%
  n\in\mathbb{N}\dsep F\colon \NatSeg{n}\bij A\dsep z\in A \vdash
  \exists x\in \NatSeg{n}\cup\{n\}\,\ExtMap{\iota_{A,A\cup\{a\}}\circ F}{n}{a}(x)=z
}{%
  \DeltaRow{Mengen}{A\dsep a\dsep n\dsep z\dsep x}
  \DeltaRow{Funktionen}{F}
}
\begin{tabproof}
  \proofstep{1}{n\in\mathbb{N}}{\rA}
  \proofstep{2}{F\colon \NatSeg{n}\bij A}{\rA}
  \proofstep{3}{z\in A}{\rA}

  \proofstep{2}{F\colon \NatSeg{n}\sur A}{%
    \FormulaRefAuto{F\colon A\bij B \vdash F\colon A\sur B}{2}}
  \proofstep{2,3}{\exists x\in \NatSeg{n}\,F(x)=z}{%
    \FormulaRefAuto{Surjektivität}{4,3}}

  \proofstep{6}{x\in \NatSeg{n}\land F(x)=z}{\rA}
  \proofstep{6}{x\in \NatSeg{n}}{\rAEa{6}}
  \proofstep{6}{F(x)=z}{\rAEb{6}}
  \proofstep{6}{x\in \NatSeg{n}\cup\{n\}}{%
    \FormulaRefAuto{z\in A\vdash z\in A\cup B}{7}}
  \proofstep{1,2,6}{\ExtMap{\iota_{A,A\cup\{a\}}\circ F}{n}{a}(x)=F(x)}{%
    \FormulaRefAuto{n\in\mathbb{N}\dsep F\colon \NatSeg{n}\bij A\dsep u\in \NatSeg{n} \vdash \ExtMap{\iota_{A,A\cup\{a\}}\circ F}{n}{a}(u)=F(u)}{1,2,7}}
  \proofstep{1,2,6}{\ExtMap{\iota_{A,A\cup\{a\}}\circ F}{n}{a}(x)=z}{\rIE{10,8}}
  \proofstep{1,2,6}{\exists u\in \NatSeg{n}\cup\{n\}\,\ExtMap{\iota_{A,A\cup\{a\}}\circ F}{n}{a}(u)=z}{%
    \rEI{\rAI{9,11}}}

  \proofstep{1,2,3}{\exists u\in \NatSeg{n}\cup\{n\}\,\ExtMap{\iota_{A,A\cup\{a\}}\circ F}{n}{a}(u)=z}{%
    \rEE{5,6,12}}
\end{tabproof}

\FormulaThmDelta[Fall $z\in A$ der Surjektivität der adjungierten Abbildung in Ordnungssprache]{%
  n\in\mathbb{N}\dsep F\colon \NatSeg{n}\bij A\dsep z\in A \vdash
  \exists x\in\mathbb{N}\,\bigl(x<n\land
  \ExtMap{\iota_{A,A\cup\{a\}}\circ F}{n}{a}(x)=z\bigr)
}{%
  \DeltaRow{Mengen}{A\dsep a\dsep n\dsep z\dsep x}
  \DeltaRow{Funktionen}{F}
  \DeltaRow{Relationen}{<_{\mathbb{N}}}
}
\begin{tabproof}
  \proofstep{1}{n\in\mathbb{N}}{\rA}
  \proofstep{2}{F\colon \NatSeg{n}\bij A}{\rA}
  \proofstep{3}{z\in A}{\rA}
  \proofstep{2}{F\colon \NatSeg{n}\sur A}{%
    \FormulaRefAuto{F\colon A\bij B \vdash F\colon A\sur B}{2}}
  \proofstep{2,3}{\exists x\in \NatSeg{n}\,F(x)=z}{%
    \FormulaRefAuto{Surjektivität}{4,3}}

  \proofstep{6}{x\in \NatSeg{n}\land F(x)=z}{\rA}
  \proofstep{6}{x\in \NatSeg{n}}{\rAEa{6}}
  \proofstep{6}{F(x)=z}{\rAEb{6}}
  \proofstep{1}{\NatSeg{n}\subseteq\mathbb{N}}{%
    \FormulaRefAuto{PeanoNatSegSubsetN}{1}}
  \proofstep{1,6}{x\in\mathbb{N}}{%
    \FormulaRefAuto{A\subseteq B,\, x\in A\vdash x\in B}{9,7}}
  \proofstep{1,6}{x<n}{%
    \FormulaRefAuto{PeanoNatSegStrictMembership}{10,1,7}}
  \proofstep{1,2,6}{\ExtMap{\iota_{A,A\cup\{a\}}\circ F}{n}{a}(x)=F(x)}{%
    \FormulaRefAuto{n\in\mathbb{N}\dsep F\colon \NatSeg{n}\bij A\dsep u\in \NatSeg{n} \vdash \ExtMap{\iota_{A,A\cup\{a\}}\circ F}{n}{a}(u)=F(u)}{1,2,7}}
  \proofstep{1,2,6}{\ExtMap{\iota_{A,A\cup\{a\}}\circ F}{n}{a}(x)=z}{\rIE{12,8}}
  \proofstep{1,2,6}{x<n\land \ExtMap{\iota_{A,A\cup\{a\}}\circ F}{n}{a}(x)=z}{\rAI{11,13}}
  \proofstep{1,2,6}{\exists u\in\mathbb{N}\,\bigl(u<n\land \ExtMap{\iota_{A,A\cup\{a\}}\circ F}{n}{a}(u)=z\bigr)}{%
    \rEI{\rAI{10,14}}}
  \proofstep{1,2,3}{\exists u\in\mathbb{N}\,\bigl(u<n\land \ExtMap{\iota_{A,A\cup\{a\}}\circ F}{n}{a}(u)=z\bigr)}{%
    \rEE{5,6,15}}
\end{tabproof}

\FormulaThmDelta[Fall $z=a$ der Surjektivität der adjungierten Abbildung]{%
  n\in\mathbb{N}\dsep F\colon \NatSeg{n}\bij A\dsep z=a \vdash
  \exists x\in \NatSeg{n}\cup\{n\}\,\ExtMap{\iota_{A,A\cup\{a\}}\circ F}{n}{a}(x)=z
}{%
  \DeltaRow{Mengen}{A\dsep a\dsep n\dsep z}
  \DeltaRow{Funktionen}{F}
}
\begin{tabproof}
  \proofstep{1}{n\in\mathbb{N}}{\rA}
  \proofstep{2}{F\colon \NatSeg{n}\bij A}{\rA}
  \proofstep{3}{z=a}{\rA}

  \proofstep{}{n\in \NatSeg{n}\cup\{n\}}{\FormulaRefAuto{a\in A\cup\{a\}}}
  \proofstep{1,2}{\ExtMap{\iota_{A,A\cup\{a\}}\circ F}{n}{a}(n)=a}{%
    \FormulaRefAuto{n\in\mathbb{N}\dsep F\colon \NatSeg{n}\bij A \vdash \ExtMap{\iota_{A,A\cup\{a\}}\circ F}{n}{a}(n)=a}{1,2}}
  \proofstep{3}{a=z}{\FormulaRefAuto{a=b\vdash b=a}{3}}
  \proofstep{1,2,3}{\ExtMap{\iota_{A,A\cup\{a\}}\circ F}{n}{a}(n)=z}{\rIE{5,6}}
  \proofstep{1,2,3}{\exists u\in \NatSeg{n}\cup\{n\}\,\ExtMap{\iota_{A,A\cup\{a\}}\circ F}{n}{a}(u)=z}{\rEI{\rAI{4,7}}}
\end{tabproof}

\FormulaThmDelta[Fall $z=a$ der Surjektivität der adjungierten Abbildung in Ordnungssprache]{%
  n\in\mathbb{N}\dsep F\colon \NatSeg{n}\bij A\dsep z=a \vdash
  \exists x\in\mathbb{N}\,\bigl(x\leq n\land
  \ExtMap{\iota_{A,A\cup\{a\}}\circ F}{n}{a}(x)=z\bigr)
}{%
  \DeltaRow{Mengen}{A\dsep a\dsep n\dsep z}
  \DeltaRow{Funktionen}{F}
  \DeltaRow{Relationen}{\leq_{\mathbb{N}}}
}
\begin{tabproof}
  \proofstep{1}{n\in\mathbb{N}}{\rA}
  \proofstep{2}{F\colon \NatSeg{n}\bij A}{\rA}
  \proofstep{3}{z=a}{\rA}
  \proofstep{}{\NatSeg{n}\subseteq \NatSeg{n}}{\FormulaRefAuto{A\subseteq A}}
  \proofstep{1}{n\leq n}{%
    \FormulaRefAuto{m\in\mathbb{N}\dsep n\in\mathbb{N}\vdash m\leq n \leftrightarrow \NatSeg{m}\subseteq \NatSeg{n}}{1,1,4}}
  \proofstep{1,2}{\ExtMap{\iota_{A,A\cup\{a\}}\circ F}{n}{a}(n)=a}{%
    \FormulaRefAuto{n\in\mathbb{N}\dsep F\colon \NatSeg{n}\bij A \vdash \ExtMap{\iota_{A,A\cup\{a\}}\circ F}{n}{a}(n)=a}{1,2}}
  \proofstep{3}{a=z}{\FormulaRefAuto{a=b\vdash b=a}{3}}
  \proofstep{1,2,3}{\ExtMap{\iota_{A,A\cup\{a\}}\circ F}{n}{a}(n)=z}{\rIE{6,7}}
  \proofstep{1,2,3}{n\leq n\land \ExtMap{\iota_{A,A\cup\{a\}}\circ F}{n}{a}(n)=z}{\rAI{5,8}}
  \proofstep{1,2,3}{\exists u\in\mathbb{N}\,\bigl(u\leq n\land \ExtMap{\iota_{A,A\cup\{a\}}\circ F}{n}{a}(u)=z\bigr)}{%
    \rEI{\rAI{1,9}}}
\end{tabproof}

\FormulaThmDelta[Surjektivität der adjungierten Abbildung]{%
  \begin{aligned}[t]
    &a\notin A \dsep n\in\mathbb{N}\dsep F\colon \NatSeg{n}\bij A\dsep z\in A\cup\{a\}\\
    &\vdash \exists x\in \NatSeg{n}\cup\{n\}\,
      \ExtMap{\iota_{A,A\cup\{a\}}\circ F}{n}{a}(x)=z
  \end{aligned}
}{%
  \DeltaRow{Mengen}{A\dsep a\dsep n\dsep z\dsep x}
  \DeltaRow{Funktionen}{F}
}
\begin{tabproofsplitwide}
  \proofpartwideindR[Fall $z\in A$]{%
    \begin{aligned}[t]
      &a\notin A\dsep n\in\mathbb{N}\dsep F\colon \NatSeg{n}\bij A\dsep z\in A\cup\{a\}\dsep z\in A\\
      &\vdash \exists x\in \NatSeg{n}\cup\{n\}\,
      \ExtMap{\iota_{A,A\cup\{a\}}\circ F}{n}{a}(x)=z
    \end{aligned}
  }{%
    \begin{aligned}[t]
      &a\notin A\dsep n\in\mathbb{N}\dsep F\colon \NatSeg{n}\bij A\dsep z\in A\cup\{a\}\dsep z\in A\\
      &\vdash \exists x\in \NatSeg{n}\cup\{n\}\,
      \ExtMap{\iota_{A,A\cup\{a\}}\circ F}{n}{a}(x)=z
    \end{aligned}
  }
    \proofstepwidestar[1]{a\notin A}{\rA}
    \proofstepwidestar[2]{n\in\mathbb{N}}{\rA}
    \proofstepwidestar[3]{F\colon \NatSeg{n}\bij A}{\rA}
    \proofstepwidestar[4]{z\in A\cup\{a\}}{\rA}
    \proofstepwidestar[5]{z\in A}{\rA}

    \proofstepwidestar[2,3,5]{\exists x\in \NatSeg{n}\cup\{n\}\,\ExtMap{\iota_{A,A\cup\{a\}}\circ F}{n}{a}(x)=z}{%
      \FormulaRefAuto{n\in\mathbb{N}\dsep F\colon \NatSeg{n}\bij A\dsep z\in A \vdash \exists x\in \NatSeg{n}\cup\{n\}\,\ExtMap{\iota_{A,A\cup\{a\}}\circ F}{n}{a}(x)=z}{2,3,5}}
  \closeproofpartwideindR

  \proofpartwideindR[Fall $z=a$]{%
    \begin{aligned}[t]
      &a\notin A\dsep n\in\mathbb{N}\dsep F\colon \NatSeg{n}\bij A\dsep z\in A\cup\{a\}\dsep z=a\\
      &\vdash \exists x\in \NatSeg{n}\cup\{n\}\,
      \ExtMap{\iota_{A,A\cup\{a\}}\circ F}{n}{a}(x)=z
    \end{aligned}
  }{%
    \begin{aligned}[t]
      &a\notin A\dsep n\in\mathbb{N}\dsep F\colon \NatSeg{n}\bij A\dsep z\in A\cup\{a\}\dsep z=a\\
      &\vdash \exists x\in \NatSeg{n}\cup\{n\}\,
      \ExtMap{\iota_{A,A\cup\{a\}}\circ F}{n}{a}(x)=z
    \end{aligned}
  }
    \proofstepwidestar[1]{a\notin A}{\rA}
    \proofstepwidestar[2]{n\in\mathbb{N}}{\rA}
    \proofstepwidestar[3]{F\colon \NatSeg{n}\bij A}{\rA}
    \proofstepwidestar[4]{z\in A\cup\{a\}}{\rA}
    \proofstepwidestar[5]{z=a}{\rA}

    \proofstepwidestar[2,3,5]{\exists x\in \NatSeg{n}\cup\{n\}\,\ExtMap{\iota_{A,A\cup\{a\}}\circ F}{n}{a}(x)=z}{%
      \FormulaRefAuto{n\in\mathbb{N}\dsep F\colon \NatSeg{n}\bij A\dsep z=a \vdash \exists x\in \NatSeg{n}\cup\{n\}\,\ExtMap{\iota_{A,A\cup\{a\}}\circ F}{n}{a}(x)=z}{2,3,5}}
  \closeproofpartwideindR

  \proofpartwide{Schluss}
    \proofstepwidestar[1]{a\notin A}{\rA}
    \proofstepwidestar[2]{n\in\mathbb{N}}{\rA}
    \proofstepwidestar[3]{F\colon \NatSeg{n}\bij A}{\rA}
    \proofstepwidestar[4]{z\in A\cup\{a\}}{\rA}

    \proofstepwidestar[4]{z\in A\lor z=a}{%
      \FormulaRefAuto{x\in A\cup\{a\} \vdash x\in A\lor x=a}{4}}

    \proofstepwidestar[6]{z\in A}{\rA}
    \proofstepwidestar[1,2,3,4,6]{\exists x\in \NatSeg{n}\cup\{n\}\,\ExtMap{\iota_{A,A\cup\{a\}}\circ F}{n}{a}(x)=z}{%
      \FormulaRefAuto{%
        \begin{aligned}[t]
          &a\notin A\dsep n\in\mathbb{N}\dsep F\colon \NatSeg{n}\bij A\dsep z\in A\cup\{a\}\dsep z\in A\\
          &\vdash \exists x\in \NatSeg{n}\cup\{n\}\,
          \ExtMap{\iota_{A,A\cup\{a\}}\circ F}{n}{a}(x)=z
        \end{aligned}
      }{1,2,3,4,6}}

    \proofstepwidestar[8]{z=a}{\rA}
    \proofstepwidestar[1,2,3,4,8]{\exists x\in \NatSeg{n}\cup\{n\}\,\ExtMap{\iota_{A,A\cup\{a\}}\circ F}{n}{a}(x)=z}{%
      \FormulaRefAuto{%
        \begin{aligned}[t]
          &a\notin A\dsep n\in\mathbb{N}\dsep F\colon \NatSeg{n}\bij A\dsep z\in A\cup\{a\}\dsep z=a\\
          &\vdash \exists x\in \NatSeg{n}\cup\{n\}\,
          \ExtMap{\iota_{A,A\cup\{a\}}\circ F}{n}{a}(x)=z
        \end{aligned}
      }{1,2,3,4,8}}

    \proofstepwidestar[1,2,3,4]{\exists x\in \NatSeg{n}\cup\{n\}\,\ExtMap{\iota_{A,A\cup\{a\}}\circ F}{n}{a}(x)=z}{%
      \rOE{5,6,7,8,9}}
  \closeproofpartwide
\end{tabproofsplitwide}

\FormulaThmDelta[Surjektivität der adjungierten Abbildung in Ordnungssprache]{%
  \begin{aligned}[t]
    &a\notin A \dsep n\in\mathbb{N}\dsep F\colon \NatSeg{n}\bij A\dsep z\in A\cup\{a\}\\
    &\vdash \exists x\in\mathbb{N}\,\bigl(x\leq n\land
      \ExtMap{\iota_{A,A\cup\{a\}}\circ F}{n}{a}(x)=z\bigr)
  \end{aligned}
}{%
  \DeltaRow{Mengen}{A\dsep a\dsep n\dsep z\dsep x}
  \DeltaRow{Funktionen}{F}
  \DeltaRow{Relationen}{<_{\mathbb{N}}\dsep \leq_{\mathbb{N}}}
}
\begin{tabproofsplitwide}
  \proofpartwideindR[Fall $z\in A$]{%
    \begin{aligned}[t]
      &a\notin A\dsep n\in\mathbb{N}\dsep F\colon \NatSeg{n}\bij A\dsep z\in A\cup\{a\}\dsep z\in A\\
      &\vdash \exists x\in\mathbb{N}\,\bigl(x\leq n\land
      \ExtMap{\iota_{A,A\cup\{a\}}\circ F}{n}{a}(x)=z\bigr)
    \end{aligned}
  }{%
    \begin{aligned}[t]
      &a\notin A\dsep n\in\mathbb{N}\dsep F\colon \NatSeg{n}\bij A\dsep z\in A\cup\{a\}\dsep z\in A\\
      &\vdash \exists x\in\mathbb{N}\,\bigl(x\leq n\land
      \ExtMap{\iota_{A,A\cup\{a\}}\circ F}{n}{a}(x)=z\bigr)
    \end{aligned}
  }
    \proofstepwidestar[1]{a\notin A}{\rA}
    \proofstepwidestar[2]{n\in\mathbb{N}}{\rA}
    \proofstepwidestar[3]{F\colon \NatSeg{n}\bij A}{\rA}
    \proofstepwidestar[4]{z\in A\cup\{a\}}{\rA}
    \proofstepwidestar[5]{z\in A}{\rA}
    \proofstepwidestar[2,3,5]{%
      \exists x\in\mathbb{N}\,\bigl(x<n\land
      \ExtMap{\iota_{A,A\cup\{a\}}\circ F}{n}{a}(x)=z\bigr)}{%
      \FormulaRefAuto{%
        n\in\mathbb{N}\dsep F\colon \NatSeg{n}\bij A\dsep z\in A \vdash
        \exists x\in\mathbb{N}\,\bigl(x<n\land
        \ExtMap{\iota_{A,A\cup\{a\}}\circ F}{n}{a}(x)=z\bigr)}{2,3,5}}

    \proofstepwidestar[7]{%
      x\in\mathbb{N}\land\bigl(x<n\land
      \ExtMap{\iota_{A,A\cup\{a\}}\circ F}{n}{a}(x)=z\bigr)}{\rA}
    \proofstepwidestar[7]{x\in\mathbb{N}}{\rAEa{7}}
    \proofstepwidestar[7]{x<n\land
      \ExtMap{\iota_{A,A\cup\{a\}}\circ F}{n}{a}(x)=z}{\rAEb{7}}
    \proofstepwidestar[7]{x<n}{\rAEa{9}}
    \proofstepwidestar[7]{\ExtMap{\iota_{A,A\cup\{a\}}\circ F}{n}{a}(x)=z}{\rAEb{9}}
    \proofstepwidestar[2,7]{x\leq n}{%
      \FormulaRefAuto{m\in\mathbb{N}\dsep n\in\mathbb{N}\dsep m<n\vdash m\leq n}{8,2,10}}
    \proofstepwidestar[2,7]{x\leq n\land
      \ExtMap{\iota_{A,A\cup\{a\}}\circ F}{n}{a}(x)=z}{\rAI{12,11}}
    \proofstepwidestar[2,7]{%
      \exists u\in\mathbb{N}\,\bigl(u\leq n\land
      \ExtMap{\iota_{A,A\cup\{a\}}\circ F}{n}{a}(u)=z\bigr)}{%
      \rEI{\rAI{8,13}}}
    \proofstepwidestar[2,3,5]{%
      \exists u\in\mathbb{N}\,\bigl(u\leq n\land
      \ExtMap{\iota_{A,A\cup\{a\}}\circ F}{n}{a}(u)=z\bigr)}{%
      \rEE{6,7,14}}
  \closeproofpartwideindR

  \proofpartwideindR[Fall $z=a$]{%
    \begin{aligned}[t]
      &a\notin A\dsep n\in\mathbb{N}\dsep F\colon \NatSeg{n}\bij A\dsep z\in A\cup\{a\}\dsep z=a\\
      &\vdash \exists x\in\mathbb{N}\,\bigl(x\leq n\land
      \ExtMap{\iota_{A,A\cup\{a\}}\circ F}{n}{a}(x)=z\bigr)
    \end{aligned}
  }{%
    \begin{aligned}[t]
      &a\notin A\dsep n\in\mathbb{N}\dsep F\colon \NatSeg{n}\bij A\dsep z\in A\cup\{a\}\dsep z=a\\
      &\vdash \exists x\in\mathbb{N}\,\bigl(x\leq n\land
      \ExtMap{\iota_{A,A\cup\{a\}}\circ F}{n}{a}(x)=z\bigr)
    \end{aligned}
  }
    \proofstepwidestar[1]{a\notin A}{\rA}
    \proofstepwidestar[2]{n\in\mathbb{N}}{\rA}
    \proofstepwidestar[3]{F\colon \NatSeg{n}\bij A}{\rA}
    \proofstepwidestar[4]{z\in A\cup\{a\}}{\rA}
    \proofstepwidestar[5]{z=a}{\rA}
    \proofstepwidestar[2,3,5]{%
      \exists x\in\mathbb{N}\,\bigl(x\leq n\land
      \ExtMap{\iota_{A,A\cup\{a\}}\circ F}{n}{a}(x)=z\bigr)}{%
      \FormulaRefAuto{%
        n\in\mathbb{N}\dsep F\colon \NatSeg{n}\bij A\dsep z=a \vdash
        \exists x\in\mathbb{N}\,\bigl(x\leq n\land
        \ExtMap{\iota_{A,A\cup\{a\}}\circ F}{n}{a}(x)=z\bigr)}{2,3,5}}
  \closeproofpartwideindR

  \proofpartwide{Schluss}
    \proofstepwidestar[1]{a\notin A}{\rA}
    \proofstepwidestar[2]{n\in\mathbb{N}}{\rA}
    \proofstepwidestar[3]{F\colon \NatSeg{n}\bij A}{\rA}
    \proofstepwidestar[4]{z\in A\cup\{a\}}{\rA}

    \proofstepwidestar[4]{z\in A\lor z=a}{%
      \FormulaRefAuto{x\in A\cup\{a\} \vdash x\in A\lor x=a}{4}}

    \proofstepwidestar[6]{z\in A}{\rA}
    \proofstepwidestar[1,2,3,4,6]{%
      \exists x\in\mathbb{N}\,\bigl(x\leq n\land
      \ExtMap{\iota_{A,A\cup\{a\}}\circ F}{n}{a}(x)=z\bigr)}{%
      \FormulaRefAuto{%
        \begin{aligned}[t]
          &a\notin A\dsep n\in\mathbb{N}\dsep F\colon \NatSeg{n}\bij A\dsep z\in A\cup\{a\}\dsep z\in A\\
          &\vdash \exists x\in\mathbb{N}\,\bigl(x\leq n\land
          \ExtMap{\iota_{A,A\cup\{a\}}\circ F}{n}{a}(x)=z\bigr)
        \end{aligned}
      }{1,2,3,4,6}}

    \proofstepwidestar[8]{z=a}{\rA}
    \proofstepwidestar[1,2,3,4,8]{%
      \exists x\in\mathbb{N}\,\bigl(x\leq n\land
      \ExtMap{\iota_{A,A\cup\{a\}}\circ F}{n}{a}(x)=z\bigr)}{%
      \FormulaRefAuto{%
        \begin{aligned}[t]
          &a\notin A\dsep n\in\mathbb{N}\dsep F\colon \NatSeg{n}\bij A\dsep z\in A\cup\{a\}\dsep z=a\\
          &\vdash \exists x\in\mathbb{N}\,\bigl(x\leq n\land
          \ExtMap{\iota_{A,A\cup\{a\}}\circ F}{n}{a}(x)=z\bigr)
        \end{aligned}
      }{1,2,3,4,8}}

    \proofstepwidestar[1,2,3,4]{%
      \exists x\in\mathbb{N}\,\bigl(x\leq n\land
      \ExtMap{\iota_{A,A\cup\{a\}}\circ F}{n}{a}(x)=z\bigr)}{%
      \rOE{5,6,7,8,9}}
  \closeproofpartwide
\end{tabproofsplitwide}

\paragraph{Bijektivität}
\FormulaThmDelta[Adjunktion einer Bijektion]{%
  a\notin A \dsep n\in\mathbb{N}\dsep F\colon \NatSeg{n}\bij A \vdash
  \ExtMap{\iota_{A,A\cup\{a\}}\circ F}{n}{a}\colon (\NatSeg{n}\cup\{n\})\bij (A\cup\{a\})
}{%
  \DeltaRow{Mengen}{A\dsep a\dsep n\dsep x\dsep y\dsep z}
  \DeltaRow{Funktionen}{F}
}
\begin{tabproof}
  \proofstep{1}{a\notin A}{\rA}
  \proofstep{2}{n\in\mathbb{N}}{\rA}
  \proofstep{3}{F\colon \NatSeg{n}\bij A}{\rA}
  \proofstep{2,3}{\ExtMap{\iota_{A,A\cup\{a\}}\circ F}{n}{a}\colon (\NatSeg{n}\cup\{n\})\to A\cup\{a\}}{%
    \FormulaRefAuto{n\in\mathbb{N}\dsep F\colon \NatSeg{n}\bij A \vdash \ExtMap{\iota_{A,A\cup\{a\}}\circ F}{n}{a}\colon (\NatSeg{n}\cup\{n\})\to A\cup\{a\}}{2,3}}
  \proofstep{1,2,3}{\ExtMap{\iota_{A,A\cup\{a\}}\circ F}{n}{a}\colon (\NatSeg{n}\cup\{n\})\inj (A\cup\{a\})}{%
    \FormulaRefAuto{Injektive Funktion}{4,\FormulaRefAuto{%
      \begin{aligned}[t]
        &a\notin A \dsep n\in\mathbb{N}\dsep F\colon \NatSeg{n}\bij A\dsep{}\\
        &x\in \NatSeg{n}\cup\{n\}\dsep y\in \NatSeg{n}\cup\{n\}\dsep{}\\
        &\ExtMap{\iota_{A,A\cup\{a\}}\circ F}{n}{a}(x)= {}\\
        &\qquad \ExtMap{\iota_{A,A\cup\{a\}}\circ F}{n}{a}(y)\vdash x=y
      \end{aligned}
    }{1,2,3}}}
  \proofstep{1,2,3}{\ExtMap{\iota_{A,A\cup\{a\}}\circ F}{n}{a}\colon (\NatSeg{n}\cup\{n\})\sur (A\cup\{a\})}{%
    \FormulaRefAuto{Surjektive Funktion}{4,\FormulaRefAuto{%
      \begin{aligned}[t]
        &a\notin A \dsep n\in\mathbb{N}\dsep F\colon \NatSeg{n}\bij A\dsep z\in A\cup\{a\}\\
        &\vdash \exists x\in \NatSeg{n}\cup\{n\}\,
          \ExtMap{\iota_{A,A\cup\{a\}}\circ F}{n}{a}(x)=z
      \end{aligned}
    }{1,2,3}}}
  \proofstep{1,2,3}{%
    \ExtMap{\iota_{A,A\cup\{a\}}\circ F}{n}{a}\colon (\NatSeg{n}\cup\{n\})\bij (A\cup\{a\})
  }{%
    \FormulaRefAuto{F\colon A\inj B\dsep F\colon A\sur B \vdash F\colon A\bij B}{5,6}}
\end{tabproof}

\FormulaThmDelta[Adjunktion einer Bijektion auf der Nachfolgermenge]{%
  a\notin A \dsep n\in\mathbb{N}\dsep F\colon \NatSeg{n}\bij A \vdash
  \ExtMap{\iota_{A,A\cup\{a\}}\circ F}{n}{a}\colon \NatSeg{\Succ(n)}\bij (A\cup\{a\})
}{%
  \DeltaRow{Mengen}{A\dsep a\dsep n}
  \DeltaRow{Funktionen}{F}
}
\begin{tabproof}
  \proofstep{1}{a\notin A}{\rA}
  \proofstep{2}{n\in\mathbb{N}}{\rA}
  \proofstep{3}{F\colon \NatSeg{n}\bij A}{\rA}
  \proofstep{1,2,3}{\ExtMap{\iota_{A,A\cup\{a\}}\circ F}{n}{a}\colon (\NatSeg{n}\cup\{n\})\bij (A\cup\{a\})}{%
    \FormulaRefAuto{a\notin A \dsep n\in\mathbb{N}\dsep F\colon \NatSeg{n}\bij A \vdash \ExtMap{\iota_{A,A\cup\{a\}}\circ F}{n}{a}\colon (\NatSeg{n}\cup\{n\})\bij (A\cup\{a\})}{1,2,3}}
  \proofstep{2}{\NatSeg{\Succ(n)}=\NatSeg{n}\cup\{n\}}{%
    \FormulaRefAuto{PeanoNatSegSucc}{2}}
  \proofstep{2}{\NatSeg{n}\cup\{n\}=\NatSeg{\Succ(n)}}{\FormulaRefAuto{a=b\vdash b=a}{5}}
  \proofstep{1,2,3}{\ExtMap{\iota_{A,A\cup\{a\}}\circ F}{n}{a}\colon \NatSeg{\Succ(n)}\bij (A\cup\{a\})}{\rIE{6,4}}
\end{tabproof}

\paragraph{Erhaltung der Endlichkeit}

\iffalse
\FormulaThmDelta[Adjunktion eines neuen Elements erhält Endlichkeit]{%
  a\notin A \dsep \Finite{A} \vdash \Finite{A\cup\{a\}}
}{%
  \DeltaRow{Mengen}{A\dsep a\dsep n\dsep m}
  \DeltaRow{Funktionen}{F\dsep G}
}
\begin{tabproof}
  \proofstep{1}{a\notin A}{\rA}
  \proofstep{2}{\Finite{A}}{\rA}

  \proofstep{2}{\exists n\in\mathbb{N}\,\exists F\colon \NatSeg{n}\bij A}{%
    \text{inaktiver Altbeweis}}

  \proofstep{4}{n\in\mathbb{N}\land \exists F\colon \NatSeg{n}\bij A}{\rA}
  \proofstep{4}{n\in\mathbb{N}}{\rAEa{4}}
  \proofstep{4}{\exists F\colon \NatSeg{n}\bij A}{\rAEb{4}}

  \proofstep{7}{F\colon \NatSeg{n}\bij A}{\rA}

  \proofstep{1,5,7}{\ExtMap{\iota_{A,A\cup\{a\}}\circ F}{n}{a}\colon \NatSeg{\Succ(n)}\bij (A\cup\{a\})}{%
    \FormulaRefAuto{a\notin A \dsep n\in\mathbb{N}\dsep F\colon \NatSeg{n}\bij A \vdash \ExtMap{\iota_{A,A\cup\{a\}}\circ F}{n}{a}\colon \NatSeg{\Succ(n)}\bij (A\cup\{a\})}{1,5,7}}

  \proofstep{1,5,7}{\exists G\colon \NatSeg{\Succ(n)}\bij (A\cup\{a\})}{\rEI{8}}
  \proofstep{1,5,7}{\NatSeg{\Succ(n)} \EqCard (A\cup\{a\})}{%
    \FormulaRefAuto{A \EqCard B \coloneqq \exists F\colon A\bij B}{9}}
  \proofstep{5}{\Succ(n)\in\mathbb{N}}{%
    \FormulaRefAuto{PeanoSuccClosure}{5}}

  \proofstep{1,5,7}{\exists m\in \mathbb{N}\,\NatSeg{m} \EqCard (A\cup\{a\})}{%
    \rEI{\rAI{11,10}}}
  \proofstep{1,5,7}{\Finite{A\cup\{a\}}}{%
    \FormulaRefAuto{\Finite{M}\coloneqq \exists n\in\mathbb{N}\,\bigl(\NatSeg{n} \EqCard M\bigr)}{12}}

  \proofstep{1,4}{\Finite{A\cup\{a\}}}{\rEE{6,7,13}}
  \proofstep{1,2}{\Finite{A\cup\{a\}}}{\rEE{3,4,14}}
\end{tabproof}
\fi

\fi

\subsection{Adjunktion endlich vieler Elemente}

\iffalse

\FormulaThmDeltaR[Zerlegung des Bildes einer Nachfolgermenge]{%
  \begin{aligned}[t]
    &n\in\mathbb{N}\dsep F\colon \NatSeg{\Succ(n)}\to M \vdash{}\\
    &F[\NatSeg{\Succ(n)}]=(F\restriction_{\NatSeg{n}})[\NatSeg{n}]\cup\{F(n)\}
  \end{aligned}
}{%
  n\in\mathbb{N}\dsep F\colon \NatSeg{\Succ(n)}\to M \vdash F[\NatSeg{\Succ(n)}]=(F\restriction_{\NatSeg{n}})[\NatSeg{n}]\cup\{F(n)\}
}{%
  \DeltaRow{Mengen}{M\dsep n\dsep y}
  \DeltaRow{Funktionen}{F}
}
\begin{tabproofwide}
  \proofstepwidestar[1]{n\in\mathbb{N}}{\rA}
  \proofstepwidestar[2]{F\colon \NatSeg{\Succ(n)}\to M}{\rA}
  \proofstepwidestar[1]{\NatSeg{\Succ(n)}=\NatSeg{n}\cup\{n\}}{%
    \FormulaRefAuto{PeanoNatSegSucc}{1}}
  \proofstepwidestar[1]{\NatSeg{n}\subseteq \NatSeg{\Succ(n)}}{%
    \FormulaRefAuto{n\in\mathbb{N}\vdash \NatSeg{n}\subseteq \NatSeg{\Succ(n)}}{1}}
  \proofstepwidestar[1]{n\in \NatSeg{\Succ(n)}}{%
    \FormulaRefAuto{n\in\mathbb{N}\vdash n\in \NatSeg{\Succ(n)}}{1}}
  \proofstepwidestar[1]{\{n\}\subseteq \NatSeg{\Succ(n)}}{%
    \FormulaRefAuto{x\in A \vdash \{x\}\subseteq A}{5}}
  \proofstepwidestar[1,2]{F[\NatSeg{n}\cup\{n\}]=F[\NatSeg{n}]\cup F[\{n\}]}{%
    \FormulaRefAuto{F\colon M\to N \dsep A\subseteq M \dsep B\subseteq M \vdash F[A\cup B]=F[A]\cup F[B]}{2,4,6}}
  \proofstepwidestar[1,2]{(F\restriction_{\NatSeg{n}})[\NatSeg{n}]=F[\NatSeg{n}]}{%
    \FormulaRefAuto{C\subseteq A \dsep F\colon A\to B \vdash (F\restriction_C)[C]=F[C]}{4,2}}
  \proofstepwidestar[1,2]{F[\{n\}]=\{F(n)\}}{%
    \FormulaRefAuto{x\in A \dsep F\colon A\to B \vdash F[\{x\}]=\{F(x)\}}{5,2}}
  \proofstepwidestar[1,2]{F[\NatSeg{n}]=(F\restriction_{\NatSeg{n}})[\NatSeg{n}]}{%
    \FormulaRefAuto{a=b\vdash b=a}{8}}

  \proofstepwide[]{F[\NatSeg{\Succ(n)}]}{=}{F[\NatSeg{\Succ(n)}]}{\rII}
  \proofstepwide[1]{}{=}{F[\NatSeg{n}\cup\{n\}]}{\rIE{3}}
  \proofstepwide[1,2]{}{=}{F[\NatSeg{n}]\cup F[\{n\}]}{\rIE{7}}
  \proofstepwide[1,2]{}{=}{(F\restriction_{\NatSeg{n}})[\NatSeg{n}]\cup F[\{n\}]}{\rIE{10}}
  \proofstepwide[1,2]{}{=}{(F\restriction_{\NatSeg{n}})[\NatSeg{n}]\cup\{F(n)\}}{\rIE{9}}

  \proofstepwidestar[1,2]{F[\NatSeg{\Succ(n)}]=(F\restriction_{\NatSeg{n}})[\NatSeg{n}]\cup\{F(n)\}}{\rChain{11,12,13,14,15}}
\end{tabproofwide}

\FormulaThmDeltaR[Hilfssatz zum Induktionsanfang der Adjunktion endlich vieler Elemente]{%
  \Finite{A}\dsep F\colon \varnothing\to M \vdash \Finite{A\cup F[\varnothing]}
}{%
  \Finite{A}\dsep F\colon \varnothing\to M \vdash \Finite{A\cup F[\varnothing]}
}{%
  \DeltaRow{Mengen}{A\dsep M}
  \DeltaRow{Funktionen}{F}
}
\begin{tabproofwide}
  \proofstepwidestar[1]{\Finite{A}}{\rA}
  \proofstepwidestar[2]{F\colon \varnothing\to M}{\rA}
  \proofstepwidestar[2]{F[\varnothing]=\varnothing}{%
    \FormulaRefAuto{F\colon A\to B \vdash F[\varnothing]=\varnothing}{2}}

  \proofstepwide[2]{A\cup F[\varnothing]}{=}{A\cup\varnothing}{\rIE{3}}
  \proofstepwide[]{}{=}{\varnothing\cup A}{%
    \FormulaRefAuto{A \cup B = B \cup A}}
  \proofstepwide[]{}{=}{A}{%
    \FormulaRefAuto{a=b\vdash b=a}{\FormulaRefAuto{A = \varnothing \cup A}}}
  \proofstepwidestar[2]{A\cup F[\varnothing]=A}{\rChain{4,5,6}}

  \proofstepwidestar[1,2]{\Finite{A\cup F[\varnothing]}}{\rIE{7,1}}
\end{tabproofwide}

\FormulaThmDeltaR[Hilfssatz zum Induktionsschritt der Adjunktion endlich vieler Elemente]{%
  \Finite{A}\dsep n\in\mathbb{N}\dsep F\colon \NatSeg{\Succ(n)}\to M\dsep \Finite{A\cup (F\restriction_{\NatSeg{n}})[\NatSeg{n}]} \vdash \Finite{A\cup F[\NatSeg{\Succ(n)}]}
}{%
  \Finite{A}\dsep n\in\mathbb{N}\dsep F\colon \NatSeg{\Succ(n)}\to M\dsep \Finite{A\cup (F\restriction_{\NatSeg{n}})[\NatSeg{n}]} \vdash \Finite{A\cup F[\NatSeg{\Succ(n)}]}
}{%
  \DeltaRow{Mengen}{A\dsep M\dsep n}
  \DeltaRow{Funktionen}{F}
}
\begin{tabproofwide}
  \proofstepwidestar[1]{\Finite{A}}{\rA}
  \proofstepwidestar[2]{n\in\mathbb{N}}{\rA}
  \proofstepwidestar[3]{F\colon \NatSeg{\Succ(n)}\to M}{\rA}
  \proofstepwidestar[4]{\Finite{A\cup (F\restriction_{\NatSeg{n}})[\NatSeg{n}]}}{\rA}

  \proofstepwidestar[2,3]{F[\NatSeg{\Succ(n)}]=(F\restriction_{\NatSeg{n}})[\NatSeg{n}]\cup\{F(n)\}}{%
    \FormulaRefAuto{n\in\mathbb{N}\dsep F\colon \NatSeg{\Succ(n)}\to M \vdash F[\NatSeg{\Succ(n)}]=(F\restriction_{\NatSeg{n}})[\NatSeg{n}]\cup\{F(n)\}}{2,3}}

  \proofstepwide[2,3]{A\cup F[\NatSeg{\Succ(n)}]}{=}{A\cup\bigl((F\restriction_{\NatSeg{n}})[\NatSeg{n}]\cup\{F(n)\}\bigr)}{\rIE{5}}
  \proofstepwide[]{}{=}{(A\cup(F\restriction_{\NatSeg{n}})[\NatSeg{n}])\cup\{F(n)\}}{%
    \FormulaRefAuto{A\cup (B\cup C) = (A\cup B)\cup C}}
  \proofstepwidestar[2,3]{A\cup F[\NatSeg{\Succ(n)}]=(A\cup(F\restriction_{\NatSeg{n}})[\NatSeg{n}])\cup\{F(n)\}}{\rChain{6,7}}

  \proofstepwidestar[]{F(n)\in A\cup(F\restriction_{\NatSeg{n}})[\NatSeg{n}]\lor F(n)\notin A\cup(F\restriction_{\NatSeg{n}})[\NatSeg{n}]}{\FormulaRefAuto{P\lor \neg P}}

  \proofstepwidestar[10]{F(n)\in A\cup(F\restriction_{\NatSeg{n}})[\NatSeg{n}]}{\rA}
  \proofstepwidestar[10]{(A\cup(F\restriction_{\NatSeg{n}})[\NatSeg{n}])\cup\{F(n)\}=A\cup(F\restriction_{\NatSeg{n}})[\NatSeg{n}]}{%
    \FormulaRefAuto{A\cup \{a\}=A\eqvdash a\in A}{10}}
  \proofstepwidestar[4,10]{%
    \Finite{(A\cup(F\restriction_{\NatSeg{n}})[\NatSeg{n}])\cup\{F(n)\}}
  }{\rIE{11,4}}

  \proofstepwidestar[13]{F(n)\notin A\cup(F\restriction_{\NatSeg{n}})[\NatSeg{n}]}{\rA}
  \proofstepwidestar[4,13]{%
    \Finite{(A\cup(F\restriction_{\NatSeg{n}})[\NatSeg{n}])\cup\{F(n)\}}
  }{%
    \FormulaRefAuto{FiniteAdjunction}[thm]{13,4}}

  \proofstepwidestar[4]{%
    \Finite{(A\cup(F\restriction_{\NatSeg{n}})[\NatSeg{n}])\cup\{F(n)\}}
  }{\rOE{9,10,12,13,14}}
  \proofstepwidestar[1,2,3,4]{\Finite{A\cup F[\NatSeg{\Succ(n)}]}}{\rIE{8,15}}
\end{tabproofwide}

\FormulaThmDeltaR[Adjunktion endlich vieler Elemente]{%
  \begin{aligned}[t]
    &\Finite{A}\dsep n\in\mathbb{N}\dsep F\colon \NatSeg{n}\to M \vdash{}\\
    &\Finite{A\cup F[\NatSeg{n}]}
  \end{aligned}
}{%
  \Finite{A}\dsep n\in\mathbb{N}\dsep F\colon \NatSeg{n}\to M \vdash \Finite{A\cup F[\NatSeg{n}]}
}{%
  \DeltaRow{Mengen}{A\dsep M\dsep n}
  \DeltaRow{Funktionen}{F\dsep G}
}
\begin{tabproofsplitwide}
  \proofpartwideind[Induktionsanfang]{%
    \Finite{A}\vdash \forall F\,\bigl(F\colon \varnothing\to M \rightarrow \Finite{A\cup F[\varnothing]}\bigr)}
    \proofstepwidestar[1]{\Finite{A}}{\rA}
    \proofstepwidestar[2]{F\colon \varnothing\to M}{\rA}
    \proofstepwidestar[1,2]{\Finite{A\cup F[\varnothing]}}{%
      \FormulaRefAuto{\Finite{A}\dsep F\colon \varnothing\to M \vdash \Finite{A\cup F[\varnothing]}}{1,2}}
    \proofstepwidestar[1]{F\colon \varnothing\to M \rightarrow \Finite{A\cup F[\varnothing]}}{\rII{2,3}}
    \proofstepwidestar[1]{\forall F\,\bigl(F\colon \varnothing\to M \rightarrow \Finite{A\cup F[\varnothing]}\bigr)}{\rUI{\rRI{2,4}}}
  \closeproofpartwideind

  \proofpartwideindR[Induktionsschritt]{%
    \begin{aligned}[t]
      &\Finite{A}\dsep n\in\mathbb{N}\dsep
      \forall G\,\bigl(G\colon \NatSeg{n}\to M \rightarrow \Finite{A\cup G[\NatSeg{n}]}\bigr)\vdash{}\\
      &\forall F\,\bigl(F\colon \NatSeg{\Succ(n)}\to M \rightarrow \Finite{A\cup F[\NatSeg{\Succ(n)}]}\bigr)
    \end{aligned}
  }{%
    \Finite{A}\dsep n\in\mathbb{N}\dsep \forall G\,\bigl(G\colon \NatSeg{n}\to M \rightarrow \Finite{A\cup G[\NatSeg{n}]}\bigr)\vdash \forall F\,\bigl(F\colon \NatSeg{\Succ(n)}\to M \rightarrow \Finite{A\cup F[\NatSeg{\Succ(n)}]}\bigr)
  }
    \proofstepwidestar[1]{\Finite{A}}{\rA}
    \proofstepwidestar[2]{n\in\mathbb{N}}{\rA}
    \proofstepwidestar[3]{\forall G\,\bigl(G\colon \NatSeg{n}\to M \rightarrow \Finite{A\cup G[\NatSeg{n}]}\bigr)}{\rA}
    \proofstepwidestar[4]{F\colon \NatSeg{\Succ(n)}\to M}{\rA}

    \proofstepwidestar[2]{\NatSeg{n}\subseteq \NatSeg{\Succ(n)}}{%
      \FormulaRefAuto{n\in\mathbb{N}\vdash \NatSeg{n}\subseteq \NatSeg{\Succ(n)}}{2}}
    \proofstepwidestar[2,4]{F\restriction_{\NatSeg{n}}\colon \NatSeg{n}\to M}{%
      \FormulaRefAuto{C\subseteq A \vdash F\restriction_C\colon C\to B}{5,4}}
    \proofstepwidestar[3]{F\restriction_{\NatSeg{n}}\colon \NatSeg{n}\to M \rightarrow \Finite{A\cup (F\restriction_{\NatSeg{n}})[\NatSeg{n}]}}{\rUE{3}}
    \proofstepwidestar[1,2,3,4]{\Finite{A\cup (F\restriction_{\NatSeg{n}})[\NatSeg{n}]}}{\rIE{6,7}}

    \proofstepwidestar[1,2,4,8]{\Finite{A\cup F[\NatSeg{\Succ(n)}]}}{%
      \FormulaRefAuto{\Finite{A}\dsep n\in\mathbb{N}\dsep F\colon \NatSeg{\Succ(n)}\to M\dsep \Finite{A\cup (F\restriction_{\NatSeg{n}})[\NatSeg{n}]} \vdash \Finite{A\cup F[\NatSeg{\Succ(n)}]}}{1,2,4,8}}

    \proofstepwidestar[1,2,3]{F\colon \NatSeg{\Succ(n)}\to M \rightarrow \Finite{A\cup F[\NatSeg{\Succ(n)}]}}{\rII{4,9}}
    \proofstepwidestar[1,2,3]{\forall F\,\bigl(F\colon \NatSeg{\Succ(n)}\to M \rightarrow \Finite{A\cup F[\NatSeg{\Succ(n)}]}\bigr)}{\rUI{\rRI{4,10}}}
  \closeproofpartwideindR

  \proofpartwide{Induktionsschluss}
    \proofstepwidestar[1]{\Finite{A}}{\rA}
    \proofstepwidestar[2]{n\in\mathbb{N}}{\rA}
    \proofstepwidestar[3]{F\colon \NatSeg{n}\to M}{\rA}

    \proofstepwidestar[1,2]{\forall F\,\bigl(F\colon \NatSeg{n}\to M \rightarrow \Finite{A\cup F[\NatSeg{n}]}\bigr)}{%
      \FormulaRefAuto{PeanoInductionRule}{IA,IS,2}}
    \proofstepwidestar[1,2]{F\colon \NatSeg{n}\to M \rightarrow \Finite{A\cup F[\NatSeg{n}]}}{\rUE{4}}
    \proofstepwidestar[1,2,3]{\Finite{A\cup F[\NatSeg{n}]}}{\rIE{3,5}}
  \closeproofpartwide
\end{tabproofsplitwide}

\fi

\iffalse
\FormulaThmDeltaK[Adjunktion endlich vieler Elemente]{%
  \Finite{A}\dsep n\in\mathbb{N}\dsep
  F\colon \NatSeg{n}\to M
  \vdash \Finite{A\cup F[\NatSeg{n}]}%
}{FiniteUnionNatSegImage}{%
  \DeltaRow{Mengen}{A\dsep M\dsep n}%
  \DeltaRow{Funktionen}{F}%
}
\begin{tabproof}
  \proofstep{1}{\Finite{A}}{\rA}
  \proofstep{2}{n\in\mathbb{N}}{\rA}
  \proofstep{3}{F\colon \NatSeg{n}\to M}{\rA}
  \proofstep{2}{\Finite{\NatSeg{n}}}{%
    \FormulaRefAuto{PeanoNatSegFinite}[thm]{2}}
  \proofstep{}{\NatSeg{n}\subseteq\NatSeg{n}}{%
    \FormulaRefAuto{A\subseteq A}}
  \proofstep{2,3}{\Finite{F[\NatSeg{n}]}}{%
    \FormulaRefAuto{FiniteImage}[thm]{5,3,4}}
  \proofstep{1,2,3}{\Finite{A\cup F[\NatSeg{n}]}}{%
    \FormulaRefAuto{FiniteUnion}[thm]{1,6}}
\end{tabproof}
\fi

\subsection{Vereinigung endlicher Mengen ist endlich}

\iffalse
\FormulaThmDelta[Vereinigung endlicher Mengen ist endlich]{%
  \Finite{A}\dsep \Finite{B} \vdash \Finite{A\cup B}
}{%
  \DeltaRow{Mengen}{A\dsep B\dsep n}
  \DeltaRow{Funktionen}{F}
}
\begin{tabproof}
  \proofstep{1}{\Finite{A}}{\rA}
  \proofstep{2}{\Finite{B}}{\rA}

  \proofstep{2}{\exists n\in\mathbb{N}\,\exists F\colon \NatSeg{n}\bij B}{%
    \text{inaktiver Altbeweis}}

  \proofstep{4}{n\in\mathbb{N}\land \exists F\colon \NatSeg{n}\bij B}{\rA}
  \proofstep{4}{n\in\mathbb{N}}{\rAEa{4}}
  \proofstep{4}{\exists F\colon \NatSeg{n}\bij B}{\rAEb{4}}

  \proofstep{7}{F\colon \NatSeg{n}\bij B}{\rA}
  \proofstep{7}{F\colon \NatSeg{n}\to B}{\FormulaRefAuto{Bijektive Funktion}{7}}
  \proofstep{1,5,7}{\Finite{A\cup F[\NatSeg{n}]}}{%
    \FormulaRefAuto{\Finite{A}\dsep n\in\mathbb{N}\dsep F\colon \NatSeg{n}\to M \vdash \Finite{A\cup F[\NatSeg{n}]}}{1,5,8}}
  \proofstep{7}{F[\NatSeg{n}]=B}{%
    \FormulaRefAuto{F\colon A\bij B \vdash F[A]=B}{7}}
  \proofstep{7}{A\cup F[\NatSeg{n}]=A\cup B}{%
    \FormulaRefAuto{A=B\vdash C\cup A = C\cup B}{10}}
  \proofstep{1,5,7}{\Finite{A\cup B}}{\rIE{11,9}}

  \proofstep{1,4}{\Finite{A\cup B}}{\rEE{6,7,12}}
  \proofstep{1,2}{\Finite{A\cup B}}{\rEE{3,4,13}}
\end{tabproof}
\fi

\subsection{Ein- und Paarmengen}

\FormulaThmDeltaK[Einermengen sind endlich]{%
  \Finite{\{A\}}
}{FiniteSingleton}{%
  \DeltaRow{Mengen}{A}
}
\begin{tabproof}
  \proofstep{}{\Finite{\varnothing}}{\FormulaRefAuto{FiniteEmpty}[thm]}
  \proofstep{}{A\notin\varnothing}{\FormulaRefAuto{x \not\in \varnothing}}
  \proofstep{}{\Finite{\varnothing\cup\{A\}}}{%
    \FormulaRefAuto{FiniteAdjunction}[thm]{2,1}}
  \proofstep{}{\{A\}=\varnothing\cup\{A\}}{%
    \FormulaRefAuto{A = \varnothing \cup A}}
  \proofstep{}{\varnothing\cup\{A\}=\{A\}}{%
    \FormulaRefAuto{a=b\vdash b=a}{4}}
  \proofstep{}{\Finite{\{A\}}}{\rIE{5,3}}
\end{tabproof}

\FormulaThmDeltaK[Paarmengen sind endlich]{%
  \Finite{\{A,B\}}
}{FinitePair}{%
  \DeltaRow{Mengen}{A\dsep B}
}
\begin{tabproof}
  \proofstep{}{\Finite{\{A\}}}{\FormulaRefAuto{\Finite{\{A\}}}}
  \proofstep{}{\Finite{\{B\}}}{\FormulaRefAuto{\Finite{\{A\}}}}
  \proofstep{}{\Finite{\{A\}\cup\{B\}}}{%
    \FormulaRefAuto{FiniteUnion}[thm]{1,2}}
  \proofstep{}{\{A,B\}=\{A\}\cup\{B\}}{%
    \FormulaRefAuto{\{a,b\} = \{a\} \cup \{b\}}}
  \proofstep{}{\{A\}\cup\{B\}=\{A,B\}}{%
    \FormulaRefAuto{a=b\vdash b=a}{4}}
  \proofstep{}{\Finite{\{A,B\}}}{\rIE{5,3}}
\end{tabproof}

\subsection{Teilmengen natürlicher Zahlen sind endlich}

Die Induktion erfolgt einheitlich über die Peano-Anfangsabschnitte
\(\NatSeg{n}\); natürliche Zahlen werden dabei nicht als Mengen verwendet.

% Die leere Menge ist nach dem Endlichkeitsaxiom endlich.

\FormulaThmDelta[Induktionsanfang für endliche Teilmengen natürlicher Zahlen]{%
  C\subseteq \varnothing \vdash \Finite{C}
}{%
  \DeltaRow{Mengen}{C}
}
\begin{tabproof}
  \proofstep{1}{C\subseteq \varnothing}{\rA}
  \proofstep{1}{C=\varnothing}{\FormulaRefAuto{A\subseteq \varnothing \vdash A=\varnothing}{1}}
  \proofstep{}{\Finite{\varnothing}}{\FormulaRefAuto{FiniteEmpty}[thm]}
  \proofstep{1}{\varnothing=C}{\FormulaRefAuto{a=b \vdash b=a}{2}}
  \proofstep{1}{\Finite{C}}{\rIE{4,3}}
\end{tabproof}

\FormulaThmDelta[Fall \(n\in C\) im Induktionsschritt für endliche Teilmengen natürlicher Zahlen]{%
  \begin{aligned}[t]
    &n\in\mathbb{N}\dsep \forall C\,\bigl(C\subseteq \NatSeg{n} \rightarrow \Finite{C}\bigr)\dsep{}\\
    &C\subseteq \NatSeg{\Succ(n)}\dsep n\in C \vdash \Finite{C}
  \end{aligned}
}{%
  \DeltaRow{Mengen}{n\dsep C}
}
\begin{tabproof}
  \proofstep{1}{n\in\mathbb{N}}{\rA}
  \proofstep{2}{\forall C\,\bigl(C\subseteq \NatSeg{n} \rightarrow \Finite{C}\bigr)}{\rA}
  \proofstep{3}{C\subseteq \NatSeg{\Succ(n)}}{\rA}
  \proofstep{4}{n\in C}{\rA}

  \proofstep{}{C\cap \NatSeg{n}\subseteq \NatSeg{n}}{\FormulaRefAuto{A\cap B \subseteq B}}
  \proofstep{2}{(C\cap \NatSeg{n})\subseteq \NatSeg{n} \rightarrow \Finite{C\cap \NatSeg{n}}}{\rUE{2}}
  \proofstep{2}{\Finite{C\cap \NatSeg{n}}}{\rRE{5,6}}

  \proofstep{8}{n\in C\cap \NatSeg{n}}{\rA}
  \proofstep{8}{n\in \NatSeg{n}}{\FormulaRefAuto{x\in A\cap B\vdash x\in B}{8}}
  \proofstep{1,8}{\bot}{\rBI{9,\FormulaRefAuto{n\in\mathbb{N}\vdash n\notin \NatSeg{n}}{1}}}
  \proofstep{1}{n\notin C\cap \NatSeg{n}}{\rCI{8,10}}

  \proofstep{1,2,4}{\Finite{(C\cap \NatSeg{n})\cup\{n\}}}{%
    \FormulaRefAuto{FiniteAdjunction}[thm]{11,7}}
  \proofstep{1,3,4}{C=(C\cap \NatSeg{n})\cup\{n\}}{%
    \FormulaRefAuto{PeanoNatSegSubsetSuccWithEndpoint}{1,3,4}}
  \proofstep{1,3,4}{(C\cap \NatSeg{n})\cup\{n\}=C}{\FormulaRefAuto{a=b \vdash b=a}{13}}
  \proofstep{1,2,3,4}{\Finite{C}}{\rIE{14,12}}
\end{tabproof}

\FormulaThmDelta[Fall \(n\notin C\) im Induktionsschritt für endliche Teilmengen natürlicher Zahlen]{%
  \begin{aligned}[t]
    &n\in\mathbb{N}\dsep \forall C\,\bigl(C\subseteq \NatSeg{n} \rightarrow \Finite{C}\bigr)\dsep{}\\
    &C\subseteq \NatSeg{\Succ(n)}\dsep n\notin C \vdash \Finite{C}
  \end{aligned}
}{%
  \DeltaRow{Mengen}{n\dsep C}
}
\begin{tabproof}
  \proofstep{1}{n\in\mathbb{N}}{\rA}
  \proofstep{2}{\forall C\,\bigl(C\subseteq \NatSeg{n} \rightarrow \Finite{C}\bigr)}{\rA}
  \proofstep{3}{C\subseteq \NatSeg{\Succ(n)}}{\rA}
  \proofstep{4}{n\notin C}{\rA}

  \proofstep{1,3,4}{C\subseteq \NatSeg{n}}{%
    \FormulaRefAuto{n\in\mathbb{N}\dsep B\subseteq \NatSeg{\Succ(n)}\dsep n\notin B \vdash B\subseteq \NatSeg{n}}{1,3,4}}
  \proofstep{1,2,3,4}{C\subseteq \NatSeg{n} \rightarrow \Finite{C}}{\rUE{2}}
  \proofstep{1,2,3,4}{\Finite{C}}{\rIE{5,6}}
\end{tabproof}

\FormulaThmDelta[Induktionsschritt für endliche Teilmengen natürlicher Zahlen]{%
  n\in\mathbb{N}\dsep \forall C\,\bigl(C\subseteq \NatSeg{n} \rightarrow \Finite{C}\bigr)\dsep C\subseteq \NatSeg{\Succ(n)}\vdash
  \Finite{C}
}{%
  \DeltaRow{Mengen}{n\dsep C}
}
\begin{tabproofsplitwide}
  \proofpartwideindR[Fall \(n\in C\)]{%
    \begin{aligned}[t]
      &n\in\mathbb{N}\dsep \forall C\,\bigl(C\subseteq \NatSeg{n} \rightarrow \Finite{C}\bigr)\dsep{}\\
      &C\subseteq \NatSeg{\Succ(n)}\dsep n\in C \vdash \Finite{C}
    \end{aligned}
  }{%
    n\in\mathbb{N}\dsep \forall C\,\bigl(C\subseteq \NatSeg{n} \rightarrow \Finite{C}\bigr)\dsep
    C\subseteq \NatSeg{\Succ(n)}\dsep n\in C \vdash \Finite{C}
  }
    \proofstepwidestar[1]{n\in\mathbb{N}}{\rA}
    \proofstepwidestar[2]{\forall C\,\bigl(C\subseteq \NatSeg{n} \rightarrow \Finite{C}\bigr)}{\rA}
    \proofstepwidestar[3]{C\subseteq \NatSeg{\Succ(n)}}{\rA}
    \proofstepwidestar[4]{n\in C}{\rA}

    \proofstepwidestar[1,2,3,4]{\Finite{C}}{%
      \FormulaRefAuto{%
        \begin{aligned}[t]
          &n\in\mathbb{N}\dsep \forall C\,\bigl(C\subseteq \NatSeg{n} \rightarrow \Finite{C}\bigr)\dsep{}\\
          &C\subseteq \NatSeg{\Succ(n)}\dsep n\in C \vdash \Finite{C}
        \end{aligned}
      }{1,2,3,4}}
  \closeproofpartwideindR

  \proofpartwideindR[Fall \(n\notin C\)]{%
    \begin{aligned}[t]
      &n\in\mathbb{N}\dsep \forall C\,\bigl(C\subseteq \NatSeg{n} \rightarrow \Finite{C}\bigr)\dsep{}\\
      &C\subseteq \NatSeg{\Succ(n)}\dsep n\notin C \vdash \Finite{C}
    \end{aligned}
  }{%
    n\in\mathbb{N}\dsep \forall C\,\bigl(C\subseteq \NatSeg{n} \rightarrow \Finite{C}\bigr)\dsep
    C\subseteq \NatSeg{\Succ(n)}\dsep n\notin C \vdash \Finite{C}
  }
    \proofstepwidestar[1]{n\in\mathbb{N}}{\rA}
    \proofstepwidestar[2]{\forall C\,\bigl(C\subseteq \NatSeg{n} \rightarrow \Finite{C}\bigr)}{\rA}
    \proofstepwidestar[3]{C\subseteq \NatSeg{\Succ(n)}}{\rA}
    \proofstepwidestar[4]{n\notin C}{\rA}

    \proofstepwidestar[1,2,3,4]{\Finite{C}}{%
      \FormulaRefAuto{%
        \begin{aligned}[t]
          &n\in\mathbb{N}\dsep \forall C\,\bigl(C\subseteq \NatSeg{n} \rightarrow \Finite{C}\bigr)\dsep{}\\
          &C\subseteq \NatSeg{\Succ(n)}\dsep n\notin C \vdash \Finite{C}
        \end{aligned}
      }{1,2,3,4}}
  \closeproofpartwideindR

  \proofpartwide{Schluss}
    \proofstepwidestar[1]{n\in\mathbb{N}}{\rA}
    \proofstepwidestar[2]{\forall C\,\bigl(C\subseteq \NatSeg{n} \rightarrow \Finite{C}\bigr)}{\rA}

    \proofstepwidestar[3]{C\subseteq \NatSeg{\Succ(n)}}{\rA}
    \proofstepwidestar[]{n\in C \lor n\notin C}{\FormulaRefAuto{P \lor \neg P}}

    \proofstepwidestar[5]{n\in C}{\rA}
    \proofstepwidestar[1,2,3,5]{\Finite{C}}{%
      \FormulaRefAuto{%
        n\in\mathbb{N}\dsep \forall C\,\bigl(C\subseteq \NatSeg{n} \rightarrow \Finite{C}\bigr)\dsep
        C\subseteq \NatSeg{\Succ(n)}\dsep n\in C \vdash \Finite{C}
      }{1,2,3,5}}

    \proofstepwidestar[7]{n\notin C}{\rA}
    \proofstepwidestar[1,2,3,7]{\Finite{C}}{%
      \FormulaRefAuto{%
        n\in\mathbb{N}\dsep \forall C\,\bigl(C\subseteq \NatSeg{n} \rightarrow \Finite{C}\bigr)\dsep
        C\subseteq \NatSeg{\Succ(n)}\dsep n\notin C \vdash \Finite{C}
      }{1,2,3,7}}

    \proofstepwidestar[1,2,3]{\Finite{C}}{\rOE{4,5,6,7,8}}
  \closeproofpartwide
\end{tabproofsplitwide}

\FormulaThmDeltaK[Teilmengen natürlicher Zahlen sind endlich]{%
  n\in\mathbb{N}\dsep C\subseteq \NatSeg{n} \vdash \Finite{C}
}{FiniteNatSegSubset}{%
  \DeltaRow{Mengen}{n\dsep C}
}
\begin{tabproofsplitwide}
  \proofpartwideind[Induktionsanfang]{%
    \forall C\,\bigl(C\subseteq \NatSeg{0} \rightarrow \Finite{C}\bigr)}
    \proofstepwidestar[1]{C\subseteq \NatSeg{0}}{\rA}
    \proofstepwidestar[]{\NatSeg{0}=\varnothing}{%
      \FormulaRefAuto{PeanoNatSegZero}}
    \proofstepwidestar[1]{C\subseteq \varnothing}{\rIE{2,1}}
    \proofstepwidestar[1]{\Finite{C}}{%
      \FormulaRefAuto{C\subseteq \varnothing \vdash \Finite{C}}{3}}
    \proofstepwidestar[]{C\subseteq \NatSeg{0} \rightarrow \Finite{C}}{\rII{1,4}}
    \proofstepwidestar[]{%
      \forall C\,\bigl(C\subseteq \NatSeg{0} \rightarrow \Finite{C}\bigr)}{%
      \rUI{\rRI{1,5}}}
  \closeproofpartwideind

  \proofpartwideind[Induktionsschritt]{%
    n\in\mathbb{N}\dsep \forall C\,\bigl(C\subseteq \NatSeg{n} \rightarrow \Finite{C}\bigr)\vdash
    \forall C\,\bigl(C\subseteq \NatSeg{\Succ(n)} \rightarrow \Finite{C}\bigr)}
    \proofstepwidestar[1]{n\in\mathbb{N}}{\rA}
    \proofstepwidestar[2]{\forall C\,\bigl(C\subseteq \NatSeg{n} \rightarrow \Finite{C}\bigr)}{\rA}

    \proofstepwidestar[3]{C\subseteq \NatSeg{\Succ(n)}}{\rA}
    \proofstepwidestar[1,2,3]{\Finite{C}}{%
      \FormulaRefAuto{n\in\mathbb{N}\dsep \forall C\,\bigl(C\subseteq \NatSeg{n} \rightarrow \Finite{C}\bigr)\dsep C\subseteq \NatSeg{\Succ(n)}\vdash \Finite{C}}{1,2,3}}
    \proofstepwidestar[1,2]{C\subseteq \NatSeg{\Succ(n)}\rightarrow \Finite{C}}{\rII{3,4}}
    \proofstepwidestar[1,2]{%
      \forall C\,\bigl(C\subseteq \NatSeg{\Succ(n)} \rightarrow \Finite{C}\bigr)}{%
      \rUI{\rRI{3,5}}}
  \closeproofpartwideind

  \proofpartwide{Induktionsschluss}
    \proofstepwidestar[1]{n\in\mathbb{N}}{\rA}
    \proofstepwidestar[2]{C\subseteq \NatSeg{n}}{\rA}

    \proofstepwidestar[1]{%
      \forall C\,\bigl(C\subseteq \NatSeg{n} \rightarrow \Finite{C}\bigr)}{%
      \FormulaRefAuto{PeanoInductionRule}{%
        \FormulaRefAuto{\forall C\,\bigl(C\subseteq \NatSeg{0} \rightarrow \Finite{C}\bigr)},%
        \FormulaRefAuto{n\in\mathbb{N}\dsep \forall C\,\bigl(C\subseteq \NatSeg{n} \rightarrow \Finite{C}\bigr)\vdash \forall C\,\bigl(C\subseteq \NatSeg{\Succ(n)} \rightarrow \Finite{C}\bigr)},%
        1}}

    \proofstepwidestar[1]{C\subseteq \NatSeg{n} \rightarrow \Finite{C}}{\rUE{3}}
    \proofstepwidestar[1,2]{\Finite{C}}{\rIE{2,4}}
  \closeproofpartwide
\end{tabproofsplitwide}

\FormulaThmDelta[Durch eine strikte natuerliche Schranke begrenzte Teilmengen sind endlich]{%
  n\in\mathbb{N}\dsep C\subseteq\mathbb{N}\dsep
  \forall x\in C\,\bigl(x<n\bigr)\vdash \Finite{C}
}{%
  \DeltaRow{Mengen}{n\dsep C\dsep x}
  \DeltaRow{Relationen}{<_{\mathbb{N}}}
}
\begin{tabproof}
  \proofstep{1}{n\in\mathbb{N}}{\rA}
  \proofstep{2}{C\subseteq\mathbb{N}}{\rA}
  \proofstep{3}{\forall x\in C\,\bigl(x<n\bigr)}{\rA}
  \proofstep{4}{x\in C}{\rA}
  \proofstep{2,4}{x\in\mathbb{N}}{%
    \FormulaRefAuto{A\subseteq B,\, x\in A\vdash x\in B}{2,4}}
  \proofstep{3,4}{x<n}{\rRE{4,\rUE{3}}}
  \proofstep{1,2,3,4}{x\in \NatSeg{n}}{%
    \FormulaRefAuto{PeanoNatSegStrictMembership}{5,1,6}}
  \proofstep{1,2,3}{\forall x\,\bigl(x\in C \rightarrow x\in \NatSeg{n}\bigr)}{\rUI{\rRI{4,7}}}
  \proofstep{1,2,3}{C\subseteq \NatSeg{n}}{%
    \FormulaRefAuto{A \subseteq B \coloneq \forall x\,(x\in A \rightarrow x\in B)}{8}}
  \proofstep{1,2,3}{\Finite{C}}{%
    \FormulaRefAuto{FiniteNatSegSubset}[thm]{1,9}}
\end{tabproof}

\FormulaThmDelta[Durch eine natuerliche Schranke begrenzte Teilmengen sind endlich]{%
  n\in\mathbb{N}\dsep C\subseteq\mathbb{N}\dsep
  \forall x\in C\,\bigl(x\leq n\bigr)\vdash \Finite{C}
}{%
  \DeltaRow{Mengen}{n\dsep C\dsep x}
  \DeltaRow{Relationen}{\leq_{\mathbb{N}}}
}
\begin{tabproof}
  \proofstep{1}{n\in\mathbb{N}}{\rA}
  \proofstep{2}{C\subseteq\mathbb{N}}{\rA}
  \proofstep{3}{\forall x\in C\,\bigl(x\leq n\bigr)}{\rA}
  \proofstep{4}{x\in C}{\rA}
  \proofstep{2,4}{x\in\mathbb{N}}{%
    \FormulaRefAuto{A\subseteq B,\, x\in A\vdash x\in B}{2,4}}
  \proofstep{3,4}{x\leq n}{\rRE{4,\rUE{3}}}
  \proofstep{1,2,3,4}{x\in\NatSeg{\Succ(n)}}{%
    \FormulaRefAuto{PeanoNatSegLeqSuccMembership}{5,1,6}}
  \proofstep{1,2,3}{\forall x\,\bigl(x\in C \rightarrow x\in\NatSeg{\Succ(n)}\bigr)}{\rUI{\rRI{4,7}}}
  \proofstep{1,2,3}{C\subseteq\NatSeg{\Succ(n)}}{%
    \FormulaRefAuto{A \subseteq B \coloneq \forall x\,(x\in A \rightarrow x\in B)}{8}}
  \proofstep{1}{\Succ(n)\in\mathbb{N}}{%
    \FormulaRefAuto{PeanoSuccClosure}{1}}
  \proofstep{1,2,3}{\Finite{C}}{%
    \FormulaRefAuto{FiniteNatSegSubset}[thm]{10,9}}
\end{tabproof}


\subsection{Teilmengen endlicher Mengen sind endlich}

\FormulaThmDeltaK[Teilmengen endlicher Mengen sind endlich]{%
  T\subseteq M\dsep \Finite{M} \vdash \Finite{T}
}{FiniteSubset}{%
  \DeltaRow{Mengen}{T\dsep M\dsep n}
  \DeltaRow{Funktionen}{F}
}
\begin{tabproof}
  \proofstep{1}{T\subseteq M}{\rA}
  \proofstep{2}{\Finite{M}}{\rA}
  \proofstep{2}{\exists n\in\mathbb{N}\,(\NatSeg{n}\EqCard M)}{%
    \FormulaRefAuto{FiniteEqCardEquiv}[thm]{2}}
  \proofstep{4}{n\in\mathbb{N}\land \NatSeg{n}\EqCard M}{\rA}
  \proofstep{4}{n\in\mathbb{N}}{\rAEa{4}}
  \proofstep{4}{\NatSeg{n}\EqCard M}{\rAEb{4}}
  \proofstep{4}{\exists F\colon \NatSeg{n}\bij M}{%
    \FormulaRefAuto{A\EqCard B\coloneqq \exists F\colon A\bij B}{6}}
  \proofstep{8}{F\colon \NatSeg{n}\bij M}{\rA}
  \proofstep{1,8}{F^{-1}[T]\subseteq \NatSeg{n}}{%
    \FormulaRefAuto{C\subseteq B\dsep F\colon A\bij B \vdash F^{-1}[C]\subseteq A}{1,8}}
  \proofstep{1,4,8}{\Finite{F^{-1}[T]}}{%
    \FormulaRefAuto{FiniteNatSegSubset}[thm]{5,9}}
  \proofstep{1,8}{F\restriction_{F^{-1}[T]}\colon F^{-1}[T]\bij T}{%
    \FormulaRefAuto{F\colon A\bij B \dsep T\subseteq B \vdash F\restriction_{F^{-1}[T]}\colon F^{-1}[T]\bij T}{8,1}}
  \proofstep{1,4,8}{\Finite{T}}{%
    \FormulaRefAuto{FiniteBijectionTransport}[thm]{10,11}}
  \proofstep{1,4}{\Finite{T}}{\rEE{7,8,12}}
  \proofstep{1,2}{\Finite{T}}{\rEE{3,4,13}}
\end{tabproof}

\subsection{Bilder endlicher Mengen sind endlich}

Dies ist die bereits axiomatisch bewiesene Aussage
\FormulaRefAuto{FiniteImage}[thm].

\subsection{Potenzmengen natürlicher Zahlen sind endlich}

Die Potenzmengeninduktion verwendet ausschliesslich die Träger
\(\NatSeg{n}\) und deren Nachfolgerzerlegung.

\FormulaDefDelta[Funktionsvorschrift der Potenzmengenadjunktion]{%
  X\in\powerset(\NatSeg{n})\vdash \tPowAdj{n}(X)\coloneqq X\cup\{n\}
}{%
  \DeltaRow{Mengen}{n\dsep X}
  \DeltaRow{Funktionensymbol}{\tPowAdj{n}}
}

\FormulaThmDelta[Typisierung der Funktionsvorschrift der Potenzmengenadjunktion]{%
  n\in\mathbb{N}\dsep X\in\powerset(\NatSeg{n})\vdash \tPowAdj{n}(X)\in\powerset(\NatSeg{\Succ(n)})
}{%
  \DeltaRow{Mengen}{n\dsep X}
}
\begin{tabproof}
  \proofstep{1}{n\in\mathbb{N}}{\rA}
  \proofstep{2}{X\in\powerset(\NatSeg{n})}{\rA}
  \proofstep{2}{X\subseteq \NatSeg{n}}{%
    \FormulaRefAuto{x \in \powerset(A)\;\eqvdash\; x \subseteq A}{2}}
  \proofstep{1}{\NatSeg{n}\subseteq \NatSeg{\Succ(n)}}{%
    \FormulaRefAuto{n\in\mathbb{N}\vdash \NatSeg{n}\subseteq \NatSeg{\Succ(n)}}{1}}
  \proofstep{1,2}{X\subseteq \NatSeg{\Succ(n)}}{%
    \FormulaRefAuto{A \subseteq B,\, B \subseteq C \vdash A \subseteq C}{3,4}}
  \proofstep{1}{n\in \NatSeg{\Succ(n)}}{%
    \FormulaRefAuto{n\in\mathbb{N}\vdash n\in \NatSeg{\Succ(n)}}{1}}
  \proofstep{1}{\{n\}\subseteq \NatSeg{\Succ(n)}}{%
    \FormulaRefAuto{a \in A \vdash \{a\} \subseteq A}{6}}
  \proofstep{1,2}{X\cup\{n\}\subseteq \NatSeg{\Succ(n)}}{%
    \FormulaRefAuto{A \subseteq C,\, B \subseteq C \vdash A \cup B \subseteq C}{5,7}}
  \proofstep{2}{\tPowAdj{n}(X)\coloneqq X\cup\{n\}}{%
    \FormulaRefAuto{X\in\powerset(\NatSeg{n})\vdash \tPowAdj{n}(X)\coloneqq X\cup\{n\}}{2}}
  \proofstep{1,2}{\tPowAdj{n}(X)\subseteq \NatSeg{\Succ(n)}}{\rIE{9,8}}
  \proofstep{1,2}{\tPowAdj{n}(X)\in\powerset(\NatSeg{\Succ(n)})}{%
    \FormulaRefAuto{x \in \powerset(A)\;\eqvdash\; x \subseteq A}{10}}
\end{tabproof}

\FormulaThmDeltaR[Existenz der Potenzmengenadjunktion]{%
  \begin{aligned}[t]
    &n\in\mathbb{N}\vdash{}\\
    &\exists! F\colon \powerset(\NatSeg{n})\to \powerset(\NatSeg{\Succ(n)})\,{}\\
    &\forall X\in\powerset(\NatSeg{n})\,F(X)=\tPowAdj{n}(X)
  \end{aligned}
}{%
  n\in\mathbb{N}\vdash
  \exists! F\colon \powerset(\NatSeg{n})\to \powerset(\NatSeg{\Succ(n)})\,
  \forall X\in\powerset(\NatSeg{n})\,F(X)=\tPowAdj{n}(X)
}{%
  \DeltaRow{Mengen}{n\dsep X}
  \DeltaRow{Funktionen}{F}
}
\begin{tabproof}
  \proofstep{1}{n\in\mathbb{N}}{\rA}
  \proofstep{1}{%
    \begin{aligned}[t]
      &\exists! F\colon \powerset(\NatSeg{n})\to \powerset(\NatSeg{\Succ(n)})\,{}\\
      &\forall X\in\powerset(\NatSeg{n})\,F(X)=\tPowAdj{n}(X)
    \end{aligned}
  }{%
    \FormulaRefAuto{%
      x\in A\vdash t(x)\coloneq y \dsep x\in A\vdash t(x)\in B
      \exists! F\colon A\to B\,\forall x\in A\,F(x)=t(x)
    }{%
      \FormulaRefAuto{%
        X\in\powerset(\NatSeg{n})\vdash \tPowAdj{n}(X)\coloneqq X\cup\{n\}
      },%
      \FormulaRefAuto{%
        n\in\mathbb{N}\dsep X\in\powerset(\NatSeg{n})\vdash
        \tPowAdj{n}(X)\in\powerset(\NatSeg{\Succ(n)})
      }{1}%
    }%
  }
\end{tabproof}

\FormulaDefDelta[Potenzmengenadjunktion]{%
  n\in\mathbb{N}\vdash
  \PowAdj{n}\coloneqq \iota F\Bigl(
    F\colon \powerset(\NatSeg{n})\to \powerset(\NatSeg{\Succ(n)})
    \land
    \forall X\in\powerset(\NatSeg{n})\,F(X)=\tPowAdj{n}(X)
  \Bigr)
}{%
  \DeltaRow{Mengen}{n\dsep X}
  \DeltaRow{Funktionen}{\PowAdj{n}}
}

\FormulaThmDelta[Typisierung der Potenzmengenadjunktion]{%
  n\in\mathbb{N}\vdash \PowAdj{n}\colon \powerset(\NatSeg{n})\to \powerset(\NatSeg{\Succ(n)})
}{%
  \DeltaRow{Mengen}{n\dsep X}
  \DeltaRow{Funktionen}{\PowAdj{n}}
}
\begin{tabproof}
  \proofstep{1}{n\in\mathbb{N}}{\rA}
  \proofstep{1}{\PowAdj{n}\colon \powerset(\NatSeg{n})\to \powerset(\NatSeg{\Succ(n)})}{%
    \rAEa{\FormulaRefAuto{%
      n\in\mathbb{N}\vdash
      \PowAdj{n}\coloneqq \iota F\Bigl(
        F\colon \powerset(\NatSeg{n})\to \powerset(\NatSeg{\Succ(n)})
        \land
        \forall X\in\powerset(\NatSeg{n})\,F(X)=\tPowAdj{n}(X)
      \Bigr)
    }{1}}}
\end{tabproof}

\FormulaThmDelta[Werte der Potenzmengenadjunktion]{%
  n\in\mathbb{N}\dsep X\in\powerset(\NatSeg{n})\vdash \PowAdj{n}(X)=\tPowAdj{n}(X)
}{%
  \DeltaRow{Mengen}{n\dsep X}
  \DeltaRow{Funktionen}{\PowAdj{n}}
}
\begin{tabproof}
  \proofstep{1}{n\in\mathbb{N}}{\rA}
  \proofstep{2}{X\in\powerset(\NatSeg{n})}{\rA}
  \proofstep{1}{\forall X\in\powerset(\NatSeg{n})\,\PowAdj{n}(X)=\tPowAdj{n}(X)}{%
    \rAEb{\FormulaRefAuto{%
      n\in\mathbb{N}\vdash
      \PowAdj{n}\coloneqq \iota F\Bigl(
        F\colon \powerset(\NatSeg{n})\to \powerset(\NatSeg{\Succ(n)})
        \land
        \forall X\in\powerset(\NatSeg{n})\,F(X)=\tPowAdj{n}(X)
      \Bigr)
    }{1}}}
  \proofstep{1,2}{\PowAdj{n}(X)=\tPowAdj{n}(X)}{\rUE{3}}
\end{tabproof}

\FormulaThmDelta[Explizite Werte der Potenzmengenadjunktion]{%
  n\in\mathbb{N}\dsep X\in\powerset(\NatSeg{n})\vdash \PowAdj{n}(X)=X\cup\{n\}
}{%
  \DeltaRow{Mengen}{n\dsep X}
  \DeltaRow{Funktionen}{\PowAdj{n}}
}
\begin{tabproof}
  \proofstep{1}{n\in\mathbb{N}}{\rA}
  \proofstep{2}{X\in\powerset(\NatSeg{n})}{\rA}
  \proofstep{1,2}{\PowAdj{n}(X)=\tPowAdj{n}(X)}{%
    \FormulaRefAuto{n\in\mathbb{N}\dsep X\in\powerset(\NatSeg{n})\vdash \PowAdj{n}(X)=\tPowAdj{n}(X)}{1,2}}
  \proofstep{2}{\tPowAdj{n}(X)\coloneqq X\cup\{n\}}{%
    \FormulaRefAuto{X\in\powerset(\NatSeg{n})\vdash \tPowAdj{n}(X)\coloneqq X\cup\{n\}}{2}}
  \proofstep{1,2}{\PowAdj{n}(X)=X\cup\{n\}}{\rIE{3,4}}
\end{tabproof}

\FormulaThmDelta[Fall \(n\notin Y\) der Zerlegung der Potenzmenge]{%
  n\in\mathbb{N}\dsep Y\in\powerset(\NatSeg{\Succ(n)})\dsep n\notin Y \vdash Y\in\powerset(\NatSeg{n})
}{%
  \DeltaRow{Mengen}{n\dsep Y}
}
\begin{tabproof}
  \proofstep{1}{n\in\mathbb{N}}{\rA}
  \proofstep{2}{Y\in\powerset(\NatSeg{\Succ(n)})}{\rA}
  \proofstep{3}{n\notin Y}{\rA}
  \proofstep{2}{Y\subseteq \NatSeg{\Succ(n)}}{%
    \FormulaRefAuto{x \in \powerset(A)\;\eqvdash\; x \subseteq A}{2}}
  \proofstep{1,2,3}{Y\subseteq \NatSeg{n}}{%
    \FormulaRefAuto{n\in\mathbb{N}\dsep B\subseteq \NatSeg{\Succ(n)}\dsep n\notin B \vdash B\subseteq \NatSeg{n}}{1,4,3}}
  \proofstep{1,2,3}{Y\in\powerset(\NatSeg{n})}{%
    \FormulaRefAuto{x \in \powerset(A)\;\eqvdash\; x \subseteq A}{5}}
\end{tabproof}

\FormulaThmDelta[Fall \(n\in Y\) der Zerlegung der Potenzmenge]{%
  n\in\mathbb{N}\dsep Y\in\powerset(\NatSeg{\Succ(n)})\dsep n\in Y \vdash Y\in \PowAdj{n}[\powerset(\NatSeg{n})]
}{%
  \DeltaRow{Mengen}{n\dsep Y\dsep X}
  \DeltaRow{Funktionen}{\PowAdj{n}}
}
\begin{tabproof}
  \proofstep{1}{n\in\mathbb{N}}{\rA}
  \proofstep{2}{Y\in\powerset(\NatSeg{\Succ(n)})}{\rA}
  \proofstep{3}{n\in Y}{\rA}
  \proofstep{2}{Y\subseteq \NatSeg{\Succ(n)}}{%
    \FormulaRefAuto{x \in \powerset(A)\eqvdash x \subseteq A}{2}}
  \proofstep{1,2,3}{Y=(Y\cap \NatSeg{n})\cup\{n\}}{%
    \FormulaRefAuto{PeanoNatSegSubsetSuccWithEndpoint}{1,4,3}}
  \proofstep{}{Y\cap \NatSeg{n}\subseteq \NatSeg{n}}{\FormulaRefAuto{A\cap B \subseteq B}}
  \proofstep{}{Y\cap \NatSeg{n}\in\powerset(\NatSeg{n})}{%
    \FormulaRefAuto{x \in \powerset(A)\eqvdash x \subseteq A}{6}}
  \proofstep{1}{\PowAdj{n}\colon \powerset(\NatSeg{n})\to \powerset(\NatSeg{\Succ(n)})}{%
    \FormulaRefAuto{n\in\mathbb{N}\vdash \PowAdj{n}\colon \powerset(\NatSeg{n})\to \powerset(\NatSeg{\Succ(n)})}{1}}
  \proofstep{1}{\PowAdj{n}(Y\cap \NatSeg{n})=(Y\cap \NatSeg{n})\cup\{n\}}{%
    \FormulaRefAuto{n\in\mathbb{N}\dsep X\in\powerset(\NatSeg{n})\vdash \PowAdj{n}(X)=X\cup\{n\}}{1,7}}
  \proofstep{1}{(Y\cap \NatSeg{n})\cup\{n\}=\PowAdj{n}(Y\cap \NatSeg{n})}{%
    \FormulaRefAuto{a=b\vdash b=a}{9}}
  \proofstep{1,2,3}{Y=\PowAdj{n}(Y\cap \NatSeg{n})}{\rIE{5,10}}
  \proofstep{1,2,3}{\exists X\in\powerset(\NatSeg{n})\,Y=\PowAdj{n}(X)}{\rEI{\rAI{7,11}}}
  \proofstep{1,2,3}{Y\in \PowAdj{n}[\powerset(\NatSeg{n})]}{%
    \FormulaRefAuto{F\colon A\to B\dsep \exists x\in A\,y=F(x)\vdash y\in F[A]}{8,12}}
\end{tabproof}

\FormulaThmDeltaR[Zerlegung der Potenzmenge einer Nachfolgermenge]{%
  n\in\mathbb{N}\vdash \powerset(\NatSeg{\Succ(n)})=\powerset(\NatSeg{n})\cup \PowAdj{n}[\powerset(\NatSeg{n})]
}{%
  n\in\mathbb{N}\vdash \powerset(\NatSeg{\Succ(n)})=\powerset(\NatSeg{n})\cup \PowAdj{n}[\powerset(\NatSeg{n})]
}{%
  \DeltaRow{Mengen}{n\dsep Y\dsep X}
  \DeltaRow{Funktionen}{\PowAdj{n}}
}
\begin{tabproofsplitwide}
  \proofpartwideindR[Teilmengenrichtung \(\powerset(\NatSeg{\Succ(n)})\subseteq \powerset(\NatSeg{n})\cup \PowAdj{n}[\powerset(\NatSeg{n})]\)]{%
    n\in\mathbb{N}\dsep Y\in\powerset(\NatSeg{\Succ(n)})\vdash Y\in\powerset(\NatSeg{n})\cup \PowAdj{n}[\powerset(\NatSeg{n})]
  }{%
    n\in\mathbb{N}\dsep Y\in\powerset(\NatSeg{\Succ(n)})\vdash Y\in\powerset(\NatSeg{n})\cup \PowAdj{n}[\powerset(\NatSeg{n})]
  }
    \proofstepwidestar[1]{n\in\mathbb{N}}{\rA}
    \proofstepwidestar[2]{Y\in\powerset(\NatSeg{\Succ(n)})}{\rA}
    \proofstepwidestar[]{n\in Y\lor n\notin Y}{\FormulaRefAuto{P \lor \neg P}}

    \proofstepwidestar[4]{n\in Y}{\rA}
    \proofstepwidestar[1,2,4]{Y\in \PowAdj{n}[\powerset(\NatSeg{n})]}{%
      \FormulaRefAuto{n\in\mathbb{N}\dsep Y\in\powerset(\NatSeg{\Succ(n)})\dsep n\in Y \vdash Y\in \PowAdj{n}[\powerset(\NatSeg{n})]}{1,2,4}}
    \proofstepwidestar[1,2,4]{Y\in\powerset(\NatSeg{n})\cup \PowAdj{n}[\powerset(\NatSeg{n})]}{%
      \FormulaRefAuto{z \in B \vdash z \in A \cup B}{5}}

    \proofstepwidestar[7]{n\notin Y}{\rA}
    \proofstepwidestar[1,2,7]{Y\in\powerset(\NatSeg{n})}{%
      \FormulaRefAuto{n\in\mathbb{N}\dsep Y\in\powerset(\NatSeg{\Succ(n)})\dsep n\notin Y \vdash Y\in\powerset(\NatSeg{n})}{1,2,7}}
    \proofstepwidestar[1,2,7]{Y\in\powerset(\NatSeg{n})\cup \PowAdj{n}[\powerset(\NatSeg{n})]}{%
      \FormulaRefAuto{z \in A \vdash z \in A \cup B}{8}}

    \proofstepwidestar[1,2]{Y\in\powerset(\NatSeg{n})\cup \PowAdj{n}[\powerset(\NatSeg{n})]}{\rOE{3,4,6,7,9}}
  \closeproofpartwideindR

  \proofpartwideindR[Teilmengenrichtung \(\powerset(\NatSeg{n})\cup \PowAdj{n}[\powerset(\NatSeg{n})]\subseteq \powerset(\NatSeg{\Succ(n)})\)]{%
    n\in\mathbb{N}\dsep Y\in\powerset(\NatSeg{n})\cup \PowAdj{n}[\powerset(\NatSeg{n})]\vdash Y\in\powerset(\NatSeg{\Succ(n)})
  }{%
    n\in\mathbb{N}\dsep Y\in\powerset(\NatSeg{n})\cup \PowAdj{n}[\powerset(\NatSeg{n})]\vdash Y\in\powerset(\NatSeg{\Succ(n)})
  }
    \proofstepwidestar[1]{n\in\mathbb{N}}{\rA}
    \proofstepwidestar[2]{Y\in\powerset(\NatSeg{n})\cup \PowAdj{n}[\powerset(\NatSeg{n})]}{\rA}
    \proofstepwidestar[2]{Y\in\powerset(\NatSeg{n})\lor Y\in \PowAdj{n}[\powerset(\NatSeg{n})]}{%
      \FormulaRefAuto{x\in A\cup B \eqvdash x\in A\lor x\in B}{2}}

    \proofstepwidestar[4]{Y\in\powerset(\NatSeg{n})}{\rA}
    \proofstepwidestar[1]{\NatSeg{n}\subseteq \NatSeg{\Succ(n)}}{%
      \FormulaRefAuto{n\in\mathbb{N}\vdash \NatSeg{n}\subseteq \NatSeg{\Succ(n)}}{1}}
    \proofstepwidestar[1,4]{\powerset(\NatSeg{n})\subseteq \powerset(\NatSeg{\Succ(n)})}{%
      \FormulaRefAuto{A \subseteq B \vdash \powerset(A) \subseteq \powerset(B)}{5}}
    \proofstepwidestar[1,4]{Y\in\powerset(\NatSeg{\Succ(n)})}{\rRE{\rUE{6},4}}

    \proofstepwidestar[8]{Y\in \PowAdj{n}[\powerset(\NatSeg{n})]}{\rA}
    \proofstepwidestar[1]{\PowAdj{n}\colon \powerset(\NatSeg{n})\to \powerset(\NatSeg{\Succ(n)})}{%
      \FormulaRefAuto{n\in\mathbb{N}\vdash \PowAdj{n}\colon \powerset(\NatSeg{n})\to \powerset(\NatSeg{\Succ(n)})}{1}}
    \proofstepwidestar[1,8]{\exists X\in\powerset(\NatSeg{n})\,Y=\PowAdj{n}(X)}{%
      \FormulaRefAuto{F\colon A\to B\dsep y\in F[A]\vdash \exists x\in A\,y=F(x)}{9,8}}

    \proofstepwidestar[11]{X\in\powerset(\NatSeg{n})\land Y=\PowAdj{n}(X)}{\rA}
    \proofstepwidestar[11]{X\in\powerset(\NatSeg{n})}{\rAEa{11}}
    \proofstepwidestar[11]{Y=\PowAdj{n}(X)}{\rAEb{11}}
    \proofstepwidestar[1,11]{\PowAdj{n}(X)\in\powerset(\NatSeg{\Succ(n)})}{%
      \FormulaRefAuto{F\colon A\to B\dsep x\in A\vdash F(x)\in B}{9,12}}
    \proofstepwidestar[1,11]{Y\in\powerset(\NatSeg{\Succ(n)})}{\rIE{13,14}}

    \proofstepwidestar[1,8]{Y\in\powerset(\NatSeg{\Succ(n)})}{\rEE{10,11,15}}
    \proofstepwidestar[1,2]{Y\in\powerset(\NatSeg{\Succ(n)})}{\rOE{3,4,7,8,16}}
  \closeproofpartwideindR

  \proofpartwide{Schluss}
    \proofstepwidestar[1]{n\in\mathbb{N}}{\rA}
    \proofstepwidestar[1]{\powerset(\NatSeg{\Succ(n)})\subseteq \powerset(\NatSeg{n})\cup \PowAdj{n}[\powerset(\NatSeg{n})]}{%
      \FormulaRefAuto{n\in\mathbb{N}\dsep Y\in\powerset(\NatSeg{\Succ(n)})\vdash Y\in\powerset(\NatSeg{n})\cup \PowAdj{n}[\powerset(\NatSeg{n})]}{1}}
    \proofstepwidestar[1]{\powerset(\NatSeg{n})\cup \PowAdj{n}[\powerset(\NatSeg{n})]\subseteq \powerset(\NatSeg{\Succ(n)})}{%
      \FormulaRefAuto{n\in\mathbb{N}\dsep Y\in\powerset(\NatSeg{n})\cup \PowAdj{n}[\powerset(\NatSeg{n})]\vdash Y\in\powerset(\NatSeg{\Succ(n)})}{1}}
    \proofstepwidestar[1]{\powerset(\NatSeg{\Succ(n)})=\powerset(\NatSeg{n})\cup \PowAdj{n}[\powerset(\NatSeg{n})]}{%
      \FormulaRefAuto{A \subseteq B, B \subseteq A \vdash A = B}{2,3}}
  \closeproofpartwide
\end{tabproofsplitwide}

\FormulaThmDeltaK[Potenzmengen natürlicher Zahlen sind endlich]{%
  n\in\mathbb{N}\vdash \Finite{\powerset(\NatSeg{n})}
}{FiniteNatSegPowerSet}{%
  \DeltaRow{Mengen}{n}
}
\begin{tabproofsplitwide}
  \proofpartwideind[Induktionsanfang]{%
    \Finite{\powerset(\NatSeg{0})}}
    \proofstepwidestar[]{\NatSeg{0}=\varnothing}{%
      \FormulaRefAuto{PeanoNatSegZero}}
    \proofstepwidestar[]{\Finite{\varnothing}}{\FormulaRefAuto{FiniteEmpty}[thm]}
    \proofstepwidestar[]{\varnothing\notin\varnothing}{\FormulaRefAuto{x \not\in \varnothing}}
    \proofstepwidestar[]{\Finite{\varnothing\cup\{\varnothing\}}}{%
      \FormulaRefAuto{FiniteAdjunction}[thm]{3,2}}
    \proofstepwidestar[]{\powerset(\varnothing)=\{\varnothing\}}{%
      \FormulaRefAuto{\powerset(\varnothing)=\{\varnothing\}}}
    \proofstepwidestar[]{\{\varnothing\}=\varnothing\cup\{\varnothing\}}{%
      \FormulaRefAuto{A = \varnothing \cup A}}
    \proofstepwidestar[]{\powerset(\varnothing)=\varnothing\cup\{\varnothing\}}{\rIE{5,6}}
    \proofstepwidestar[]{\Finite{\powerset(\varnothing)}}{\rIE{7,4}}
    \proofstepwidestar[]{\Finite{\powerset(\NatSeg{0})}}{\rIE{1,8}}
  \closeproofpartwideind

  \proofpartwideindR[Induktionsschritt]{%
    n\in\mathbb{N}\dsep \Finite{\powerset(\NatSeg{n})} \vdash \Finite{\powerset(\NatSeg{\Succ(n)})}
  }{%
    n\in\mathbb{N}\dsep \Finite{\powerset(\NatSeg{n})} \vdash \Finite{\powerset(\NatSeg{\Succ(n)})}
  }
    \proofstepwidestar[1]{n\in\mathbb{N}}{\rA}
    \proofstepwidestar[2]{\Finite{\powerset(\NatSeg{n})}}{\rA}
    \proofstepwidestar[1]{\powerset(\NatSeg{\Succ(n)})=\powerset(\NatSeg{n})\cup \PowAdj{n}[\powerset(\NatSeg{n})]}{%
      \FormulaRefAuto{n\in\mathbb{N}\vdash \powerset(\NatSeg{\Succ(n)})=\powerset(\NatSeg{n})\cup \PowAdj{n}[\powerset(\NatSeg{n})]}{1}}
    \proofstepwidestar[]{\powerset(\NatSeg{n})\subseteq \powerset(\NatSeg{n})}{\FormulaRefAuto{A \subseteq A}}
    \proofstepwidestar[1]{\PowAdj{n}\colon \powerset(\NatSeg{n})\to \powerset(\NatSeg{\Succ(n)})}{%
      \FormulaRefAuto{n\in\mathbb{N}\vdash \PowAdj{n}\colon \powerset(\NatSeg{n})\to \powerset(\NatSeg{\Succ(n)})}{1}}
    \proofstepwidestar[1,2]{\Finite{\PowAdj{n}[\powerset(\NatSeg{n})]}}{%
      \FormulaRefAuto{FiniteImage}[thm]{4,5,2}}
    \proofstepwidestar[1,2]{\Finite{\powerset(\NatSeg{n})\cup \PowAdj{n}[\powerset(\NatSeg{n})]}}{%
      \FormulaRefAuto{FiniteUnion}[thm]{2,6}}
    \proofstepwidestar[1,2]{\Finite{\powerset(\NatSeg{\Succ(n)})}}{\rIE{3,7}}
  \closeproofpartwideindR

  \proofpartwide{Induktionsschluss}
    \proofstepwidestar[1]{n\in\mathbb{N}}{\rA}
    \proofstepwidestar[1]{\Finite{\powerset(\NatSeg{n})}}{%
      \FormulaRefAuto{PeanoInductionRule}{IA,IS,1}}
  \closeproofpartwide
\end{tabproofsplitwide}

\subsection{Potenzmengen endlicher Mengen sind endlich}

\FormulaThmDeltaK[Potenzmengen endlicher Mengen sind endlich]{%
  \Finite{A}\vdash \Finite{\powerset(A)}
}{FinitePowerSet}{%
  \DeltaRow{Mengen}{A\dsep n}
}
\begin{tabproof}
  \proofstep{1}{\Finite{A}}{\rA}
  \proofstep{1}{\exists n\,\bigl(n\in\mathbb{N}\land \NatSeg{n}\EqCard A\bigr)}{%
    \FormulaRefAuto{FiniteEqCardEquiv}[thm]{1}}

  \proofstep{3}{n\in\mathbb{N}\land \NatSeg{n}\EqCard A}{\rA}
  \proofstep{3}{n\in\mathbb{N}}{\rAEa{3}}
  \proofstep{3}{\NatSeg{n}\EqCard A}{\rAEb{3}}
  \proofstep{4}{\Finite{\powerset(\NatSeg{n})}}{%
    \FormulaRefAuto{FiniteNatSegPowerSet}[thm]{4}}
  \proofstep{5}{\powerset(\NatSeg{n})\EqCard \powerset(A)}{%
    \FormulaRefAuto{A \EqCard B \vdash \powerset(A)\EqCard \powerset(B)}{5}}
  \proofstep{3}{\Finite{\powerset(A)}}{\FormulaRefAuto{FiniteEqCardTransport}[thm]{6,7}}
  \proofstep{1}{\Finite{\powerset(A)}}{\rEE{2,3,8}}
\end{tabproof}

\section{Äquivalente Formulierungen der Endlichkeit}

\paragraph{Implikationsdiagramm}

Wir verwenden in diesem Abschnitt die folgenden Abkürzungen:
\[
\begin{aligned}
D(M)&\coloneqq \neg\exists H\colon \mathbb{N}\inj M,\\
S(M)&\coloneqq \neg\exists H\colon M\sur\mathbb{N},\\
I(M)&\coloneqq \forall F\colon M\inj M\,\bigl(F\colon M\bij M\bigr),\\
Q(M)&\coloneqq \forall F\colon M\sur M\,\bigl(F\colon M\bij M\bigr),\\
T_{\max}(M)&\coloneqq \TarskiFinite(M),\\
T_{\min}(M)&\coloneqq \TarskiMinFinite(M).
\end{aligned}
\]
Die Markierung \(\mathrm{AC}\) bedeutet, dass die angegebene Folgerung in
der hier verwendeten Beweisroute ein zuvor mit dem Auswahlaxiom bewiesenes
Resultat benutzt.
\[
\begin{array}{ccccc}
\Finite{M} & \xrightarrow{\mathrm{ZF}} & D(M)
  & \xrightarrow{\mathrm{AC}} & S(M)\\[1.1ex]
&& D(M) & \xleftarrow{\mathrm{ZF}} & S(M)\\[1.1ex]
&& D(M) & \xrightarrow{\mathrm{AC}} &
  T_{\max}(M)\overset{\mathrm{ZF}}{\Longleftrightarrow}
  T_{\min}(M)\xrightarrow{\mathrm{ZF}}\Finite{M}\\[1.1ex]
S(M) & \xrightarrow{\mathrm{ZF}} & I(M)
  & \xrightarrow{\mathrm{AC}} & Q(M)\\[1.1ex]
&& I(M) & \xleftarrow{\mathrm{ZF}} & Q(M).
\end{array}
\]

\subsection[Ausschluss von Injektionen N -> M]{%
  Ausschluss von Injektionen \(\mathbb{N}\to M\)}

\FormulaThmDeltaK[Transport des Injektionsausschlusses von einer natürlichen Zahl]{%
    n\in\mathbb{N}\dsep \NatSeg{n}\EqCard M
    \vdash \neg\exists F\colon \mathbb{N}\inj M %
}{FiniteWitnessNoNatInjection}{%
  \DeltaDecl{Mengen}{n\dsep M}%
  \DeltaDecl{Funktionen}{F\dsep G\dsep H}%
}
\begin{tabproof}
  \proofstep{1}{n\in\mathbb{N}}{\rA}
  \proofstep{2}{\NatSeg{n}\EqCard M}{\rA}

  \proofstep{3}{\exists F\colon \mathbb{N}\inj M}{\rA}

  \proofstep{2}{\exists G\colon \NatSeg{n}\bij M}{\FormulaRefAuto{A \EqCard B \coloneqq \exists F\colon A\bij B}{2}}

  \proofstep{5}{G\colon \NatSeg{n}\bij M}{\rA}

  \proofstep{5}{G^{-1}\colon M\bij \NatSeg{n}}{\FormulaRefAuto{F\colon A\bij B \vdash F^{-1}\colon B\bij A}{5}}
  \proofstep{5}{G^{-1}\colon M\inj \NatSeg{n}}{\FormulaRefAuto{F\colon A\bij B \vdash F\colon A\inj B}{6}}

  \proofstep{8}{F\colon \mathbb{N}\inj M}{\rA}

  \proofstep{5,8}{G^{-1}\circ F\colon \mathbb{N}\inj \NatSeg{n}}{\FormulaRefAuto{F\colon A\inj B \dsep G\colon B\inj C\vdash G\circ F\colon A\inj C}{8,7}}

  \proofstep{5,8}{\exists H\colon \mathbb{N}\inj \NatSeg{n}}{\rEI{9}}
  \proofstep{1}{\neg\exists H\colon \mathbb{N}\inj \NatSeg{n}}{%
    \FormulaRefAuto{PeanoNatSegNoNatInjection}[thm]{1}}
  \proofstep{1,5,8}{\bot}{\rBI{10,11}}

  \proofstep{1,3,5}{\bot}{\rEE{3,8,12}}
  \proofstep{1,2,3}{\bot}{\rEE{4,5,13}}

  \proofstep{1,2}{\neg\exists F\colon \mathbb{N}\inj M}{\rCI{3,14}}
\end{tabproof}

\subsection[Injektionen N -> M und Surjektionen M -> N]{%
  Injektionen \(\mathbb{N}\to M\) und Surjektionen \(M\to \mathbb{N}\)}
\FormulaThmDelta{%
  \exists G\colon \mathbb{N}\inj M\vdash
  \exists H\colon M\sur\mathbb{N}
}{%
  \DeltaDecl{Mengen}{M}%
  \DeltaDecl{Funktionen}{G\dsep H}%
}
\begin{tabproof}
  \proofstep{1}{\exists G\colon \mathbb{N}\inj M}{\rA}

  \proofstep{2}{G\colon \mathbb{N}\inj M}{\rA}
  \proofstep{}{\varnothing\in\mathbb{N}}{\FormulaRefAuto{\varnothing \in \mathbb{N}}}

  \proofstep{2}{\Gsurjfrominj{G}{\varnothing}\colon M\sur\mathbb{N}}{%
    \FormulaRefAuto{F\colon A\inj B \dsep a_0\in A \vdash \Gsurjfrominj{F}{a_0}\colon B\sur A}{2,3}}

  \proofstep{2}{\exists H\colon M\sur\mathbb{N}}{\rEI{4}}

  \proofstep{1}{\exists H\colon M\sur\mathbb{N}}{\rEE{1,2,5}}
\end{tabproof}

\FormulaThmDelta{%
  \neg\exists F\colon \mathbb{N}\inj M
  \vdash
  \neg\exists F\colon M\sur \mathbb{N}%
}{%
  \DeltaDecl{Mengen}{M}%
  \DeltaDecl{Funktionen}{F\dsep H}%
}
\begin{tabproof}
  \proofstep{1}{\neg\exists F\colon \mathbb{N}\inj M}{\rA}

  \proofstep{2}{\exists F\colon M\sur \mathbb{N}}{\rA}
  \proofstep{3}{F\colon M\sur \mathbb{N}}{\rA}

  \proofstep{3}{\exists H\colon \mathbb{N}\inj M}{%
    \FormulaRefAuto{F\colon A\sur B\vdash \exists H\colon B\inj A}{3}}

  \proofstep{1,3}{\bot}{\rBI{1,4}}
  \proofstep{1,2}{\bot}{\rEE{2,3,5}}
  \proofstep{1}{\neg\exists F\colon M\sur \mathbb{N}}{\rCI{2,6}}
\end{tabproof}

\FormulaThmDeltaR{%
  \exists F\colon \mathbb{N}\inj M
  \eqvdash
  \exists F\colon M\sur\mathbb{N}
}{%
  \exists F\colon \mathbb{N}\inj M
  \eqvdash
  \exists F\colon M\sur\mathbb{N}
}{%
  \DeltaDecl{Mengen}{M}%
  \DeltaDecl{Funktionen}{F\dsep G\dsep H}%
}
\begin{tabproofsplit}
  \proofpart{\(\vdash\)}
    \proofstep{1}{\exists F\colon \mathbb{N}\inj M}{\rA}
    \proofstep{1}{\exists F\colon M\sur\mathbb{N}}{%
      \FormulaRefAuto{%
        \exists G\colon \mathbb{N}\inj M\vdash
        \exists H\colon M\sur\mathbb{N}
      }{1}}
  \closeproofpart

  \proofpart{\(\dashv\)}
    \proofstep{1}{\exists F\colon M\sur\mathbb{N}}{\rA}

    \proofstep{2}{\neg\exists F\colon \mathbb{N}\inj M}{\rA}
    \proofstep{2}{\neg\exists F\colon M\sur\mathbb{N}}{%
      \FormulaRefAuto{%
        \neg\exists F\colon \mathbb{N}\inj M
        \vdash
        \neg\exists F\colon M\sur \mathbb{N}
      }{2}}

    \proofstep{1,2}{\bot}{\rBI{1,3}}
    \proofstep{1}{\neg\neg\exists F\colon \mathbb{N}\inj M}{\rCI{2,4}}
    \proofstep{1}{\exists F\colon \mathbb{N}\inj M}{\FormulaRefAuto{rDN}{5}}
  \closeproofpart
\end{tabproofsplit}

\subsection{Dedekindsche Charakterisierung endlicher Mengen}

\subsubsection{Orbit einer injektiven Selbstabbildung}

\FormulaThmDeltaR[Kein echter Folgewert ist der Startwert]{%
  \begin{aligned}[t]
    &m_0\in M\dsep f\colon M\inj M\dsep{}\\
    &\forall x\in M\,\bigl(f(x)\neq m_0\bigr)\vdash{}\\
    &\forall n\in\mathbb{N}\,\bigl(
      \RecFun{m_0}{f}(\Succ(n))\neq m_0
    \bigr)
  \end{aligned}
}{%
  m_0\in M\dsep f\colon M\inj M\dsep
  \forall x\in M\,\bigl(f(x)\neq m_0\bigr)\vdash
  \forall n\in\mathbb{N}\,\bigl(
    \RecFun{m_0}{f}(\Succ(n))\neq m_0
  \bigr)
}{%
  \DeltaDecl{Mengen}{M\dsep n\dsep x\dsep m_0}%
  \DeltaDecl{Funktionen}{f}%
}
\begin{tabproof}
  \proofstep{1}{m_0\in M}{\rA}
  \proofstep{2}{f\colon M\inj M}{\rA}
  \proofstep{3}{\forall x\in M\,\bigl(f(x)\neq m_0\bigr)}{\rA}

  \proofstep{2}{f\colon M\to M}{\FormulaRefAuto{Injektive Funktion}{2}}

  \proofstep{5}{n\in\mathbb{N}}{\rA}

  \proofstep{1,2}{\RecFun{m_0}{f}\colon \mathbb{N}\to M}{%
    \FormulaRefAuto{m_0\in M\dsep f\colon M\to M \vdash \RecFun{m_0}{f}\colon \mathbb{N}\to M}{1,4}}

  \proofstep{1,2,5}{\RecFun{m_0}{f}(n)\in M}{%
    \FormulaRefAuto{F\colon A\to B\dsep x\in A\vdash F(x)\in B}{6,5}}

  \proofstep{1,2,5}{%
    \RecFun{m_0}{f}(\Succ(n))=f(\RecFun{m_0}{f}(n))
  }{%
    \FormulaRefAuto{m_0\in M\dsep f\colon M\to M\dsep n\in\mathbb{N}\vdash \RecFun{m_0}{f}(\Succ(n))=f(\RecFun{m_0}{f}(n))}{1,4,5}}

  \proofstep{9}{\RecFun{m_0}{f}(\Succ(n))=m_0}{\rA}

  \proofstep{1,2,5}{%
    f(\RecFun{m_0}{f}(n))=\RecFun{m_0}{f}(\Succ(n))
  }{%
    \FormulaRefAuto{a=b\vdash b=a}{8}}

  \proofstep{1,2,5,9}{f(\RecFun{m_0}{f}(n))=m_0}{%
    \FormulaRefAuto{a=b,\, b=c \vdash a=c}{10,9}}

  \proofstep{1,2,3,5}{f(\RecFun{m_0}{f}(n))\neq m_0}{%
    \rRE{\rUE{3},7}}

  \proofstep{1,2,3,5,9}{\bot}{\rBI{11,12}}

  \proofstep{1,2,3,5}{\RecFun{m_0}{f}(\Succ(n))\neq m_0}{\rCI{9,13}}

  \proofstep{1,2,3}{%
    \forall n\in\mathbb{N}\,\bigl(
      \RecFun{m_0}{f}(\Succ(n))\neq m_0
    \bigr)
  }{%
    \rUI{\rRI{5,14}}}
\end{tabproof}

\subsubsection{Dedekindsche Charakterisierung}

\FormulaThmDelta{%
  \neg\exists H\colon M\sur\mathbb{N}\vdash
  \forall F\colon M\inj M\,\bigl(F\colon M\bij M\bigr)
}{%
  \DeltaDecl{Mengen}{M\dsep m_0\dsep x}%
  \DeltaDecl{Funktionen}{F\dsep G\dsep H}%
}
\begin{tabproof}
  \proofstep{1}{\neg\exists H\colon M\sur\mathbb{N}}{\rA}

  \proofstep{2}{F\colon M\inj M}{\rA}
  \proofstep{2}{F\colon M\to M}{%
    \FormulaRefAuto{Injektive Funktion}{2}}

  \proofstep{4}{\neg(F\colon M\sur M)}{\rA}

  \proofstep{2,4}{%
    \exists m_0\in M\,\forall x\in M\,\bigl(F(x)\neq m_0\bigr)
  }{%
    \FormulaRefAuto{%
      F\colon M\to M\dsep \neg(F\colon M\sur M)\vdash
      \exists m_0\in M\,\forall x\in M\,\bigl(F(x)\neq m_0\bigr)
    }{3,4}}

  \proofstep{6}{m_0\in M}{\rA}
  \proofstep{7}{%
    \forall x\in M\,\bigl(F(x)\neq m_0\bigr)
  }{\rA}

  \proofstep{2,6,7}{%
    \RecFun{m_0}{F}\colon \mathbb{N}\inj M
  }{%
    \FormulaRefAuto{%
      m_0\in M\dsep f\colon M\inj M\dsep
      \forall x\in M\,\bigl(f(x)\neq m_0\bigr)\vdash
      \RecFun{m_0}{f}\colon \mathbb{N}\inj M
    }{6,2,7}}

  \proofstep{2,6,7}{\exists G\colon \mathbb{N}\inj M}{\rEI{8}}

  \proofstep{2,6,7}{\exists H\colon M\sur\mathbb{N}}{%
    \FormulaRefAuto{%
      \exists G\colon \mathbb{N}\inj M\vdash
      \exists H\colon M\sur\mathbb{N}
    }{9}}

  \proofstep{1,2,6,7}{\bot}{\rBI{1,10}}

  \proofstep{1,2,4}{\bot}{\rEE{5,6,7,11}}

  \proofstep{1,2}{F\colon M\sur M}{\rCI{4,12}}

  \proofstep{1,2}{F\colon M\bij M}{%
    \FormulaRefAuto{F\colon A\inj B\dsep F\colon A\sur B \vdash F\colon A\bij B}{2,13}}

  \proofstep{1}{%
    \forall F\colon M\inj M\,\bigl(F\colon M\bij M\bigr)
  }{\rUI{\rRI{2,14}}}
\end{tabproof}

\FormulaThmDelta{%
  \forall F\colon M\inj M\,\bigl(F\colon M\bij M\bigr)\vdash
  \forall F\colon M\sur M\,\bigl(F\colon M\bij M\bigr)
}{%
  \DeltaDecl{Mengen}{M\dsep x\dsep y}%
  \DeltaDecl{Funktionen}{F\dsep H\dsep K}%
}
\begin{tabproofsplitwide}
  \proofpartwideindR[Fixierter injektiver Schnitt]{%
    \begin{aligned}[t]
      &\forall K\colon M\inj M\,\bigl(K\colon M\bij M\bigr)
      \dsep F\colon M\sur M\dsep H\colon M\inj M\dsep{}\\
      &\forall y\in M\,\bigl(F(H(y))=y\bigr)
      \vdash F\colon M\bij M
    \end{aligned}%
  }{%
    \forall K\colon M\inj M\,\bigl(K\colon M\bij M\bigr)
    \dsep F\colon M\sur M\dsep H\colon M\inj M
    \dsep \forall y\in M\,\bigl(F(H(y))=y\bigr)
    \vdash F\colon M\bij M%
  }[SelfSurjectionFixedSectionBijective]
    \proofstepwidestar[1]{%
      \forall K\colon M\inj M\,\bigl(K\colon M\bij M\bigr)}{\rA}
    \proofstepwidestar[2]{F\colon M\sur M}{\rA}
    \proofstepwidestar[2]{F\colon M\to M}{%
      \FormulaRefAuto{Surjektive Funktion}{2}}
    \proofstepwidestar[4]{H\colon M\inj M}{\rA}
    \proofstepwidestar[5]{%
      \forall y\in M\,\bigl(F(H(y))=y\bigr)}{\rA}
    \proofstepwidestar[1,4]{H\colon M\bij M}{\rRE{\rUE{1},4}}
    \proofstepwidestar[1,4]{H^{-1}\colon M\to M}{%
      \FormulaRefAuto{F\colon A\bij B\vdash F^{-1}\colon B\to A}{6}}
    \proofstepwidestar[8]{x\in M}{\rA}
    \proofstepwidestar[1,4,8]{H^{-1}(x)\in M}{%
      \FormulaRefAuto{F\colon A\to B\dsep x\in A\vdash F(x)\in B}{7,8}}
    \proofstepwidestar[1,4,5,8]{%
      F\bigl(H(H^{-1}(x))\bigr)=H^{-1}(x)}{\rRE{9,\rUE{5}}}
    \proofstepwidestar[1,4,8]{H(H^{-1}(x))=x}{%
      \FormulaRefAuto{F\colon A\bij B \dsep y\in B
      \vdash F(F^{-1}(y)) = y}{6,8}}
    \proofstepwidestar[1,3,4,5,8]{F(x)=H^{-1}(x)}{\rIE{11,10}}
    \proofstepwidestar[1,3,4,5]{%
      \forall x\in M\,\bigl(F(x)=H^{-1}(x)\bigr)}{%
      \rUI{\rRI{8,12}}}
    \proofstepwidestar[1,3,4,5]{F=H^{-1}}{%
      \FormulaRefAuto{\forall x\in A (F(x)=G(x)) \eqvdash F=G}{13}}
    \proofstepwidestar[1,4]{H^{-1}\colon M\bij M}{%
      \FormulaRefAuto{F\colon A\bij B \vdash F^{-1}\colon B\bij A}{6}}
    \proofstepwidestar[1,3,4,5]{F\colon M\bij M}{\rIE{14,15}}
  \closeproofpartwideindR

  \proofpartwide{Schluss}
    \proofstepwidestar[1]{%
      \forall F\colon M\inj M\,\bigl(F\colon M\bij M\bigr)}{\rA}
    \proofstepwidestar[2]{F\colon M\sur M}{\rA}
    \proofstepwidestar[2]{%
      \exists H\colon M\inj M\,\forall y\in M\,\bigl(F(H(y))=y\bigr)
    }{%
      \FormulaRefAuto{%
        F\colon A\sur B\vdash
        \exists H\colon B\inj A\,\forall y\in B\,\bigl(F(H(y))=y\bigr)
      }{2}}
    \proofstepwidestar[4]{H\colon M\inj M}{\rA}
    \proofstepwidestar[5]{\forall y\in M\,\bigl(F(H(y))=y\bigr)}{\rA}
    \proofstepwidestar[1,2,4,5]{F\colon M\bij M}{%
      \FormulaRefAuto{SelfSurjectionFixedSectionBijective}{1,2,4,5}}
    \proofstepwidestar[1,2]{F\colon M\bij M}{\rEE{3,4,5,6}}
    \proofstepwidestar[1]{%
      \forall F\colon M\sur M\,\bigl(F\colon M\bij M\bigr)}{%
      \rUI{\rRI{2,7}}}
  \closeproofpartwide
\end{tabproofsplitwide}

\FormulaThmDelta{%
  \forall F\colon M\sur M\,\bigl(F\colon M\bij M\bigr)\vdash
  \forall F\colon M\inj M\,\bigl(F\colon M\bij M\bigr)
}{%
  \DeltaDecl{Mengen}{M\dsep y}%
  \DeltaDecl{Funktionen}{F\dsep H}%
}
\begin{tabproofsplitwide}
  \proofpartwideindR[Fall \(M=\varnothing\)]{%
    \forall F\colon M\sur M\,\bigl(F\colon M\bij M\bigr)\dsep
    F\colon M\inj M\dsep
    M=\varnothing
    \vdash F\colon M\bij M
  }{%
    \forall F\colon M\sur M\,\bigl(F\colon M\bij M\bigr)\dsep
    F\colon M\inj M\dsep
    M=\varnothing
    \vdash F\colon M\bij M
  }
    \proofstepwidestar[1]{\forall F\colon M\sur M\,\bigl(F\colon M\bij M\bigr)}{\rA}
    \proofstepwidestar[2]{F\colon M\inj M}{\rA}
    \proofstepwidestar[3]{M=\varnothing}{\rA}
    \proofstepwidestar[2]{F\colon M\to M}{\FormulaRefAuto{Injektive Funktion}{2}}
    \proofstepwidestar[2,3]{F\colon \varnothing\to \varnothing}{%
      \FormulaRefAuto{a=b\dsep c=d \dsep P(b,d)\vdash P(a,c)}{3,3,4}}
    \proofstepwidestar[2,3]{F=\varnothing}{%
      \FormulaRefAuto{f\colon \varnothing\to M \vdash f=\varnothing}{5}}
    \proofstepwidestar[]{\varnothing\colon \varnothing\bij \varnothing}{%
      \FormulaRefAuto{\varnothing\colon \varnothing\bij \varnothing}}
    \proofstepwidestar[3]{\varnothing\colon M\bij M}{%
      \FormulaRefAuto{a=b\dsep c=d \dsep P(b,d)\vdash P(a,c)}{3,3,7}}
    \proofstepwidestar[2,3]{F\colon M\bij M}{\rIE{6,8}}
  \closeproofpartwideindR

  \proofpartwideindR[Fall \(M\neq\varnothing\)]{%
    \forall F\colon M\sur M\,\bigl(F\colon M\bij M\bigr)\dsep
    F\colon M\inj M\dsep
    M\neq\varnothing
    \vdash F\colon M\bij M
  }{%
    \forall F\colon M\sur M\,\bigl(F\colon M\bij M\bigr)\dsep
    F\colon M\inj M\dsep
    M\neq\varnothing
    \vdash F\colon M\bij M
  }
    \proofstepwidestar[1]{\forall F\colon M\sur M\,\bigl(F\colon M\bij M\bigr)}{\rA}
    \proofstepwidestar[2]{F\colon M\inj M}{\rA}
    \proofstepwidestar[3]{M\neq\varnothing}{\rA}

    \proofstepwidestar[2,3]{%
      \exists H\colon M\sur M\,\forall y\in M\,\bigl(H(F(y))=y\bigr)
    }{%
      \FormulaRefAuto{F\colon A\inj B \dsep A\neq \varnothing \vdash \exists H\colon B\sur A\,\forall y\in A\,\bigl(H(F(y))=y\bigr)}{2,3}}

    \proofstepwidestar[4]{H\colon M\sur M}{\rA}
    \proofstepwidestar[5]{\forall y\in M\,\bigl(H(F(y))=y\bigr)}{\rA}

    \proofstepwidestar[1,4]{H\colon M\bij M}{\rRE{\rUE{1},5}}
    \proofstepwidestar[1,4]{H\colon M\inj M}{%
      \FormulaRefAuto{F\colon A\bij B \vdash F\colon A\inj B}{7}}

    \proofstepwidestar[9]{y\in M}{\rA}
    \proofstepwidestar[1,4,9]{H(y)\in M}{%
      \FormulaRefAuto{F\colon A\to B\dsep x\in A\vdash F(x)\in B}{%
        \FormulaRefAuto{Bijektive Funktion}{7},9}}
    \proofstepwidestar[1,4,9]{H(F(H(y)))=H(y)}{\rRE{10,\rUE{5}}}
    \proofstepwidestar[1,4,9]{F(H(y))\in M}{%
      \FormulaRefAuto{F\colon A\to B\dsep x\in A\vdash F(x)\in B}{%
        \FormulaRefAuto{Injektive Funktion}{2},10}}
    \proofstepwidestar[1,4,9]{F(H(y))=y}{%
      \FormulaRefAuto{F\colon A\inj B\dsep x\in A\dsep y\in A\dsep F(x)=F(y)\vdash x=y}{8,12,9,11}}
    \proofstepwidestar[1,4,9]{\exists x\in M\,\bigl(F(x)=y\bigr)}{\rEI{\rAI{10,13}}}
    \proofstepwidestar[1,4]{y\in M\rightarrow \exists x\in M\,\bigl(F(x)=y\bigr)}{\rRI{9,14}}
    \proofstepwidestar[1,2,4]{F\colon M\sur M}{%
      \FormulaRefAuto{Surjektive Funktion}{\FormulaRefAuto{Injektive Funktion}{2},15}}
    \proofstepwidestar[1,2,4]{F\colon M\bij M}{%
      \FormulaRefAuto{F\colon A\inj B\dsep F\colon A\sur B \vdash F\colon A\bij B}{2,16}}

    \proofstepwidestar[1,2,3]{F\colon M\bij M}{\rEE{4,5,17}}
  \closeproofpartwideindR

  \proofpartwide{Schluss}
    \proofstepwidestar[1]{\forall F\colon M\sur M\,\bigl(F\colon M\bij M\bigr)}{\rA}
    \proofstepwidestar[2]{F\colon M\inj M}{\rA}
    \proofstepwidestar[]{M=\varnothing \lor M\neq\varnothing}{\FormulaRefAuto{P \lor \neg P}}

    \proofstepwidestar[4]{M=\varnothing}{\rA}
    \proofstepwidestar[1,2,4]{F\colon M\bij M}{%
      \FormulaRefAuto{\forall F\colon M\sur M\,\bigl(F\colon M\bij M\bigr)\dsep F\colon M\inj M\dsep M=\varnothing \vdash F\colon M\bij M}{1,2,4}}

    \proofstepwidestar[6]{M\neq\varnothing}{\rA}
    \proofstepwidestar[1,2,6]{F\colon M\bij M}{%
      \FormulaRefAuto{\forall F\colon M\sur M\,\bigl(F\colon M\bij M\bigr)\dsep F\colon M\inj M\dsep M\neq\varnothing \vdash F\colon M\bij M}{1,2,6}}

    \proofstepwidestar[1,2]{F\colon M\bij M}{\rOE{3,4,5,6,7}}
    \proofstepwidestar[1]{%
      \forall F\colon M\inj M\,\bigl(F\colon M\bij M\bigr)
    }{\rUI{\rRI{2,8}}}
  \closeproofpartwide
\end{tabproofsplitwide}

\FormulaThmDelta{%
  \forall F\colon M\inj M\,\bigl(F\colon M\bij M\bigr)\vdash
  \neg\exists H\colon \mathbb{N}\inj M
}{%
  \DeltaDecl{Mengen}{M\dsep x}%
  \DeltaDecl{Funktionen}{F\dsep H}%
  \DeltaDecl{Funktionensymbol}{\SuccShift{H}}%
}
\begin{tabproof}
  \proofstep{1}{%
    \forall F\colon M\inj M\,\bigl(F\colon M\bij M\bigr)
  }{\rA}

  \proofstep{2}{\exists H\colon \mathbb{N}\inj M}{\rA}
  \proofstep{3}{H\colon \mathbb{N}\inj M}{\rA}

  \proofstep{3}{\SuccShift{H}\colon M\inj M}{%
    \FormulaRefAuto{H\colon \mathbb{N}\inj M \vdash \SuccShift{H}\colon M\inj M}{3}}

  \proofstep{1,3}{\SuccShift{H}\colon M\bij M}{%
    \rRE{\rUE{1},4}}

  \proofstep{1,3}{\SuccShift{H}\colon M\sur M}{%
    \FormulaRefAuto{F\colon A\bij B \vdash F\colon A\sur B}{5}}

  \proofstep{3}{\neg(\SuccShift{H}\colon M\sur M)}{%
    \FormulaRefAuto{H\colon \mathbb{N}\inj M \vdash \neg(\SuccShift{H}\colon M\sur M)}{3}}

  \proofstep{1,3}{\bot}{\rBI{6,7}}
  \proofstep{1,2}{\bot}{\rEE{2,3,8}}

  \proofstep{1}{\neg\exists H\colon \mathbb{N}\inj M}{\rCI{2,9}}
\end{tabproof}

\subsection{Tarskische Charakterisierung endlicher Mengen}

\subsubsection{Definition der Tarski-Endlichkeit}

\FormulaDefDeltaK[Begriff der Tarski-Endlichkeit]{%
  \TarskiFinite(M)
}{Tarski-endliche Menge}{
  \DeltaRow{Mengen}{M\dsep B\dsep x\dsep y}
  \DeltaRow{\textbf{Definitionsbedingung}}{}
  \DeltaRow{\makecell[l]{Tarskisches\\Maximalitätsaxiom}}
  {%
    \begin{aligned}[t]
      &B\subseteq \powerset(M)\dsep B\neq\varnothing \vdash{}\\
      &\exists x\in B\,\forall y\in B\,\bigl(x\subseteq y \rightarrow y=x\bigr)
    \end{aligned}
  }
  \DeltaRow{\textbf{Neue Symbole}}{}
  \DeltaPrem{Tarski-endliche Mengen}{ }
}

\subsubsection{Definition der Tarski-Endlichkeit über Minimalität}

\FormulaDefDeltaK[Begriff der minimalen Tarski-Endlichkeit]{%
  \TarskiMinFinite(M)
}{minimal Tarski-endliche Menge}{
  \DeltaRow{Mengen}{M\dsep B\dsep x\dsep y}
  \DeltaRow{\textbf{Definitionsbedingung}}{}
  \DeltaRow{\makecell[l]{Tarskisches\\Minimalitätsaxiom}}
  {%
    \begin{aligned}[t]
      &B\subseteq \powerset(M)\dsep B\neq\varnothing \vdash{}\\
      &\exists x\in B\,\forall y\in B\,\bigl(y\subseteq x \rightarrow y=x\bigr)
    \end{aligned}
  }
  \DeltaRow{\textbf{Neue Symbole}}{}
  \DeltaPrem{minimal Tarski-endliche Mengen}{ }
}

\subsubsection{Dedekind-Endlichkeit impliziert Tarski-Endlichkeit}

\paragraph{Beweisskizze}

\begingroup
\newcommand{\sketchref}[1]{\text{\scriptsize(\FormulaRefAuto{#1})}}
\newcommand{\sketchdefref}[1]{\text{\scriptsize(\FormulaRefAuto[env=definition]{#1})}}
\[
\begin{array}{c}
\text{\small Symbollegende}
\\[-.1ex]
\boxed{\begin{aligned}
\GrowRec(B,b_0,S)&:\quad
  b_0\in B,\quad S\colon B\to B,\quad
  \forall X\in B\,\bigl(X\subset S(X)\bigr)
\\[-.2ex]
&\qquad\sketchdefref{\GrowRec(B,b_0,S)\coloneqq b_0\in B\land S\colon B\to B\land \forall X\in B\,\bigl(X\subset S(X)\bigr)}
\\[.3ex]
\GrowFresh(M,B,b_0,S,D)&:\quad
  B\subseteq\powerset(M),\quad \GrowRec(B,b_0,S),\quad D\colon B\to M
\\[-.2ex]
&\quad\ \forall X\in B\,\bigl(D(X)\in S(X)\setminus X\bigr)
\quad
\sketchdefref{\GrowFresh(M,B,b_0,S,D)\coloneqq B\subseteq \powerset(M)\land \GrowRec(B,b_0,S)\land D\colon B\to M\land \forall X\in B\,\bigl(D(X)\in S(X)\setminus X\bigr)}
\end{aligned}}
\\[1.1ex]
\text{\small Kernbeweis}
\\[-.1ex]
\boxed{\begin{array}{@{}c@{}}
\neg\exists H\colon\mathbb{N}\inj M
\\[.6ex]
\boxed{\begin{gathered}
B\subseteq\powerset(M),\quad B\neq\varnothing\\
\text{kein Tarski-Maximum}
\end{gathered}}
\\[.4ex]
\Downarrow
\\[-.2ex]
\forall X\in B\,\exists Y\in B\,\bigl(X\subset Y\bigr)
\\[-.2ex]
\sketchref{B\subseteq \powerset(M)\dsep B\neq\varnothing\dsep \neg\exists x\in B\,\forall y\in B\,\bigl(x\subseteq y\rightarrow y=x\bigr)\vdash \forall x\in B\,\exists y\in B\,\bigl(x\subset y\bigr)}
\\[.4ex]
\Downarrow
\\[-.2ex]
\exists b_0\,\exists S\,\exists D\,\GrowFresh(M,B,b_0,S,D)
\\[-.2ex]
\sketchref{B\subseteq \powerset(M)\dsep B\neq\varnothing\dsep \forall x\in B\,\exists y\in B\,\bigl(x\subset y\bigr)\vdash \exists b_0\,\exists S\,\exists D\,\GrowFresh(M,B,b_0,S,D)}
\\[.4ex]
\Downarrow
\\[-.2ex]
\FreshSeq{b_0}{S}{D}\colon\mathbb{N}\inj M
\\[-.2ex]
\sketchref{\GrowFresh(M,B,b_0,S,D)\vdash \FreshSeq{b_0}{S}{D}\colon \mathbb{N}\inj M}
\\[.4ex]
\Downarrow\ \bot
\\[-.2ex]
\boxed{\TarskiFinite(M)}
\\[-.2ex]
\sketchref{\neg\exists H\colon \mathbb{N}\inj M\vdash \TarskiFinite(M)}
\end{array}}
\\[1.1ex]
\text{\small Schlussbeweis}
\\[-.1ex]
\boxed{\begin{array}{@{}c@{}}
\neg\exists H\colon\mathbb{N}\inj M
\\[.3ex]
\left[
\begin{gathered}
B\subseteq\powerset(M),\quad B\neq\varnothing,\\
\neg\exists x\in B\,\forall y\in B\,
\bigl(x\subseteq y\rightarrow y=x\bigr)
\end{gathered}
\right]
\Longrightarrow \exists H\colon\mathbb{N}\inj M
\\[-.2ex]
\sketchref{B\subseteq \powerset(M)\dsep B\neq\varnothing\dsep \neg\exists x\in B\,\forall y\in B\,\bigl(x\subseteq y\rightarrow y=x\bigr)\vdash \exists H\colon \mathbb{N}\inj M}
\\[.5ex]
\Downarrow\ \bot
\\[-.2ex]
B\subseteq\powerset(M),\ B\neq\varnothing
\vdash
\exists x\in B\,\forall y\in B\,
\bigl(x\subseteq y\rightarrow y=x\bigr)
\\[.4ex]
\Downarrow
\\[-.2ex]
\TarskiFinite(M)
\\[-.2ex]
\sketchref{\neg\exists H\colon \mathbb{N}\inj M\vdash \TarskiFinite(M)}
\end{array}}
\\[1.1ex]
\text{\small Stufenmechanik der frischen Werte}
\\[-.1ex]
\boxed{\begin{array}{@{}c@{}}
F_n\coloneqq\RecFun{b_0}{S}(n),\quad d_n\coloneqq D(F_n)
\\[.4ex]
F_0\subseteq F_1\subseteq F_2\subseteq \cdots
\subseteq F_n\subseteq F_{\Succ(n)}\subseteq\cdots
\\[-.1ex]
d_0\in F_1\setminus F_0,\quad
d_1\in F_2\setminus F_1,\quad
d_n\in F_{\Succ(n)}\setminus F_n
\\[-.2ex]
\sketchref{\GrowRec(B,b_0,S)\dsep u\in\mathbb{N}\vdash \RecFun{b_0}{S}(u)\subseteq \RecFun{b_0}{S}(\Succ(u))}
\quad
\sketchref{\GrowFresh(M,B,b_0,S,D)\dsep u\in\mathbb{N}\vdash D(\RecFun{b_0}{S}(u))\in \RecFun{b_0}{S}(\Succ(u))\land D(\RecFun{b_0}{S}(u))\notin \RecFun{b_0}{S}(u)}
\\[.6ex]
k<n\Longrightarrow
d_k\in F_n\land d_n\notin F_n
\Longrightarrow d_k\neq d_n
\\[-.2ex]
\sketchref{\GrowFresh(M,B,b_0,S,D)\dsep k\in\mathbb{N}\dsep n\in\mathbb{N}\dsep k < n\vdash D(\RecFun{b_0}{S}(k))\in \RecFun{b_0}{S}(n)}
\quad
\sketchref{\GrowFresh(M,B,b_0,S,D)\dsep k\in\mathbb{N}\dsep n\in\mathbb{N}\dsep k < n\vdash \FreshSeq{b_0}{S}{D}(k)\neq \FreshSeq{b_0}{S}{D}(n)}
\end{array}}
\end{array}
\]
\endgroup

\paragraph{Auswahlfunktionen aus nichtmaximalen Familien}

\subparagraph{Entfaltungslemmata fuer Obermengenwahlen}

\FormulaThmDeltaR[Obermengenwahl entpacken]{%
  \begin{aligned}[t]
    &B\subseteq \powerset(M)\dsep
    S\colon B\to B\dsep
    \forall x\in B\,\bigl((x,S(x))\in \StrictSupRel{B}\bigr)
    \vdash{}\\
    &S\colon B\to B\land \forall x\in B\,\bigl(x\subset S(x)\bigr)
  \end{aligned}
}{%
  B\subseteq \powerset(M)\dsep
  S\colon B\to B\dsep
  \forall x\in B\,\bigl((x,S(x))\in \StrictSupRel{B}\bigr)
  \vdash
  S\colon B\to B\land \forall x\in B\,\bigl(x\subset S(x)\bigr)
}{%
  \DeltaRow{Mengen}{M\dsep B\dsep x}
  \DeltaRow{Funktionen}{S}
  \DeltaRow{Relationen}{\StrictSupRel{B}}
}
\begin{tabproof}
  \proofstep{1}{B\subseteq \powerset(M)}{\rA}
  \proofstep{2}{S\colon B\to B}{\rA}
  \proofstep{3}{\forall x\in B\,\bigl((x,S(x))\in \StrictSupRel{B}\bigr)}{\rA}
  \proofstep{4}{x\in B}{\rA}
  \proofstep{3,4}{(x,S(x))\in \StrictSupRel{B}}{\rRE{4,\rUE{3}}}
  \proofstep{1,3,4}{x\in B\land S(x)\in B\land x\subset S(x)}{%
    \FormulaRefAuto{B\subseteq \powerset(M)\dsep (x,y)\in \StrictSupRel{B}\vdash x\in B\land y\in B\land x\subset y}{1,5}}
  \proofstep{1,3,4}{x\subset S(x)}{\rAEb{6}}
  \proofstep{1,3}{\forall x\in B\,\bigl(x\subset S(x)\bigr)}{\rUI{\rRI{4,7}}}
  \proofstep{1,2,3}{S\colon B\to B\land \forall x\in B\,\bigl(x\subset S(x)\bigr)}{\rAI{2,8}}
\end{tabproof}

\subparagraph{Konstruktives Auswahllemma fuer Obermengen}

\FormulaThmDelta[Existenz einer strikten Obermengenwahl]{%
  B\subseteq \powerset(M)\dsep
  \forall x\in B\,\exists y\in B\,\bigl(x\subset y\bigr)
  \vdash
  \exists S\colon B\to B\,\forall x\in B\,\bigl(x\subset S(x)\bigr)
}{%
  \DeltaRow{Mengen}{M\dsep B\dsep x\dsep y}
  \DeltaRow{Funktionen}{S}
  \DeltaRow{Relationen}{\StrictSupRel{B}}
}
\begin{tabproof}
  \proofstep{1}{B\subseteq \powerset(M)}{\rA}
  \proofstep{2}{\forall x\in B\,\exists y\in B\,\bigl(x\subset y\bigr)}{\rA}
  \proofstep{1,2}{\TotRel{\StrictSupRel{B},B,B}}{%
    \FormulaRefAuto{B\subseteq \powerset(M)\dsep \forall x\in B\,\exists y\in B\,\bigl(x\subset y\bigr)\vdash \TotRel{\StrictSupRel{B},B,B}}{1,2}}
  \proofstep{1,2}{%
    \exists S\colon B\to B\,\forall x\in B\,\bigl((x,S(x))\in \StrictSupRel{B}\bigr)
  }{%
    \FormulaRefAuto{\exists F\colon A \to B\,\forall x\in A\,\bigl((x,F(x))\in R\bigr)}{3}}

  \proofstep{5}{S\colon B\to B\land \forall x\in B\,\bigl((x,S(x))\in \StrictSupRel{B}\bigr)}{\rA}
  \proofstep{5}{S\colon B\to B}{\rAEa{5}}
  \proofstep{5}{\forall x\in B\,\bigl((x,S(x))\in \StrictSupRel{B}\bigr)}{\rAEb{5}}
  \proofstep{1,5}{S\colon B\to B\land \forall x\in B\,\bigl(x\subset S(x)\bigr)}{%
    \FormulaRefAuto{B\subseteq \powerset(M)\dsep S\colon B\to B\dsep \forall x\in B\,\bigl((x,S(x))\in \StrictSupRel{B}\bigr)\vdash S\colon B\to B\land \forall x\in B\,\bigl(x\subset S(x)\bigr)}{1,6,7}}
  \proofstep{1,5}{\exists S\colon B\to B\,\forall x\in B\,\bigl(x\subset S(x)\bigr)}{\rEI{8}}
  \proofstep{1,2}{\exists S\colon B\to B\,\forall x\in B\,\bigl(x\subset S(x)\bigr)}{\rEE{4,5,9}}
\end{tabproof}

\subparagraph{Frischelementrelation}

\FormulaDefDelta[Frischelementrelation zu einer strikten Obermengenwahl]{%
  B\subseteq \powerset(M)\dsep S\colon B\to B \vdash
  \FreshRel{S}\coloneqq
  \{(X,a)\in B\times M\mid a\in S(X)\setminus X\}%
}{%
  \DeltaRow{Mengen}{M\dsep B\dsep X\dsep a}
  \DeltaRow{Funktionen}{S}
  \DeltaRow{Relationen}{\FreshRel{S}}
}

\subparagraph{Entfaltungslemmata der Frischelementrelation}

\FormulaThmDelta[Typisierung der Frischelementrelation]{%
  B\subseteq \powerset(M)\dsep
  S\colon B\to B
  \vdash
  \FreshRel{S}\subseteq B\times M
}{%
  \DeltaRow{Mengen}{M\dsep B}
  \DeltaRow{Funktionen}{S}
  \DeltaRow{Relationen}{\FreshRel{S}}
}
\begin{tabproof}
  \proofstep{1}{B\subseteq \powerset(M)}{\rA}
  \proofstep{2}{S\colon B\to B}{\rA}
  \proofstep{1,2}{\FreshRel{S}\subseteq B\times M}{%
    \FormulaRefAuto{A=\{ x \in B \mid P(x)\}\vdash  A\subseteq B}{%
      \FormulaRefAuto{B\subseteq \powerset(M)\dsep S\colon B\to B \vdash \FreshRel{S}\coloneqq \{(X,a)\in B\times M\mid a\in S(X)\setminus X\}}{1,2}}}
\end{tabproof}

\FormulaThmDeltaR[Frischelementdaten typisieren Paar]{%
  \begin{aligned}[t]
    &B\subseteq \powerset(M)\dsep
    S\colon B\to B\dsep X\in B\dsep
    a\in S(X)\setminus X
    \vdash{}\\
    &(X,a)\in B\times M
  \end{aligned}
}{%
  B\subseteq \powerset(M)\dsep
  S\colon B\to B\dsep X\in B\dsep
  a\in S(X)\setminus X
  \vdash (X,a)\in B\times M
}{%
  \DeltaRow{Mengen}{M\dsep B\dsep X\dsep a}
  \DeltaRow{Funktionen}{S}
}
\begin{tabproof}
  \proofstep{1}{B\subseteq \powerset(M)}{\rA}
  \proofstep{2}{S\colon B\to B}{\rA}
  \proofstep{3}{X\in B}{\rA}
  \proofstep{4}{a\in S(X)\setminus X}{\rA}
  \proofstep{2,3}{S(X)\in B}{%
    \FormulaRefAuto{F\colon A\to B\dsep x\in A\vdash F(x)\in B}{2,3}}
  \proofstep{1,5}{S(X)\in \powerset(M)}{%
    \FormulaRefAuto{A\subseteq B,\, x\in A \vdash x\in B}{1,5}}
  \proofstep{6}{S(X)\subseteq M}{%
    \FormulaRefAuto{x \in \powerset(A)\eqvdash x \subseteq A}{6}}
  \proofstep{4}{a\in S(X)}{\FormulaRefAuto{x\in A\setminus B\vdash x\in A}{4}}
  \proofstep{1,2,3,4}{a\in M}{%
    \FormulaRefAuto{A\subseteq B,\, x\in A \vdash x\in B}{7,8}}
  \proofstep{3,9}{(X,a)\in B\times M}{%
    \FormulaRefAuto{a\in A\dsep b\in B\vdash (a,b)\in A\times B}{3,9}}
\end{tabproof}

\FormulaThmDelta[Frischelementrelation einfuehren]{%
  B\subseteq \powerset(M)\dsep
  S\colon B\to B\dsep
  X\in B\dsep
  a\in S(X)\setminus X
  \vdash
  (X,a)\in \FreshRel{S}
}{%
  \DeltaRow{Mengen}{M\dsep B\dsep X\dsep a}
  \DeltaRow{Funktionen}{S}
  \DeltaRow{Relationen}{\FreshRel{S}}
}
\begin{tabproof}
  \proofstep{1}{B\subseteq \powerset(M)}{\rA}
  \proofstep{2}{S\colon B\to B}{\rA}
  \proofstep{3}{X\in B}{\rA}
  \proofstep{4}{a\in S(X)\setminus X}{\rA}
  \proofstep{1,2,3,4}{(X,a)\in B\times M}{%
    \FormulaRefAuto{B\subseteq \powerset(M)\dsep S\colon B\to B\dsep X\in B\dsep a\in S(X)\setminus X\vdash (X,a)\in B\times M}{1,2,3,4}}
  \proofstep{1,2,3,4}{(X,a)\in B\times M\land a\in S(X)\setminus X}{\rAI{5,4}}
  \proofstep{1,2,3,4}{(X,a)\in \FreshRel{S}}{%
    \rIE{\FormulaRefAuto{B\subseteq \powerset(M)\dsep S\colon B\to B \vdash \FreshRel{S}\coloneqq \{(X,a)\in B\times M\mid a\in S(X)\setminus X\}}{1,2},6}}
\end{tabproof}

\FormulaThmDelta[Frischelementrelation liefert Frischelement]{%
  B\subseteq \powerset(M)\dsep
  S\colon B\to B\dsep
  (X,a)\in \FreshRel{S}
  \vdash
  a\in S(X)\setminus X
}{%
  \DeltaRow{Mengen}{M\dsep B\dsep X\dsep a}
  \DeltaRow{Funktionen}{S}
  \DeltaRow{Relationen}{\FreshRel{S}}
}
\begin{tabproof}
  \proofstep{1}{B\subseteq \powerset(M)}{\rA}
  \proofstep{2}{S\colon B\to B}{\rA}
  \proofstep{3}{(X,a)\in \FreshRel{S}}{\rA}
  \proofstep{1,2,3}{(X,a)\in \{(X,a)\in B\times M\mid a\in S(X)\setminus X\}}{%
    \rIE{\FormulaRefAuto{B\subseteq \powerset(M)\dsep S\colon B\to B \vdash \FreshRel{S}\coloneqq \{(X,a)\in B\times M\mid a\in S(X)\setminus X\}}{1,2},3}}
  \proofstep{1,2,3}{a\in S(X)\setminus X}{%
    \FormulaRefAuto{x \in \{x \in A \mid P(x)\}\vdash P(x)}{4}}
\end{tabproof}

\subparagraph{Konstruktive Frischelementlemmata}

\FormulaThmDelta[Existenz eines Frischelements]{%
  B\subseteq \powerset(M)\dsep
  S\colon B\to B\dsep
  \forall X\in B\,\bigl(X\subset S(X)\bigr)\dsep
  X\in B
  \vdash
  \exists a\,\bigl((X,a)\in \FreshRel{S}\bigr)
}{%
  \DeltaRow{Mengen}{M\dsep B\dsep X\dsep a}
  \DeltaRow{Funktionen}{S}
  \DeltaRow{Relationen}{\FreshRel{S}}
}
\begin{tabproof}
  \proofstep{1}{B\subseteq \powerset(M)}{\rA}
  \proofstep{2}{S\colon B\to B}{\rA}
  \proofstep{3}{\forall X\in B\,\bigl(X\subset S(X)\bigr)}{\rA}
  \proofstep{4}{X\in B}{\rA}
  \proofstep{3,4}{X\subset S(X)}{\rRE{4,\rUE{3}}}
  \proofstep{3,4}{\exists a\in S(X)\setminus X}{%
    \FormulaRefAuto{A\subset B\vdash \exists x\in B\setminus A}{5}}
  \proofstep{7}{a\in S(X)\setminus X}{\rA}
  \proofstep{1,2,4,7}{(X,a)\in \FreshRel{S}}{%
    \FormulaRefAuto{B\subseteq \powerset(M)\dsep S\colon B\to B\dsep X\in B\dsep a\in S(X)\setminus X\vdash (X,a)\in \FreshRel{S}}{1,2,4,7}}
  \proofstep{1,2,4,7}{\exists a\,\bigl((X,a)\in \FreshRel{S}\bigr)}{\rEI{8}}
  \proofstep{1,2,3,4}{\exists a\,\bigl((X,a)\in \FreshRel{S}\bigr)}{\rEE{6,7,9}}
\end{tabproof}

\FormulaThmDelta[Totalitaet der Frischelementrelation]{%
  B\subseteq \powerset(M)\dsep
  S\colon B\to B\dsep
  \forall X\in B\,\bigl(X\subset S(X)\bigr)
  \vdash
  \TotRel{\FreshRel{S},B,M}
}{%
  \DeltaRow{Mengen}{M\dsep B\dsep X\dsep a}
  \DeltaRow{Funktionen}{S}
  \DeltaRow{Relationen}{\FreshRel{S}}
}
\begin{tabproof}
  \proofstep{1}{B\subseteq \powerset(M)}{\rA}
  \proofstep{2}{S\colon B\to B}{\rA}
  \proofstep{3}{\forall X\in B\,\bigl(X\subset S(X)\bigr)}{\rA}
  \proofstep{1,2}{\FreshRel{S}\subseteq B\times M}{%
    \FormulaRefAuto{B\subseteq \powerset(M)\dsep S\colon B\to B\vdash \FreshRel{S}\subseteq B\times M}{1,2}}
  \proofstep{5}{X\in B}{\rA}
  \proofstep{1,2,3,5}{\exists a\,\bigl((X,a)\in \FreshRel{S}\bigr)}{%
    \FormulaRefAuto{B\subseteq \powerset(M)\dsep S\colon B\to B\dsep \forall X\in B\,\bigl(X\subset S(X)\bigr)\dsep X\in B\vdash \exists a\,\bigl((X,a)\in \FreshRel{S}\bigr)}{1,2,3,5}}
  \proofstep{1,2,3}{X\in B\rightarrow \exists a\,(X,a)\in \FreshRel{S}}{\rRI{5,6}}
  \proofstep{1,2,3}{\TotRel{\FreshRel{S},B,M}}{%
    \FormulaRefAuto{Totale Relation}{4,7}}
\end{tabproof}

\subparagraph{Auswahl frischer Elemente}

\FormulaThmDeltaR[Frischelementauswahl entpacken]{%
  \begin{aligned}[t]
    &B\subseteq \powerset(M)\dsep
    S\colon B\to B\dsep D\colon B\to M\dsep{}\\
    &\forall X\in B\,\bigl((X,D(X))\in \FreshRel{S}\bigr)
    \vdash{}\\
    &D\colon B\to M\land
    \forall X\in B\,\bigl(D(X)\in S(X)\setminus X\bigr)
  \end{aligned}
}{%
  B\subseteq \powerset(M)\dsep
  S\colon B\to B\dsep
  D\colon B\to M\dsep
  \forall X\in B\,\bigl((X,D(X))\in \FreshRel{S}\bigr)
  \vdash
  D\colon B\to M\land
  \forall X\in B\,\bigl(D(X)\in S(X)\setminus X\bigr)
}{%
  \DeltaRow{Mengen}{M\dsep B\dsep X}
  \DeltaRow{Funktionen}{S\dsep D}
  \DeltaRow{Relationen}{\FreshRel{S}}
}
\begin{tabproof}
  \proofstep{1}{B\subseteq \powerset(M)}{\rA}
  \proofstep{2}{S\colon B\to B}{\rA}
  \proofstep{3}{D\colon B\to M}{\rA}
  \proofstep{4}{\forall X\in B\,\bigl((X,D(X))\in \FreshRel{S}\bigr)}{\rA}
  \proofstep{5}{X\in B}{\rA}
  \proofstep{4,5}{(X,D(X))\in \FreshRel{S}}{\rRE{5,\rUE{4}}}
  \proofstep{1,2,4,5}{D(X)\in S(X)\setminus X}{%
    \FormulaRefAuto{B\subseteq \powerset(M)\dsep S\colon B\to B\dsep (X,a)\in \FreshRel{S}\vdash a\in S(X)\setminus X}{1,2,6}}
  \proofstep{1,2,4}{\forall X\in B\,\bigl(D(X)\in S(X)\setminus X\bigr)}{\rUI{\rRI{5,7}}}
  \proofstep{1,2,3,4}{D\colon B\to M\land \forall X\in B\,\bigl(D(X)\in S(X)\setminus X\bigr)}{\rAI{3,8}}
\end{tabproof}

\FormulaThmDeltaR[Ausgewaehlte Frischrelation liefert Frischelementauswahl]{%
  \begin{aligned}[t]
    &B\subseteq \powerset(M)\dsep S\colon B\to B\dsep{}\\
    &\exists D\colon B\to M\,
      \forall X\in B\,\bigl((X,D(X))\in \FreshRel{S}\bigr)
    \vdash{}\\
    &\exists D\colon B\to M\,
      \forall X\in B\,\bigl(D(X)\in S(X)\setminus X\bigr)
  \end{aligned}
}{%
  B\subseteq \powerset(M)\dsep S\colon B\to B\dsep
  \exists D\colon B\to M\,\forall X\in B\,\bigl((X,D(X))\in \FreshRel{S}\bigr)
  \vdash
  \exists D\colon B\to M\,\forall X\in B\,\bigl(D(X)\in S(X)\setminus X\bigr)
}{%
  \DeltaRow{Mengen}{M\dsep B\dsep X}
  \DeltaRow{Funktionen}{S\dsep D}
  \DeltaRow{Relationen}{\FreshRel{S}}
}
\begin{tabproof}
  \proofstep{1}{B\subseteq \powerset(M)}{\rA}
  \proofstep{2}{S\colon B\to B}{\rA}
  \proofstep{3}{\exists D\colon B\to M\,\forall X\in B\,\bigl((X,D(X))\in \FreshRel{S}\bigr)}{\rA}
  \proofstep{4}{D\colon B\to M\land \forall X\in B\,\bigl((X,D(X))\in \FreshRel{S}\bigr)}{\rA}
  \proofstep{4}{D\colon B\to M}{\rAEa{4}}
  \proofstep{4}{\forall X\in B\,\bigl((X,D(X))\in \FreshRel{S}\bigr)}{\rAEb{4}}
  \proofstep{1,2,4}{D\colon B\to M\land \forall X\in B\,\bigl(D(X)\in S(X)\setminus X\bigr)}{%
    \FormulaRefAuto{B\subseteq \powerset(M)\dsep S\colon B\to B\dsep D\colon B\to M\dsep \forall X\in B\,\bigl((X,D(X))\in \FreshRel{S}\bigr)\vdash D\colon B\to M\land \forall X\in B\,\bigl(D(X)\in S(X)\setminus X\bigr)}{1,2,5,6}}
  \proofstep{1,2,4}{\exists D\colon B\to M\,\forall X\in B\,\bigl(D(X)\in S(X)\setminus X\bigr)}{\rEI{7}}
  \proofstep{1,2,3}{\exists D\colon B\to M\,\forall X\in B\,\bigl(D(X)\in S(X)\setminus X\bigr)}{\rEE{3,4,8}}
\end{tabproof}

\FormulaThmDelta[Existenz einer Frischelementauswahl]{%
  B\subseteq \powerset(M)\dsep
  S\colon B\to B\dsep
  \forall X\in B\,\bigl(X\subset S(X)\bigr)
  \vdash
  \exists D\colon B\to M\,\forall X\in B\,\bigl(D(X)\in S(X)\setminus X\bigr)
}{%
  \DeltaRow{Mengen}{M\dsep B\dsep X\dsep a}
  \DeltaRow{Funktionen}{S\dsep D}
  \DeltaRow{Relationen}{\FreshRel{S}}
}
\begin{tabproof}
  \proofstep{1}{B\subseteq \powerset(M)}{\rA}
  \proofstep{2}{S\colon B\to B}{\rA}
  \proofstep{3}{\forall X\in B\,\bigl(X\subset S(X)\bigr)}{\rA}
  \proofstep{1,2,3}{\TotRel{\FreshRel{S},B,M}}{%
    \FormulaRefAuto{B\subseteq \powerset(M)\dsep S\colon B\to B\dsep \forall X\in B\,\bigl(X\subset S(X)\bigr)\vdash \TotRel{\FreshRel{S},B,M}}{1,2,3}}
  \proofstep{1,2,3}{%
    \exists D\colon B\to M\,\forall X\in B\,\bigl((X,D(X))\in \FreshRel{S}\bigr)
  }{%
    \FormulaRefAuto{\exists F\colon A \to B\,\forall x\in A\,\bigl((x,F(x))\in R\bigr)}{4}}
  \proofstep{1,2,3}{\exists D\colon B\to M\,\forall X\in B\,\bigl(D(X)\in S(X)\setminus X\bigr)}{%
    \FormulaRefAuto{B\subseteq \powerset(M)\dsep S\colon B\to B\dsep \exists D\colon B\to M\,\forall X\in B\,\bigl((X,D(X))\in \FreshRel{S}\bigr)\vdash \exists D\colon B\to M\,\forall X\in B\,\bigl(D(X)\in S(X)\setminus X\bigr)}{1,2,5}}
\end{tabproof}

\paragraph{Nichtmaximale Familien und Dedekind-Unendlichkeit}

\subparagraph{Abkuerzungen und Projektionen}

\FormulaDefDeltaR[Wachstumssituation]{%
  \begin{aligned}[t]
    \GrowRec(B,b_0,S)\coloneqq{}&
    b_0\in B\land S\colon B\to B\land{}\\
    &\forall X\in B\,\bigl(X\subset S(X)\bigr)
  \end{aligned}
}{%
  \GrowRec(B,b_0,S)\coloneqq
  b_0\in B\land S\colon B\to B\land
  \forall X\in B\,\bigl(X\subset S(X)\bigr)
}{%
  \DeltaRow{Mengen}{B\dsep X\dsep b_0}
  \DeltaRow{Funktionen}{S}
}

\FormulaDefDeltaR[Wachstumssituation mit Frischelementauswahl]{%
  \begin{aligned}[t]
    \GrowFresh(M,B,b_0,S,D)\coloneqq{}&
    B\subseteq \powerset(M)\land \GrowRec(B,b_0,S)\land{}\\
    &D\colon B\to M\land
    \forall X\in B\,\bigl(D(X)\in S(X)\setminus X\bigr)
  \end{aligned}
}{%
  \GrowFresh(M,B,b_0,S,D)\coloneqq
  B\subseteq \powerset(M)\land \GrowRec(B,b_0,S)\land
  D\colon B\to M\land
  \forall X\in B\,\bigl(D(X)\in S(X)\setminus X\bigr)
}{%
  \DeltaRow{Mengen}{M\dsep B\dsep X\dsep b_0}
  \DeltaRow{Funktionen}{S\dsep D}
}

\FormulaDefDelta[Frischwertfolge]{%
  \FreshSeq{b_0}{S}{D}\coloneqq D\circ \RecFun{b_0}{S}
}{%
  \DeltaRow{Mengen}{B\dsep b_0}
  \DeltaRow{Funktionen}{S\dsep D}
}

\FormulaThmDelta[Frischdaten typisieren Rekursionsstufen]{%
  \GrowFresh(M,B,b_0,S,D)\dsep u\in\mathbb{N}\vdash
  \RecFun{b_0}{S}(u)\in B
}{%
  \DeltaDecl{Mengen}{M\dsep B\dsep b_0\dsep u}%
  \DeltaDecl{Funktionen}{S\dsep D}%
}
\begin{tabproof}
  \proofstep{1}{\GrowFresh(M,B,b_0,S,D)}{\rA}
  \proofstep{2}{u\in\mathbb{N}}{\rA}
  \proofstep{1}{\GrowRec(B,b_0,S)}{%
    \FormulaRefAuto{\GrowFresh(M,B,b_0,S,D)\coloneqq B\subseteq \powerset(M)\land \GrowRec(B,b_0,S)\land D\colon B\to M\land \forall X\in B\,\bigl(D(X)\in S(X)\setminus X\bigr)}{1}}
  \proofstep{1}{b_0\in B}{%
    \FormulaRefAuto{\GrowRec(B,b_0,S)\coloneqq b_0\in B\land S\colon B\to B\land \forall X\in B\,\bigl(X\subset S(X)\bigr)}{3}}
  \proofstep{1}{S\colon B\to B}{%
    \FormulaRefAuto{\GrowRec(B,b_0,S)\coloneqq b_0\in B\land S\colon B\to B\land \forall X\in B\,\bigl(X\subset S(X)\bigr)}{3}}
  \proofstep{1,2}{\RecFun{b_0}{S}(u)\in B}{%
    \FormulaRefAuto{m_0\in M\dsep f\colon M\to M\dsep n\in\mathbb{N}\vdash \RecFun{m_0}{f}(n)\in M}{4,5,2}}
\end{tabproof}

\FormulaThmDelta[Frischdaten liefern die Rekursionsgleichung]{%
  \GrowFresh(M,B,b_0,S,D)\dsep u\in\mathbb{N}\vdash
  \RecFun{b_0}{S}(\Succ(u))=S(\RecFun{b_0}{S}(u))
}{%
  \DeltaDecl{Mengen}{M\dsep B\dsep b_0\dsep u}%
  \DeltaDecl{Funktionen}{S\dsep D}%
}
\begin{tabproof}
  \proofstep{1}{\GrowFresh(M,B,b_0,S,D)}{\rA}
  \proofstep{2}{u\in\mathbb{N}}{\rA}
  \proofstep{1}{\GrowRec(B,b_0,S)}{%
    \FormulaRefAuto{\GrowFresh(M,B,b_0,S,D)\coloneqq B\subseteq \powerset(M)\land \GrowRec(B,b_0,S)\land D\colon B\to M\land \forall X\in B\,\bigl(D(X)\in S(X)\setminus X\bigr)}{1}}
  \proofstep{1}{b_0\in B}{%
    \FormulaRefAuto{\GrowRec(B,b_0,S)\coloneqq b_0\in B\land S\colon B\to B\land \forall X\in B\,\bigl(X\subset S(X)\bigr)}{3}}
  \proofstep{1}{S\colon B\to B}{%
    \FormulaRefAuto{\GrowRec(B,b_0,S)\coloneqq b_0\in B\land S\colon B\to B\land \forall X\in B\,\bigl(X\subset S(X)\bigr)}{3}}
  \proofstep{1,2}{%
    \RecFun{b_0}{S}(\Succ(u))=S(\RecFun{b_0}{S}(u))
  }{\FormulaRefAuto{m_0\in M\dsep f\colon M\to M\dsep n\in\mathbb{N}\vdash \RecFun{m_0}{f}(\Succ(n))=f(\RecFun{m_0}{f}(n))}{4,5,2}}
\end{tabproof}

\FormulaThmDeltaR[Frischdaten typisieren die Frischwertfolge]{%
  \begin{aligned}[t]
    &\GrowFresh(M,B,b_0,S,D)\vdash{}\\
    &\FreshSeq{b_0}{S}{D}\colon \mathbb{N}\to M
  \end{aligned}
}{%
  \GrowFresh(M,B,b_0,S,D)\vdash
  \FreshSeq{b_0}{S}{D}\colon \mathbb{N}\to M
}{%
  \DeltaDecl{Mengen}{M\dsep B\dsep b_0}%
  \DeltaDecl{Funktionen}{S\dsep D}%
}
\begin{tabproof}
  \proofstep{1}{\GrowFresh(M,B,b_0,S,D)}{\rA}
  \proofstep{1}{\GrowRec(B,b_0,S)}{%
    \FormulaRefAuto{\GrowFresh(M,B,b_0,S,D)\coloneqq B\subseteq \powerset(M)\land \GrowRec(B,b_0,S)\land D\colon B\to M\land \forall X\in B\,\bigl(D(X)\in S(X)\setminus X\bigr)}{1}}
  \proofstep{1}{b_0\in B}{%
    \FormulaRefAuto{\GrowRec(B,b_0,S)\coloneqq b_0\in B\land S\colon B\to B\land \forall X\in B\,\bigl(X\subset S(X)\bigr)}{2}}
  \proofstep{1}{S\colon B\to B}{%
    \FormulaRefAuto{\GrowRec(B,b_0,S)\coloneqq b_0\in B\land S\colon B\to B\land \forall X\in B\,\bigl(X\subset S(X)\bigr)}{2}}
  \proofstep{1}{D\colon B\to M}{%
    \FormulaRefAuto{\GrowFresh(M,B,b_0,S,D)\coloneqq B\subseteq \powerset(M)\land \GrowRec(B,b_0,S)\land D\colon B\to M\land \forall X\in B\,\bigl(D(X)\in S(X)\setminus X\bigr)}{1}}
  \proofstep{1}{\RecFun{b_0}{S}\colon \mathbb{N}\to B}{%
    \FormulaRefAuto{m_0\in M\dsep f\colon M\to M \vdash \RecFun{m_0}{f}\colon \mathbb{N}\to M}{3,4}}
  \proofstep{1}{D\circ \RecFun{b_0}{S}\colon \mathbb{N}\to M}{%
    \FormulaRefAuto{G\circ F\colon A\to C}{6,5}}
  \proofstep{1}{\FreshSeq{b_0}{S}{D}\colon \mathbb{N}\to M}{%
    \rIE{\FormulaRefAuto{\FreshSeq{b_0}{S}{D}\coloneqq D\circ \RecFun{b_0}{S}}{},7}}
\end{tabproof}

\subparagraph{Technischer Kern der Konstruktion}

\FormulaThmDeltaR[Ein-Schritt-Monotonie der wachsenden Rekursion]{%
  \begin{aligned}[t]
    &\GrowRec(B,b_0,S)\dsep u\in\mathbb{N}\vdash{}\\
    &\RecFun{b_0}{S}(u)\subseteq \RecFun{b_0}{S}(\Succ(u))
  \end{aligned}
}{%
  \GrowRec(B,b_0,S)\dsep u\in\mathbb{N}\vdash
  \RecFun{b_0}{S}(u)\subseteq \RecFun{b_0}{S}(\Succ(u))
}{%
  \DeltaDecl{Mengen}{B\dsep X\dsep b_0\dsep u}%
  \DeltaDecl{Funktionen}{S}%
}
\begin{tabproof}
  \proofstep{1}{\GrowRec(B,b_0,S)}{\rA}
  \proofstep{2}{u\in\mathbb{N}}{\rA}
  \proofstep{1}{b_0\in B}{%
    \FormulaRefAuto{\GrowRec(B,b_0,S)\coloneqq b_0\in B\land S\colon B\to B\land \forall X\in B\,\bigl(X\subset S(X)\bigr)}{1}}
  \proofstep{1}{S\colon B\to B}{%
    \FormulaRefAuto{\GrowRec(B,b_0,S)\coloneqq b_0\in B\land S\colon B\to B\land \forall X\in B\,\bigl(X\subset S(X)\bigr)}{1}}
  \proofstep{1}{\forall X\in B\,\bigl(X\subset S(X)\bigr)}{%
    \FormulaRefAuto{\GrowRec(B,b_0,S)\coloneqq b_0\in B\land S\colon B\to B\land \forall X\in B\,\bigl(X\subset S(X)\bigr)}{1}}
  \proofstep{1,2}{\RecFun{b_0}{S}(u)\in B}{%
    \FormulaRefAuto{m_0\in M\dsep f\colon M\to M\dsep n\in\mathbb{N}\vdash \RecFun{m_0}{f}(n)\in M}{3,4,2}}
  \proofstep{1,2}{%
    \RecFun{b_0}{S}(u)\subset S(\RecFun{b_0}{S}(u))
  }{\rRE{6,\rUE{5}}}
  \proofstep{1,2}{%
    \RecFun{b_0}{S}(u)\subseteq S(\RecFun{b_0}{S}(u))
  }{\FormulaRefAuto{A\subset B\vdash A\subseteq B}{7}}
  \proofstep{1,2}{%
    \RecFun{b_0}{S}(\Succ(u))=S(\RecFun{b_0}{S}(u))
  }{\FormulaRefAuto{m_0\in M\dsep f\colon M\to M\dsep n\in\mathbb{N}\vdash \RecFun{m_0}{f}(\Succ(n))=f(\RecFun{m_0}{f}(n))}{3,4,2}}
  \proofstep{1,2}{%
    S(\RecFun{b_0}{S}(u))=\RecFun{b_0}{S}(\Succ(u))
  }{\FormulaRefAuto{a=b\vdash b=a}{9}}
  \proofstep{1,2}{%
    \RecFun{b_0}{S}(u)\subseteq \RecFun{b_0}{S}(\Succ(u))
  }{\FormulaRefAuto{A \subseteq B, B = C \vdash A \subseteq C}{8,10}}
\end{tabproof}

\FormulaThmDeltaR[Monotonie der wachsenden Rekursion]{%
  \begin{aligned}[t]
    &\GrowRec(B,b_0,S)\dsep
    k\in\mathbb{N}\dsep n\in\mathbb{N}\dsep k < n
    \vdash{}\\
    &
    \RecFun{b_0}{S}(\Succ(k))\subseteq \RecFun{b_0}{S}(n)
  \end{aligned}
}{%
  \GrowRec(B,b_0,S)\dsep
  k\in\mathbb{N}\dsep n\in\mathbb{N}\dsep k < n
  \vdash
  \RecFun{b_0}{S}(\Succ(k))\subseteq \RecFun{b_0}{S}(n)
}{%
  \DeltaDecl{Mengen}{B\dsep X\dsep b_0\dsep k\dsep n\dsep u}%
  \DeltaDecl{Funktionen}{S}%
  \DeltaDecl{Relation}{<_{\mathbb{N}}}%
}
\begin{tabproof}
  \proofstep{1}{\GrowRec(B,b_0,S)}{\rA}
  \proofstep{2}{k\in\mathbb{N}}{\rA}
  \proofstep{3}{n\in\mathbb{N}}{\rA}
  \proofstep{4}{k<n}{\rA}
  \proofstep{5}{u\in\mathbb{N}}{\rA}
  \proofstep{1,5}{%
    \RecFun{b_0}{S}(u)\subseteq \RecFun{b_0}{S}(\Succ(u))
  }{%
    \FormulaRefAuto{\GrowRec(B,b_0,S)\dsep u\in\mathbb{N}\vdash \RecFun{b_0}{S}(u)\subseteq \RecFun{b_0}{S}(\Succ(u))}{1,5}}
  \proofstep{1}{%
    \forall u\in\mathbb{N}\,\bigl(
      \RecFun{b_0}{S}(u)\subseteq \RecFun{b_0}{S}(\Succ(u))
    \bigr)
  }{\rUI{\rRI{5,6}}}
  \proofstep{1,2,3,4}{%
    \RecFun{b_0}{S}(\Succ(k))\subseteq \RecFun{b_0}{S}(n)
  }{%
    \FormulaRefAuto{B04PeanoStrictMonotoneSucc}[thm]{7,2,3,4}}
\end{tabproof}

\FormulaThmDeltaR[Frischwert liegt im Zuwachs]{%
  \begin{aligned}[t]
    &\GrowFresh(M,B,b_0,S,D)\dsep u\in\mathbb{N}\vdash{}\\
    &D(\RecFun{b_0}{S}(u))\in \RecFun{b_0}{S}(\Succ(u))
      \land
      D(\RecFun{b_0}{S}(u))\notin \RecFun{b_0}{S}(u)
  \end{aligned}
}{%
  \GrowFresh(M,B,b_0,S,D)\dsep u\in\mathbb{N}\vdash
  D(\RecFun{b_0}{S}(u))\in \RecFun{b_0}{S}(\Succ(u))
  \land
  D(\RecFun{b_0}{S}(u))\notin \RecFun{b_0}{S}(u)
}{%
  \DeltaDecl{Mengen}{M\dsep B\dsep X\dsep b_0\dsep u}%
  \DeltaDecl{Funktionen}{S\dsep D}%
}
\begin{tabproof}
  \proofstep{1}{\GrowFresh(M,B,b_0,S,D)}{\rA}
  \proofstep{2}{u\in\mathbb{N}}{\rA}
  \proofstep{1}{\forall X\in B\,\bigl(D(X)\in S(X)\setminus X\bigr)}{%
    \FormulaRefAuto{\GrowFresh(M,B,b_0,S,D)\coloneqq B\subseteq \powerset(M)\land \GrowRec(B,b_0,S)\land D\colon B\to M\land \forall X\in B\,\bigl(D(X)\in S(X)\setminus X\bigr)}{1}}
  \proofstep{1,2}{\RecFun{b_0}{S}(u)\in B}{%
    \FormulaRefAuto{\GrowFresh(M,B,b_0,S,D)\dsep u\in\mathbb{N}\vdash \RecFun{b_0}{S}(u)\in B}{1,2}}
  \proofstep{1,2}{%
    \begin{aligned}[t]
      D(\RecFun{b_0}{S}(u))&\in{}\\
      &S(\RecFun{b_0}{S}(u))\setminus \RecFun{b_0}{S}(u)
    \end{aligned}
  }{%
    \rRE{4,\rUE{3}}}
  \proofstep{1,2}{%
    \begin{aligned}[t]
      D(\RecFun{b_0}{S}(u))&\in S(\RecFun{b_0}{S}(u))
    \end{aligned}
  }{%
    \FormulaRefAuto{x\in A\setminus B\vdash x\in A}{5}}
  \proofstep{1,2}{%
    \begin{aligned}[t]
      \RecFun{b_0}{S}(\Succ(u))
      &=S(\RecFun{b_0}{S}(u))
    \end{aligned}
  }{\FormulaRefAuto{\GrowFresh(M,B,b_0,S,D)\dsep u\in\mathbb{N}\vdash \RecFun{b_0}{S}(\Succ(u))=S(\RecFun{b_0}{S}(u))}{1,2}}
  \proofstep{1,2}{%
    \begin{aligned}[t]
      S(\RecFun{b_0}{S}(u))
      &=\RecFun{b_0}{S}(\Succ(u))
    \end{aligned}
  }{\FormulaRefAuto{a=b\vdash b=a}{7}}
  \proofstep{1,2}{%
    \begin{aligned}[t]
      D(\RecFun{b_0}{S}(u))
      &\in \RecFun{b_0}{S}(\Succ(u))
    \end{aligned}
  }{%
    \rIE{8,6}}
  \proofstep{1,2}{%
    \begin{aligned}[t]
      D(\RecFun{b_0}{S}(u))
      &\notin \RecFun{b_0}{S}(u)
    \end{aligned}
  }{%
    \FormulaRefAuto{x\in A\setminus B\vdash x\notin B}{5}}
  \proofstep{1,2}{%
    \begin{aligned}[t]
      &D(\RecFun{b_0}{S}(u))\in \RecFun{b_0}{S}(\Succ(u))\\
      &\land\,
      D(\RecFun{b_0}{S}(u))\notin \RecFun{b_0}{S}(u)
    \end{aligned}
  }{\rAI{9,10}}
\end{tabproof}

\FormulaThmDeltaR[Frueher Frischwert liegt in spaeterer Stufe]{%
  \begin{aligned}[t]
    &\GrowFresh(M,B,b_0,S,D)\dsep
    k\in\mathbb{N}\dsep n\in\mathbb{N}\dsep k < n
    \vdash{}\\
    &D(\RecFun{b_0}{S}(k))\in \RecFun{b_0}{S}(n)
  \end{aligned}
}{%
  \GrowFresh(M,B,b_0,S,D)\dsep
  k\in\mathbb{N}\dsep n\in\mathbb{N}\dsep k < n
  \vdash
  D(\RecFun{b_0}{S}(k))\in \RecFun{b_0}{S}(n)
}{%
  \DeltaDecl{Mengen}{M\dsep B\dsep b_0\dsep k\dsep n}%
  \DeltaDecl{Funktionen}{S\dsep D}%
  \DeltaDecl{Relation}{<_{\mathbb{N}}}%
}
\begin{tabproof}
  \proofstep{1}{\GrowFresh(M,B,b_0,S,D)}{\rA}
  \proofstep{2}{k\in\mathbb{N}}{\rA}
  \proofstep{3}{n\in\mathbb{N}}{\rA}
  \proofstep{4}{k<n}{\rA}
  \proofstep{1}{\GrowRec(B,b_0,S)}{%
    \FormulaRefAuto{\GrowFresh(M,B,b_0,S,D)\coloneqq B\subseteq \powerset(M)\land \GrowRec(B,b_0,S)\land D\colon B\to M\land \forall X\in B\,\bigl(D(X)\in S(X)\setminus X\bigr)}{1}}
  \proofstep{1,2}{D(\RecFun{b_0}{S}(k))\in \RecFun{b_0}{S}(\Succ(k))}{%
    \rAEa{\FormulaRefAuto{\GrowFresh(M,B,b_0,S,D)\dsep u\in\mathbb{N}\vdash D(\RecFun{b_0}{S}(u))\in \RecFun{b_0}{S}(\Succ(u))\land D(\RecFun{b_0}{S}(u))\notin \RecFun{b_0}{S}(u)}{1,2}}}
  \proofstep{1,2,3,4}{%
    \RecFun{b_0}{S}(\Succ(k))\subseteq \RecFun{b_0}{S}(n)
  }{%
    \FormulaRefAuto{\GrowRec(B,b_0,S)\dsep k\in\mathbb{N}\dsep n\in\mathbb{N}\dsep k < n\vdash \RecFun{b_0}{S}(\Succ(k))\subseteq \RecFun{b_0}{S}(n)}{5,2,3,4}}
  \proofstep{1,2,3,4}{D(\RecFun{b_0}{S}(k))\in \RecFun{b_0}{S}(n)}{%
    \FormulaRefAuto{A\subseteq B,\, x\in A \vdash x\in B}{7,6}}
\end{tabproof}

\FormulaThmDeltaR[Frische Werte bleiben verschieden]{%
  \begin{aligned}[t]
    &\GrowFresh(M,B,b_0,S,D)\dsep
    k\in\mathbb{N}\dsep n\in\mathbb{N}\dsep k < n
    \vdash{}\\
    &\FreshSeq{b_0}{S}{D}(k)\neq \FreshSeq{b_0}{S}{D}(n)
  \end{aligned}
}{%
  \GrowFresh(M,B,b_0,S,D)\dsep
  k\in\mathbb{N}\dsep n\in\mathbb{N}\dsep k < n
  \vdash
  \FreshSeq{b_0}{S}{D}(k)\neq \FreshSeq{b_0}{S}{D}(n)
}{%
  \DeltaDecl{Mengen}{M\dsep B\dsep X\dsep b_0\dsep k\dsep n}%
  \DeltaDecl{Funktionen}{S\dsep D}%
  \DeltaDecl{Relation}{<_{\mathbb{N}}}%
}
\begin{tabproof}
  \proofstep{1}{\GrowFresh(M,B,b_0,S,D)}{\rA}
  \proofstep{2}{k\in\mathbb{N}}{\rA}
  \proofstep{3}{n\in\mathbb{N}}{\rA}
  \proofstep{4}{k < n}{\rA}
  \proofstep{1,2,3,4}{D(\RecFun{b_0}{S}(k))\in \RecFun{b_0}{S}(n)}{%
    \FormulaRefAuto{\GrowFresh(M,B,b_0,S,D)\dsep k\in\mathbb{N}\dsep n\in\mathbb{N}\dsep k < n\vdash D(\RecFun{b_0}{S}(k))\in \RecFun{b_0}{S}(n)}{1,2,3,4}}
  \proofstep{1,3}{D(\RecFun{b_0}{S}(n))\notin \RecFun{b_0}{S}(n)}{%
    \rAEb{\FormulaRefAuto{\GrowFresh(M,B,b_0,S,D)\dsep u\in\mathbb{N}\vdash D(\RecFun{b_0}{S}(u))\in \RecFun{b_0}{S}(\Succ(u))\land D(\RecFun{b_0}{S}(u))\notin \RecFun{b_0}{S}(u)}{1,3}}}
  \proofstep{1,2,3,4}{D(\RecFun{b_0}{S}(k))\neq D(\RecFun{b_0}{S}(n))}{%
    \FormulaRefAuto{x\in A\dsep y\notin A \vdash x\neq y}{5,6}}
  \proofstep{1,2,3,4}{%
    (D\circ \RecFun{b_0}{S})(k)\neq
    (D\circ \RecFun{b_0}{S})(n)
  }{%
    \FormulaRefAuto{x\in A\dsep y\in A\dsep G(F(x))\neq G(F(y))\vdash (G\circ F)(x)\neq (G\circ F)(y)}{2,3,7}}
  \proofstep{1,2,3,4}{%
    \FreshSeq{b_0}{S}{D}(k)\neq \FreshSeq{b_0}{S}{D}(n)
  }{%
    \rIE{\FormulaRefAuto{\FreshSeq{b_0}{S}{D}\coloneqq D\circ \RecFun{b_0}{S}}{},8}}
\end{tabproof}

\FormulaThmDeltaR[Frischwertfolge ist ordnungsgetrennt]{%
  \begin{aligned}[t]
    &\GrowFresh(M,B,b_0,S,D)\vdash{}\\
    &\forall k,n\in\mathbb{N}\,\bigl(k<n\rightarrow
      \FreshSeq{b_0}{S}{D}(k)\neq \FreshSeq{b_0}{S}{D}(n)\bigr)
  \end{aligned}
}{%
  \GrowFresh(M,B,b_0,S,D)\vdash
  \forall k,n\in\mathbb{N}\,\bigl(k<n\rightarrow
    \FreshSeq{b_0}{S}{D}(k)\neq \FreshSeq{b_0}{S}{D}(n)\bigr)
}{%
  \DeltaDecl{Mengen}{M\dsep B\dsep b_0\dsep k\dsep n}%
  \DeltaDecl{Funktionen}{S\dsep D}%
  \DeltaDecl{Relation}{<_{\mathbb{N}}}%
}
\begin{tabproof}
  \proofstep{1}{\GrowFresh(M,B,b_0,S,D)}{\rA}
  \proofstep{2}{k\in\mathbb{N}}{\rA}
  \proofstep{3}{n\in\mathbb{N}}{\rA}
  \proofstep{4}{k<n}{\rA}
  \proofstep{1,2,3,4}{%
    \FreshSeq{b_0}{S}{D}(k)\neq \FreshSeq{b_0}{S}{D}(n)
  }{%
    \FormulaRefAuto{\GrowFresh(M,B,b_0,S,D)\dsep k\in\mathbb{N}\dsep n\in\mathbb{N}\dsep k < n\vdash \FreshSeq{b_0}{S}{D}(k)\neq \FreshSeq{b_0}{S}{D}(n)}{1,2,3,4}}
  \proofstep{1,2,3}{%
    k<n\rightarrow
    \FreshSeq{b_0}{S}{D}(k)\neq \FreshSeq{b_0}{S}{D}(n)
  }{\rRI{4,5}}
  \proofstep{1}{%
    \forall k,n\in\mathbb{N}\,\bigl(k<n\rightarrow
      \FreshSeq{b_0}{S}{D}(k)\neq \FreshSeq{b_0}{S}{D}(n)\bigr)
  }{\rUI{\rRI{2,\rRI{3,6}}}}
\end{tabproof}

\FormulaThmDeltaR[Frische Rekursion ist injektiv]{%
  \begin{aligned}[t]
    &\GrowFresh(M,B,b_0,S,D)\vdash{}\\
    &\FreshSeq{b_0}{S}{D}\colon \mathbb{N}\inj M
  \end{aligned}
}{%
  \GrowFresh(M,B,b_0,S,D)\vdash
  \FreshSeq{b_0}{S}{D}\colon \mathbb{N}\inj M
}{%
  \DeltaDecl{Mengen}{M\dsep B\dsep X\dsep b_0\dsep n\dsep k}%
  \DeltaDecl{Funktionen}{S\dsep D}%
}
\begin{tabproof}
  \proofstep{1}{\GrowFresh(M,B,b_0,S,D)}{\rA}
  \proofstep{1}{\FreshSeq{b_0}{S}{D}\colon \mathbb{N}\to M}{%
    \FormulaRefAuto{\GrowFresh(M,B,b_0,S,D)\vdash \FreshSeq{b_0}{S}{D}\colon \mathbb{N}\to M}{1}}
  \proofstep{1}{%
    \forall k,n\in\mathbb{N}\,\bigl(k<n\rightarrow
      \FreshSeq{b_0}{S}{D}(k)\neq \FreshSeq{b_0}{S}{D}(n)\bigr)
  }{%
    \FormulaRefAuto{\GrowFresh(M,B,b_0,S,D)\vdash \forall k,n\in\mathbb{N}\,\bigl(k<n\rightarrow \FreshSeq{b_0}{S}{D}(k)\neq \FreshSeq{b_0}{S}{D}(n)\bigr)}{1}}
  \proofstep{1}{\FreshSeq{b_0}{S}{D}\colon \mathbb{N}\inj M}{%
    \FormulaRefAuto{B04PeanoStrictSeparationInjection}[thm]{2,3}}
\end{tabproof}

\subparagraph{Anwendung auf Tarski-Endlichkeit}

\paragraph{Abhängigkeit vom Auswahlaxiom.}
Die im nächsten Satz gleichzeitig gewählten Erweiterungen \(S(X)\) und
Frischelemente \(D(X)\) werden mit dem bereits eingeführten
\FormulaRefAuto{Auswahlaxiom}[ax] gewonnen. Damit hängen der folgende
Frischdaten-Satz und folglich die Dedekind-Rückrichtung zur Endlichkeit in
der hier gegebenen Beweisroute vom Auswahlaxiom ab.

\FormulaThmDeltaR[Nichtmaximalitaet liefert Frischdaten]{%
  \begin{aligned}[t]
    &B\subseteq \powerset(M)\dsep B\neq\varnothing\dsep{}\\
    &\forall x\in B\,\exists y\in B\,\bigl(x\subset y\bigr)
    \vdash{}\\
    &\exists b_0\,\exists S\,\exists D\,\GrowFresh(M,B,b_0,S,D)
  \end{aligned}
}{%
  B\subseteq \powerset(M)\dsep B\neq\varnothing\dsep
  \forall x\in B\,\exists y\in B\,\bigl(x\subset y\bigr)
  \vdash \exists b_0\,\exists S\,\exists D\,\GrowFresh(M,B,b_0,S,D)
}{%
  \DeltaDecl{Mengen}{M\dsep B\dsep b_0\dsep X\dsep x\dsep y}%
  \DeltaDecl{Funktionen}{S\dsep D}%
}
\begin{tabproof}
  \proofstep{1}{B\subseteq \powerset(M)}{\rA}
  \proofstep{2}{B\neq\varnothing}{\rA}
  \proofstep{3}{\forall x\in B\,\exists y\in B\,\bigl(x\subset y\bigr)}{\rA}
  \proofstep{2}{\exists b_0\in B}{%
    \FormulaRefAuto{B\neq\varnothing\vdash \exists x\in B}{2}}
  \proofstep{1,3}{\exists S\colon B\to B\,\forall x\in B\,\bigl(x\subset S(x)\bigr)}{%
    \FormulaRefAuto{B\subseteq \powerset(M)\dsep \forall x\in B\,\exists y\in B\,\bigl(x\subset y\bigr)\vdash \exists S\colon B\to B\,\forall x\in B\,\bigl(x\subset S(x)\bigr)}{1,3},
    \FormulaRefAuto{Auswahlaxiom}[ax]}
  \proofstep{6}{b_0\in B}{\rA}
  \proofstep{7}{S\colon B\to B\land \forall x\in B\,\bigl(x\subset S(x)\bigr)}{\rA}
  \proofstep{6,7}{\GrowRec(B,b_0,S)}{%
    \rAI{6,\rAEa{7},\rAEb{7}}}
  \proofstep{1,7}{%
    \begin{aligned}[t]
      &\exists D\colon B\to M\,\forall X\in B\,\\
      &\bigl(D(X)\in S(X)\setminus X\bigr)
    \end{aligned}
  }{%
    \FormulaRefAuto{B\subseteq \powerset(M)\dsep S\colon B\to B\dsep \forall X\in B\,\bigl(X\subset S(X)\bigr)\vdash \exists D\colon B\to M\,\forall X\in B\,\bigl(D(X)\in S(X)\setminus X\bigr)}{1,\rAEa{7},\rAEb{7}},
    \FormulaRefAuto{Auswahlaxiom}[ax]}
  \proofstep{10}{%
    \begin{aligned}[t]
      &D\colon B\to M\land{}\\
      &\forall X\in B\,\bigl(D(X)\in S(X)\setminus X\bigr)
    \end{aligned}
  }{\rA}
  \proofstep{1,6,7,10}{\GrowFresh(M,B,b_0,S,D)}{%
    \rAI{1,8,\rAEa{10},\rAEb{10}}}
  \proofstep{1,6,7,10}{\exists D\,\GrowFresh(M,B,b_0,S,D)}{\rEI{11}}
  \proofstep{1,6,7,10}{\exists S\,\exists D\,\GrowFresh(M,B,b_0,S,D)}{\rEI{12}}
  \proofstep{1,6,7,10}{\exists b_0\,\exists S\,\exists D\,\GrowFresh(M,B,b_0,S,D)}{\rEI{13}}
  \proofstep{1,6,7}{\exists b_0\,\exists S\,\exists D\,\GrowFresh(M,B,b_0,S,D)}{\rEE{9,10,14}}
  \proofstep{1,3,6}{\exists b_0\,\exists S\,\exists D\,\GrowFresh(M,B,b_0,S,D)}{\rEE{5,7,15}}
  \proofstep{1,2,3}{\exists b_0\,\exists S\,\exists D\,\GrowFresh(M,B,b_0,S,D)}{\rEE{4,6,16}}
\end{tabproof}

\FormulaThmDeltaR[Nichtmaximale Potenzmengenfamilien sind Dedekind-unendlich]{%
  \begin{aligned}[t]
    &B\subseteq \powerset(M)\dsep B\neq\varnothing\dsep{}\\
    &\forall x\in B\,\exists y\in B\,\bigl(x\subset y\bigr)
    \vdash \exists H\colon \mathbb{N}\inj M
  \end{aligned}
}{%
  B\subseteq \powerset(M)\dsep B\neq\varnothing\dsep
  \forall x\in B\,\exists y\in B\,\bigl(x\subset y\bigr)
  \vdash \exists H\colon \mathbb{N}\inj M
}{%
  \DeltaDecl{Mengen}{M\dsep B\dsep b_0\dsep x\dsep y}%
  \DeltaDecl{Funktionen}{S\dsep D\dsep H}%
}
\begin{tabproof}
  \proofstep{1}{B\subseteq \powerset(M)}{\rA}
  \proofstep{2}{B\neq\varnothing}{\rA}
  \proofstep{3}{\forall x\in B\,\exists y\in B\,\bigl(x\subset y\bigr)}{\rA}
  \proofstep{1,2,3}{\exists b_0\,\exists S\,\exists D\,\GrowFresh(M,B,b_0,S,D)}{%
    \FormulaRefAuto{B\subseteq \powerset(M)\dsep B\neq\varnothing\dsep \forall x\in B\,\exists y\in B\,\bigl(x\subset y\bigr)\vdash \exists b_0\,\exists S\,\exists D\,\GrowFresh(M,B,b_0,S,D)}{1,2,3}}
  \proofstep{5}{\GrowFresh(M,B,b_0,S,D)}{\rA}
  \proofstep{5}{\FreshSeq{b_0}{S}{D}\colon \mathbb{N}\inj M}{%
    \FormulaRefAuto{\GrowFresh(M,B,b_0,S,D)\vdash \FreshSeq{b_0}{S}{D}\colon \mathbb{N}\inj M}{5}}
  \proofstep{5}{\exists H\colon \mathbb{N}\inj M}{\rEI{6}}
  \proofstep{1,2,3}{\exists H\colon \mathbb{N}\inj M}{\rEE{4,5,7}}
\end{tabproof}

\subparagraph{Tarski-Gegenbeispiele als Nichtmaximalitaet}

\FormulaThmDeltaR[Kein Tarski-Maximum erzwingt echte Erweiterungen]{%
  \begin{aligned}[t]
    &B\subseteq \powerset(M)\dsep B\neq\varnothing\dsep{}\\
    &\neg\exists x\in B\,\forall y\in B\,
      \bigl(x\subseteq y\rightarrow y=x\bigr)
    \vdash{}\\
    &\forall x\in B\,\exists y\in B\,\bigl(x\subset y\bigr)
  \end{aligned}
}{%
  B\subseteq \powerset(M)\dsep B\neq\varnothing\dsep
  \neg\exists x\in B\,\forall y\in B\,\bigl(x\subseteq y\rightarrow y=x\bigr)
  \vdash
  \forall x\in B\,\exists y\in B\,\bigl(x\subset y\bigr)
}{%
  \DeltaDecl{Mengen}{M\dsep B\dsep x\dsep y}%
}
\begin{tabproof}
  \proofstep{1}{B\subseteq \powerset(M)}{\rA}
  \proofstep{2}{B\neq\varnothing}{\rA}
  \proofstep{3}{\neg\exists x\in B\,\forall y\in B\,\bigl(x\subseteq y\rightarrow y=x\bigr)}{\rA}
  \proofstep{4}{x\in B}{\rA}
  \proofstep{5}{\neg\exists y\in B\,\bigl(x\subset y\bigr)}{\rA}
  \proofstep{6}{y\in B}{\rA}
  \proofstep{7}{x\subseteq y}{\rA}
  \proofstep{8}{y\neq x}{\rA}
  \proofstep{7,8}{x\subset y}{%
    \FormulaRefAuto{A \subset B \coloneq A \subseteq B \land A\neq B}{%
      \rAI{7,\FormulaRefAuto{a \neq b \vdash b \neq a}{8}}}}
  \proofstep{6,7,8}{\exists y\in B\,\bigl(x\subset y\bigr)}{\rEI{\rAI{6,9}}}
  \proofstep{5,6,7,8}{\bot}{\rBI{5,10}}
  \proofstep{5,6,7}{y=x}{\rCE{8,11}}
  \proofstep{5,6}{x\subseteq y\rightarrow y=x}{\rRI{7,12}}
  \proofstep{4,5}{\forall y\in B\,\bigl(x\subseteq y\rightarrow y=x\bigr)}{\rUI{\rRI{6,13}}}
  \proofstep{4,5}{\exists x\in B\,\forall y\in B\,\bigl(x\subseteq y\rightarrow y=x\bigr)}{\rEI{\rAI{4,14}}}
  \proofstep{3,4,5}{\bot}{\rBI{3,15}}
  \proofstep{3,4}{\exists y\in B\,\bigl(x\subset y\bigr)}{\rCE{5,16}}
  \proofstep{1,2,3}{\forall x\in B\,\exists y\in B\,\bigl(x\subset y\bigr)}{\rUI{\rRI{4,17}}}
\end{tabproof}

\FormulaThmDeltaR[Tarski-Gegenbeispiele erzwingen Dedekind-Unendlichkeit]{%
  \begin{aligned}[t]
    &B\subseteq \powerset(M)\dsep B\neq\varnothing\dsep{}\\
    &\neg\exists x\in B\,\forall y\in B\,
      \bigl(x\subseteq y\rightarrow y=x\bigr)
    \vdash{}\\
    &\exists H\colon \mathbb{N}\inj M
  \end{aligned}
}{%
  B\subseteq \powerset(M)\dsep B\neq\varnothing\dsep
  \neg\exists x\in B\,\forall y\in B\,\bigl(x\subseteq y\rightarrow y=x\bigr)
  \vdash \exists H\colon \mathbb{N}\inj M
}{%
  \DeltaDecl{Mengen}{M\dsep B\dsep x\dsep y}%
  \DeltaDecl{Funktionen}{H}%
}
\begin{tabproof}
  \proofstep{1}{B\subseteq \powerset(M)}{\rA}
  \proofstep{2}{B\neq\varnothing}{\rA}
  \proofstep{3}{\neg\exists x\in B\,\forall y\in B\,\bigl(x\subseteq y\rightarrow y=x\bigr)}{\rA}
  \proofstep{1,2,3}{\forall x\in B\,\exists y\in B\,\bigl(x\subset y\bigr)}{%
    \FormulaRefAuto{B\subseteq \powerset(M)\dsep B\neq\varnothing\dsep \neg\exists x\in B\,\forall y\in B\,\bigl(x\subseteq y\rightarrow y=x\bigr)\vdash \forall x\in B\,\exists y\in B\,\bigl(x\subset y\bigr)}{1,2,3}}
  \proofstep{1,2,3}{\exists H\colon \mathbb{N}\inj M}{%
    \FormulaRefAuto{B\subseteq \powerset(M)\dsep B\neq\varnothing\dsep \forall x\in B\,\exists y\in B\,\bigl(x\subset y\bigr)\vdash \exists H\colon \mathbb{N}\inj M}{1,2,4}}
\end{tabproof}

\paragraph{Schluss: Dedekind-Endlichkeit impliziert Tarski-Endlichkeit}

\FormulaThmDelta[Dedekind-Endlichkeit impliziert Tarski-Endlichkeit]{%
  \neg\exists H\colon \mathbb{N}\inj M\vdash \TarskiFinite(M)
}{%
  \DeltaDecl{Mengen}{M\dsep B\dsep x\dsep y}%
  \DeltaDecl{Funktionen}{H}%
}
\begin{tabproof}
  \proofstep{1}{\neg\exists H\colon \mathbb{N}\inj M}{\rA}
  \proofstep{2}{B\subseteq \powerset(M)\land B\neq\varnothing}{\rA}
  \proofstep{2}{B\subseteq \powerset(M)}{\rAEa{2}}
  \proofstep{2}{B\neq\varnothing}{\rAEb{2}}
  \proofstep{5}{\neg\exists x\in B\,\forall y\in B\,\bigl(x\subseteq y\rightarrow y=x\bigr)}{\rA}
  \proofstep{2,5}{\exists H\colon \mathbb{N}\inj M}{%
    \FormulaRefAuto{B\subseteq \powerset(M)\dsep B\neq\varnothing\dsep \neg\exists x\in B\,\forall y\in B\,\bigl(x\subseteq y\rightarrow y=x\bigr)\vdash \exists H\colon \mathbb{N}\inj M}{3,4,5}}
  \proofstep{1,2,5}{\bot}{\rBI{1,6}}
  \proofstep{1,2}{\exists x\in B\,\forall y\in B\,\bigl(x\subseteq y\rightarrow y=x\bigr)}{\rCE{5,7}}
  \proofstep{1}{%
    \bigl(B\subseteq \powerset(M)\land B\neq\varnothing\bigr)\rightarrow
    \exists x\in B\,\forall y\in B\,\bigl(x\subseteq y\rightarrow y=x\bigr)
  }{\rRI{2,8}}
  \proofstep{1}{\TarskiFinite(M)}{\FormulaRefAuto{Tarski-endliche Menge}[def]{9}}
\end{tabproof}

\subsubsection{Tarski-Endlichkeit impliziert Endlichkeit}

\paragraph{Beweisskizze}

\begingroup
\newcommand{\sketchdefref}[1]{\text{\scriptsize(\FormulaRefAuto[env=definition]{#1})}}
\[
\begin{array}{@{}c@{}}
\FinSubsets{M}\coloneq \{F\in\powerset(M)\mid \Finite{F}\}
\quad
\sketchdefref{\FinSubsets{M}\coloneq \{F\in\powerset(M)\mid \Finite{F}\}}
\\[1ex]
\boxed{
\TarskiFinite(M)
\Longrightarrow
\exists F\in\FinSubsets{M}\text{ maximal}
\Longrightarrow
F=M
\Longrightarrow
\Finite{M}}
\end{array}
\]
\endgroup

Der einzige eigentliche Widerspruchsschritt liegt in der Gleichung \(F=M\):
Wäre \(F\neq M\), so gäbe es ein \(x\in M\setminus F\). Dann wäre
\(F\cup\{x\}\) wieder eine endliche Teilmenge von \(M\), aber echt größer
als \(F\), im Widerspruch zur Maximalität von \(F\).

Formalisiert werden nun nur drei Punkte: \(\FinSubsets{M}\) ist eine
nichtleere Teilfamilie von \(\powerset(M)\), frische Adjunktion bleibt endlich,
und ein maximales endliches \(F\subseteq M\) muss schon ganz \(M\) sein.

\paragraph{Endliche Teilmengen}

\FormulaDefDelta[Familie der endlichen Teilmengen]{%
  \FinSubsets{M}\coloneq \{F\in\powerset(M)\mid \Finite{F}\}
}{%
  \DeltaDecl{Mengen}{M\dsep F}%
}

\FormulaThmDelta[Charakterisierung der Familie der endlichen Teilmengen]{%
  F\in\FinSubsets{M}\eqvdash F\subseteq M\land \Finite{F}
}{%
  \DeltaDecl{Mengen}{M\dsep F}%
}
\begin{tabproof}
  \proofstep{}{F\in\FinSubsets{M}\eqvdash F\in\powerset(M)\land \Finite{F}}{%
    \FormulaRefAuto{\FinSubsets{M}\coloneq \{F\in\powerset(M)\mid \Finite{F}\}}}
  \proofstep{}{F\in\powerset(M)\eqvdash F\subseteq M}{%
    \FormulaRefAuto{x \in \powerset(A)\;\eqvdash\; x \subseteq A}}
  \proofstep{}{F\in\FinSubsets{M}\eqvdash F\subseteq M\land \Finite{F}}{\rChain{1,2}}
\end{tabproof}

\FormulaThmDelta[Endliche Teilmengen bilden eine nichtleere Potenzmengenfamilie]{%
  \FinSubsets{M}\subseteq \powerset(M)\land \FinSubsets{M}\neq\varnothing
}{%
  \DeltaDecl{Mengen}{M\dsep F}%
}
\begin{tabproof}
  \proofstep{1}{F\in\FinSubsets{M}}{\rA}
  \proofstep{1}{F\subseteq M\land \Finite{F}}{%
    \FormulaRefAuto{F\in\FinSubsets{M}\eqvdash F\subseteq M\land \Finite{F}}{1}}
  \proofstep{1}{F\subseteq M}{\rAEa{2}}
  \proofstep{1}{F\in\powerset(M)}{%
    \FormulaRefAuto{x \in \powerset(A)\;\eqvdash\; x \subseteq A}{3}}
  \proofstep{}{\FinSubsets{M}\subseteq \powerset(M)}{%
    \FormulaRefAuto{A \subseteq B \coloneq \forall x\,(x\in A \rightarrow x\in B)}{\rUI{\rRI{1,4}}}}
  \proofstep{}{\varnothing\subseteq M}{\FormulaRefAuto{\varnothing\subseteq A}}
  \proofstep{}{\Finite{\varnothing}}{\FormulaRefAuto{FiniteEmpty}[thm]}
  \proofstep{}{\varnothing\in\FinSubsets{M}}{%
    \FormulaRefAuto{F\in\FinSubsets{M}\eqvdash F\subseteq M\land \Finite{F}}{\rAI{6,7}}}
  \proofstep{}{\FinSubsets{M}\neq\varnothing}{%
    \FormulaRefAuto{a \in A \vdash A \neq \varnothing}{8}}
  \proofstep{}{\FinSubsets{M}\subseteq \powerset(M)\land \FinSubsets{M}\neq\varnothing}{\rAI{5,9}}
\end{tabproof}

\FormulaThmDelta[Frische Adjunktion endlicher Teilmengen]{%
  F\in\FinSubsets{M}\dsep x\in M\setminus F
  \vdash F\cup\{x\}\in\FinSubsets{M}\land F\subset F\cup\{x\}
}{%
  \DeltaDecl{Mengen}{M\dsep F\dsep x}%
}
\begin{tabproof}
  \proofstep{1}{F\in\FinSubsets{M}}{\rA}
  \proofstep{2}{x\in M\setminus F}{\rA}
  \proofstep{1}{F\subseteq M\land \Finite{F}}{%
    \FormulaRefAuto{F\in\FinSubsets{M}\eqvdash F\subseteq M\land \Finite{F}}{1}}
  \proofstep{1}{F\subseteq M}{\rAEa{3}}
  \proofstep{1}{\Finite{F}}{\rAEb{3}}
  \proofstep{2}{x\in M}{\FormulaRefAuto{x \in A \setminus B \vdash x \in A}{2}}
  \proofstep{2}{x\notin F}{\FormulaRefAuto{x\in A\setminus B\vdash x\notin B}{2}}
  \proofstep{2}{\{x\}\subseteq M}{\FormulaRefAuto{a\in A \vdash \{a\}\subseteq A}{6}}
  \proofstep{1,2}{F\cup\{x\}\subseteq M}{%
    \FormulaRefAuto{A\subseteq M\dsep B\subseteq M \vdash A\cup B\subseteq M}{4,8}}
  \proofstep{1,2}{\Finite{F\cup\{x\}}}{%
    \FormulaRefAuto{FiniteAdjunction}[thm]{7,5}}
  \proofstep{1,2}{F\cup\{x\}\in\FinSubsets{M}}{%
    \FormulaRefAuto{F\in\FinSubsets{M}\eqvdash F\subseteq M\land \Finite{F}}{\rAI{9,10}}}
  \proofstep{}{F\subseteq F\cup\{x\}}{\FormulaRefAuto{A\subseteq A\cup B}}
  \proofstep{}{x\in\{x\}}{\FormulaRefAuto{a\in\{a\}}}
  \proofstep{}{x\in F\cup\{x\}}{%
    \FormulaRefAuto{z \in B \vdash z \in A \cup B}{13}}
  \proofstep{2}{F\cup\{x\}\neq F}{%
    \FormulaRefAuto{x \in A,\, x \not\in B \vdash A \neq B}{14,7}}
  \proofstep{2}{F\neq F\cup\{x\}}{\FormulaRefAuto{a \neq b \vdash b \neq a}{15}}
  \proofstep{2}{F\subset F\cup\{x\}}{%
    \FormulaRefAuto{A \subset B \coloneq A \subseteq B \land A\neq B}{\rAI{12,16}}}
  \proofstep{1,2}{F\cup\{x\}\in\FinSubsets{M}\land F\subset F\cup\{x\}}{\rAI{11,17}}
\end{tabproof}

\paragraph{Maximalität}

\FormulaThmDeltaK[Tarski-Endlichkeit liefert ein Maximum]{%
  \TarskiFinite(M)\dsep B\subseteq \powerset(M)\dsep B\neq\varnothing
  \vdash \exists F\in B\,\forall G\in B\,\bigl(F\subseteq G\rightarrow G=F\bigr)
}{TarskiFiniteElim}{%
  \DeltaDecl{Mengen}{M\dsep B\dsep F\dsep G}%
}
\begin{tabproof}
  \proofstep{1}{\TarskiFinite(M)}{\rA}
  \proofstep{2}{B\subseteq \powerset(M)}{\rA}
  \proofstep{3}{B\neq\varnothing}{\rA}
  \proofstep{1}{%
    \bigl(B\subseteq \powerset(M)\land B\neq\varnothing\bigr)\rightarrow
    \exists F\in B\,\forall G\in B\,\bigl(F\subseteq G\rightarrow G=F\bigr)
  }{\FormulaRefAuto{Tarski-endliche Menge}[def]{1}}
  \proofstep{2,3}{B\subseteq \powerset(M)\land B\neq\varnothing}{\rAI{2,3}}
  \proofstep{1,2,3}{\exists F\in B\,\forall G\in B\,\bigl(F\subseteq G\rightarrow G=F\bigr)}{\rIE{5,4}}
\end{tabproof}

\FormulaThmDelta[Maximale endliche Teilmenge ist die ganze Menge]{%
  F\in\FinSubsets{M}\dsep
  \forall G\in\FinSubsets{M}\,\bigl(F\subseteq G\rightarrow G=F\bigr)
  \vdash F=M
}{%
  \DeltaDecl{Mengen}{M\dsep F\dsep G\dsep x}%
}
\begin{tabproof}
  \proofstep{1}{F\in\FinSubsets{M}}{\rA}
  \proofstep{2}{\forall G\in\FinSubsets{M}\,\bigl(F\subseteq G\rightarrow G=F\bigr)}{\rA}
  \proofstep{1}{F\subseteq M\land \Finite{F}}{%
    \FormulaRefAuto{F\in\FinSubsets{M}\eqvdash F\subseteq M\land \Finite{F}}{1}}
  \proofstep{1}{F\subseteq M}{\rAEa{3}}
  \proofstep{5}{F\neq M}{\rA}
  \proofstep{1,5}{\exists x\in M\setminus F}{%
    \FormulaRefAuto{A\subseteq B\dsep A\neq B\vdash \exists x\in B\setminus A}{4,5}}
  \proofstep{7}{x\in M\setminus F}{\rA}
  \proofstep{1,7}{F\cup\{x\}\in\FinSubsets{M}\land F\subset F\cup\{x\}}{%
    \FormulaRefAuto{F\in\FinSubsets{M}\dsep x\in M\setminus F
    \vdash F\cup\{x\}\in\FinSubsets{M}\land F\subset F\cup\{x\}}{1,7}}
  \proofstep{1,7}{F\cup\{x\}\in\FinSubsets{M}}{\rAEa{8}}
  \proofstep{1,7}{F\subset F\cup\{x\}}{\rAEb{8}}
  \proofstep{1,7}{F\subseteq F\cup\{x\}}{\FormulaRefAuto{A\subset B\vdash A\subseteq B}{10}}
  \proofstep{2,7}{F\subseteq F\cup\{x\}\rightarrow F\cup\{x\}=F}{\rUE{2}}
  \proofstep{1,2,7}{F\cup\{x\}=F}{\rIE{11,12}}
  \proofstep{1,7}{F\neq F\cup\{x\}}{\FormulaRefAuto{A\subset B \vdash A\neq B}{10}}
  \proofstep{1,7}{F\cup\{x\}\neq F}{\FormulaRefAuto{a \neq b \vdash b \neq a}{14}}
  \proofstep{1,2,7}{\bot}{\rBI{13,15}}
  \proofstep{1,2,5}{\bot}{\rEE{6,7,16}}
  \proofstep{1,2}{F=M}{\rCE{5,17}}
\end{tabproof}

\paragraph{Schluss: Tarski-Endlichkeit impliziert Endlichkeit}

\FormulaThmDelta[Tarski-Endlichkeit impliziert Endlichkeit]{%
  \TarskiFinite(M)\vdash \Finite{M}
}{%
  \DeltaDecl{Mengen}{M\dsep F\dsep G}%
}
\begin{tabproof}
  \proofstep{1}{\TarskiFinite(M)}{\rA}
  \proofstep{}{\FinSubsets{M}\subseteq \powerset(M)\land \FinSubsets{M}\neq\varnothing}{%
    \FormulaRefAuto{\FinSubsets{M}\subseteq \powerset(M)\land \FinSubsets{M}\neq\varnothing}}
  \proofstep{}{\FinSubsets{M}\subseteq \powerset(M)}{\rAEa{2}}
  \proofstep{}{\FinSubsets{M}\neq\varnothing}{\rAEb{2}}
  \proofstep{1}{%
    \exists F\in\FinSubsets{M}\,
    \forall G\in\FinSubsets{M}\,\bigl(F\subseteq G\rightarrow G=F\bigr)
  }{%
    \FormulaRefAuto{TarskiFiniteElim}[thm]{1,3,4}}
  \proofstep{6}{%
    F\in\FinSubsets{M}\land
    \forall G\in\FinSubsets{M}\,\bigl(F\subseteq G\rightarrow G=F\bigr)
  }{\rA}
  \proofstep{6}{F\in\FinSubsets{M}}{\rAEa{6}}
  \proofstep{6}{\forall G\in\FinSubsets{M}\,\bigl(F\subseteq G\rightarrow G=F\bigr)}{\rAEb{6}}
  \proofstep{6}{F=M}{%
    \FormulaRefAuto{F\in\FinSubsets{M}\dsep
    \forall G\in\FinSubsets{M}\,\bigl(F\subseteq G\rightarrow G=F\bigr)
    \vdash F=M}{7,8}}
  \proofstep{6}{F\subseteq M\land \Finite{F}}{%
    \FormulaRefAuto{F\in\FinSubsets{M}\eqvdash F\subseteq M\land \Finite{F}}{7}}
  \proofstep{6}{\Finite{F}}{\rAEb{10}}
  \proofstep{6}{\Finite{M}}{\rIE{9,11}}
  \proofstep{1}{\Finite{M}}{\rEE{5,6,12}}
\end{tabproof}

\paragraph{Bemerkung zur Auswahl}

Dieser Beweis ist auswahlfrei. Er verwendet nur das Tarski-Maximum der bereits
definierten Familie \(\FinSubsets{M}\); es wird keine unendliche Folge von
Entscheidungen konstruiert. Der direkte Schluss von Dedekind-Endlichkeit auf
Endlichkeit ist damit nicht mitbewiesen.

\FormulaThmDeltaK[Endliche Mengen sind Dedekind-endlich]{%
  \Finite{M}\vdash \neg\exists H\colon \mathbb{N}\inj M
}{FiniteNoNatInjection}{%
  \DeltaDecl{Mengen}{M\dsep n}%
  \DeltaDecl{Funktionen}{H}%
}
\begin{tabproof}
  \proofstep{1}{\Finite{M}}{\rA}
  \proofstep{1}{\exists n\,\bigl(n\in\mathbb{N}\land \NatSeg{n}\EqCard M\bigr)}{%
    \FormulaRefAuto{FiniteEqCardEquiv}[thm]{1}}
  \proofstep{3}{n\in\mathbb{N}\land \NatSeg{n}\EqCard M}{\rA}
  \proofstep{3}{n\in\mathbb{N}}{\rAEa{3}}
  \proofstep{3}{\NatSeg{n}\EqCard M}{\rAEb{3}}
  \proofstep{3}{\neg\exists H\colon \mathbb{N}\inj M}{%
    \FormulaRefAuto{FiniteWitnessNoNatInjection}[thm]{4,5}}
  \proofstep{1}{\neg\exists H\colon \mathbb{N}\inj M}{\rEE{2,3,6}}
\end{tabproof}

\FormulaThmDelta[Endliche Mengen sind Tarski-endlich]{%
  \Finite{M}\vdash \TarskiFinite(M)
}{%
  \DeltaDecl{Mengen}{M}%
  \DeltaDecl{Funktionen}{H}%
}
\begin{tabproof}
  \proofstep{1}{\Finite{M}}{\rA}
  \proofstep{1}{\neg\exists H\colon \mathbb{N}\inj M}{%
    \FormulaRefAuto{\Finite{M}\vdash \neg\exists H\colon \mathbb{N}\inj M}{1}}
  \proofstep{1}{\TarskiFinite(M)}{%
    \FormulaRefAuto{\neg\exists H\colon \mathbb{N}\inj M\vdash \TarskiFinite(M)}{2}}
\end{tabproof}

\subsubsection[Max.- und Min.-Form der Tarski-Endlichkeit]{Maximalitäts- und Minimalitätsform der Tarski-Endlichkeit}

Die Minimalitätsform ist die duale Form der Maximalitätsform. Der Übergang
geschieht über relatives Komplement in \(M\): Eine maximale Menge in der
Komplementfamilie entspricht genau einer minimalen Menge in der
Ausgangsfamilie.

Für \(X,Y\subseteq M\) und die Zuordnung \(X\mapsto M\setminus X\) fasst das
folgende Wörterbuch die Übersetzung zusammen:
\[
\begin{aligned}
X\subseteq Y
  &\quad\longleftrightarrow\quad
  M\setminus Y\subseteq M\setminus X,\\
\text{maximal in }B^{\complement_M}
  &\quad\longleftrightarrow\quad
  \text{minimal in }B,\\
\text{minimal in }B^{\complement_M}
  &\quad\longleftrightarrow\quad
  \text{maximal in }B.
\end{aligned}
\]

Wir trennen den Nachweis in drei Schritte: Zuerst wird die Komplementfamilie
typisiert, dann wird die Ordnungsumkehrung des relativen Komplements benutzt,
und zuletzt werden Maxima der Komplementfamilie in Minima der Ausgangsfamilie
übersetzt sowie umgekehrt.
Die allgemeinen Rechenregeln zum relativen Komplement stehen im Kapitel
\enquote{Die Differenz}.

\paragraph{Komplementduale Familien}

\FormulaDefDelta[Komplementduale einer Potenzmengenfamilie]{%
  B^{\complement_M}\coloneq \{M\setminus X\mid X\in B\}
}{%
  \DeltaDecl{Mengen}{M\dsep B\dsep X}%
}

\FormulaThmDeltaK[Charakterisierung der Komplementdualen]{%
  Z\in B^{\complement_M}\eqvdash
  \exists X\in B\,\bigl(Z=M\setminus X\bigr)
}{TarskiKomplementdualChar}{
  \DeltaDecl{Mengen}{M\dsep B\dsep X\dsep Z}%
}
\begin{tabproof}
  \proofstep{}{Z\in B^{\complement_M}\eqvdash
    \exists X\in B\,\bigl(Z=M\setminus X\bigr)}{%
    \FormulaRefAuto{B^{\complement_M}\coloneq \{M\setminus X\mid X\in B\}}[def]}
\end{tabproof}

\FormulaThmDeltaK[Einführung in die Komplementduale]{%
  X\in B
  \vdash
  M\setminus X\in B^{\complement_M}
}{TarskiKomplementdualIntro}{%
  \DeltaDecl{Mengen}{M\dsep B\dsep X\dsep Y}%
}
\begin{tabproof}
  \proofstep{1}{X\in B}{\rA}
  \proofstep{1}{\exists Y\in B\,\bigl(M\setminus X=M\setminus Y\bigr)}{%
    \rEI{\rAI{1,\rII}}}
  \proofstep{1}{M\setminus X\in B^{\complement_M}}{%
    \FormulaRefAuto{TarskiKomplementdualChar}[thm]{2}}
\end{tabproof}

\FormulaThmDeltaK[Element einer Potenzmengenfamilie liegt in der Grundmenge]{%
  B\subseteq \powerset(M)\dsep X\in B
  \vdash
  X\subseteq M
}{PowerFamilyMemberSubsetBase}{%
  \DeltaDecl{Mengen}{M\dsep B\dsep X}%
}
\begin{tabproof}
  \proofstep{1}{B\subseteq \powerset(M)}{\rA}
  \proofstep{2}{X\in B}{\rA}
  \proofstep{1,2}{X\in\powerset(M)}{%
    \FormulaRefAuto{A\subseteq B,\, x\in A \vdash x\in B}[thm]{1,2}}
  \proofstep{1,2}{X\subseteq M}{%
    \FormulaRefAuto{x \in \powerset(A)\;\eqvdash\; x \subseteq A}[thm]{3}}
\end{tabproof}

\FormulaThmDeltaK[Komplementduale bleibt eine nichtleere Potenzmengenfamilie]{%
  B\subseteq \powerset(M)\dsep B\neq\varnothing
  \vdash
  B^{\complement_M}\subseteq \powerset(M)\land
  B^{\complement_M}\neq\varnothing
}{TarskiKomplementdualPowerFamily}{
  \DeltaDecl{Mengen}{M\dsep B\dsep X\dsep Z}%
}
\begin{tabproof}
  \proofstep{1}{B\subseteq \powerset(M)}{\rA}
  \proofstep{2}{B\neq\varnothing}{\rA}
  \proofstep{3}{Z\in B^{\complement_M}}{\rA}
  \proofstep{3}{\exists X\in B\,\bigl(Z=M\setminus X\bigr)}{%
    \FormulaRefAuto{TarskiKomplementdualChar}[thm]{3}}
  \proofstep{5}{X\in B\land Z=M\setminus X}{\rA}
  \proofstep{5}{X\in B}{\rAEa{5}}
  \proofstep{5}{Z=M\setminus X}{\rAEb{5}}
  \proofstep{1,5}{M\setminus X\subseteq M}{%
    \FormulaRefAuto{A\setminus B\subseteq A}[thm]}
  \proofstep{1,5}{Z\subseteq M}{\rIE{7,8}}
  \proofstep{1,5}{Z\in\powerset(M)}{%
    \FormulaRefAuto{x \in \powerset(A)\;\eqvdash\; x \subseteq A}[thm]{9}}
  \proofstep{1,3}{Z\in\powerset(M)}{\rEE{4,5,10}}
  \proofstep{1}{B^{\complement_M}\subseteq \powerset(M)}{%
    \FormulaRefAuto{A \subseteq B \coloneq \forall x\,(x\in A \rightarrow x\in B)}[def]{\rUI{\rRI{3,11}}}}
  \proofstep{2}{\exists X\in B}{\FormulaRefAuto{B\neq\varnothing\vdash \exists x\in B}[thm]{2}}
  \proofstep{14}{X\in B}{\rA}
  \proofstep{14}{M\setminus X\in B^{\complement_M}}{%
    \FormulaRefAuto{TarskiKomplementdualIntro}[thm]{14}}
  \proofstep{14}{B^{\complement_M}\neq\varnothing}{%
    \FormulaRefAuto{a \in A \vdash A \neq \varnothing}[thm]{15}}
  \proofstep{2}{B^{\complement_M}\neq\varnothing}{\rEE{13,14,16}}
  \proofstep{1,2}{%
    B^{\complement_M}\subseteq \powerset(M)\land
    B^{\complement_M}\neq\varnothing}{\rAI{12,17}}
\end{tabproof}

\paragraph{Dualität von Maxima und Minima}

Die beiden Übersetzungssätze sind spiegelbildlich. Wir schreiben sie dennoch
beide aus: So bleiben die späteren Äquivalenzbeweise auf jeweils einen
benannten Übersetzungssatz reduziert.

\FormulaThmDeltaK[Punktweise Übersetzung eines komplementdualen Maximums]{%
  \begin{aligned}[t]
    &B\subseteq \powerset(M)\dsep X\in B\dsep U=M\setminus X\dsep{}\\
    &\forall V\in B^{\complement_M}\,
      \bigl(U\subseteq V\rightarrow V=U\bigr)
    \vdash{}\\
    &\forall Y\in B\,\bigl(Y\subseteq X\rightarrow Y=X\bigr)
  \end{aligned}
}{TarskiKomplementdualMaxWitness}{%
  \DeltaDecl{Mengen}{M\dsep B\dsep U\dsep V\dsep X\dsep Y}%
}
\begin{tabproof}
  \proofstep{1}{B\subseteq \powerset(M)}{\rA}
  \proofstep{2}{X\in B}{\rA}
  \proofstep{3}{U=M\setminus X}{\rA}
  \proofstep{4}{\forall V\in B^{\complement_M}\,
    \bigl(U\subseteq V\rightarrow V=U\bigr)}{\rA}
  \proofstep{5}{Y\in B}{\rA}
  \proofstep{6}{Y\subseteq X}{\rA}
  \proofstep{1,2}{X\subseteq M}{%
    \FormulaRefAuto{PowerFamilyMemberSubsetBase}[thm]{1,2}}
  \proofstep{1,5}{Y\subseteq M}{%
    \FormulaRefAuto{PowerFamilyMemberSubsetBase}[thm]{1,5}}
  \proofstep{5}{M\setminus Y\in B^{\complement_M}}{%
    \FormulaRefAuto{TarskiKomplementdualIntro}[thm]{5}}
  \proofstep{1,2,5,6}{M\setminus X\subseteq M\setminus Y}{%
    \FormulaRefAuto{RelCompAntitone}[thm]{8,7,6}}
  \proofstep{1,2,3,5,6}{U\subseteq M\setminus Y}{\rIE{3,10}}
  \proofstep{4,5,6}{U\subseteq M\setminus Y\rightarrow M\setminus Y=U}{\rUE{4}}
  \proofstep{1,2,3,4,5,6}{M\setminus Y=U}{\rIE{11,12}}
  \proofstep{1,2,3,4,5,6}{M\setminus Y=M\setminus X}{\rIE{3,13}}
  \proofstep{1,2,3,4,5,6}{Y=X}{%
    \FormulaRefAuto{RelCompEqualityCancel}[thm]{7,8,14}}
  \proofstep{1,2,3,4,5}{Y\subseteq X\rightarrow Y=X}{\rRI{6,15}}
  \proofstep{1,2,3,4}{\forall Y\in B\,\bigl(Y\subseteq X\rightarrow Y=X\bigr)}{%
    \rUI{\rRI{5,16}}}
\end{tabproof}

\FormulaThmDeltaK[Punktweise Übersetzung eines komplementdualen Minimums]{%
  \begin{aligned}[t]
    &B\subseteq \powerset(M)\dsep X\in B\dsep U=M\setminus X\dsep{}\\
    &\forall V\in B^{\complement_M}\,
      \bigl(V\subseteq U\rightarrow V=U\bigr)
    \vdash{}\\
    &\forall Y\in B\,\bigl(X\subseteq Y\rightarrow Y=X\bigr)
  \end{aligned}
}{TarskiKomplementdualMinWitness}{%
  \DeltaDecl{Mengen}{M\dsep B\dsep U\dsep V\dsep X\dsep Y}%
}
\begin{tabproof}
  \proofstep{1}{B\subseteq \powerset(M)}{\rA}
  \proofstep{2}{X\in B}{\rA}
  \proofstep{3}{U=M\setminus X}{\rA}
  \proofstep{4}{\forall V\in B^{\complement_M}\,
    \bigl(V\subseteq U\rightarrow V=U\bigr)}{\rA}
  \proofstep{5}{Y\in B}{\rA}
  \proofstep{6}{X\subseteq Y}{\rA}
  \proofstep{1,2}{X\subseteq M}{%
    \FormulaRefAuto{PowerFamilyMemberSubsetBase}[thm]{1,2}}
  \proofstep{1,5}{Y\subseteq M}{%
    \FormulaRefAuto{PowerFamilyMemberSubsetBase}[thm]{1,5}}
  \proofstep{5}{M\setminus Y\in B^{\complement_M}}{%
    \FormulaRefAuto{TarskiKomplementdualIntro}[thm]{5}}
  \proofstep{1,2,5,6}{M\setminus Y\subseteq M\setminus X}{%
    \FormulaRefAuto{RelCompAntitone}[thm]{7,8,6}}
  \proofstep{1,2,3,5,6}{M\setminus Y\subseteq U}{\rIE{3,10}}
  \proofstep{4,5,6}{M\setminus Y\subseteq U\rightarrow M\setminus Y=U}{\rUE{4}}
  \proofstep{1,2,3,4,5,6}{M\setminus Y=U}{\rIE{11,12}}
  \proofstep{1,2,3,4,5,6}{M\setminus Y=M\setminus X}{\rIE{3,13}}
  \proofstep{1,2,3,4,5,6}{Y=X}{%
    \FormulaRefAuto{RelCompEqualityCancel}[thm]{7,8,14}}
  \proofstep{1,2,3,4,5}{X\subseteq Y\rightarrow Y=X}{\rRI{6,15}}
  \proofstep{1,2,3,4}{\forall Y\in B\,\bigl(X\subseteq Y\rightarrow Y=X\bigr)}{%
    \rUI{\rRI{5,16}}}
\end{tabproof}

\FormulaThmDeltaK[Komplementduales Maximum liefert ein Minimum]{%
  \begin{aligned}[t]
    &B\subseteq \powerset(M)\dsep
    \exists U\in B^{\complement_M}\,
      \forall V\in B^{\complement_M}\,
        \bigl(U\subseteq V\rightarrow V=U\bigr)
    \vdash{}\\
    &\exists X\in B\,\forall Y\in B\,
      \bigl(Y\subseteq X\rightarrow Y=X\bigr)
  \end{aligned}
}{TarskiKomplementdualMaxToMin}{%
  \DeltaDecl{Mengen}{M\dsep B\dsep U\dsep V\dsep X\dsep Y}%
}
\begin{tabproof}
  \proofstep{1}{B\subseteq \powerset(M)}{\rA}
  \proofstep{2}{%
    \begin{aligned}[t]
      &\exists U\in B^{\complement_M}\,
        \forall V\in B^{\complement_M}\,
          \bigl(U\subseteq V\rightarrow V=U\bigr)
    \end{aligned}}{\rA}
  \proofstep{3}{%
    \begin{aligned}[t]
      &U\in B^{\complement_M}\land{}\\
      &\forall V\in B^{\complement_M}\,
        \bigl(U\subseteq V\rightarrow V=U\bigr)
    \end{aligned}}{\rA}
  \proofstep{3}{U\in B^{\complement_M}}{\rAEa{3}}
  \proofstep{3}{%
    \begin{aligned}[t]
      &\forall V\in B^{\complement_M}\,
        \bigl(U\subseteq V\rightarrow V=U\bigr)
    \end{aligned}}{\rAEb{3}}
  \proofstep{3}{\exists X\in B\,\bigl(U=M\setminus X\bigr)}{%
    \FormulaRefAuto{TarskiKomplementdualChar}[thm]{4}}
  \proofstep{7}{X\in B\land U=M\setminus X}{\rA}
  \proofstep{7}{X\in B}{\rAEa{7}}
  \proofstep{7}{U=M\setminus X}{\rAEb{7}}
  \proofstep{1,3,7}{%
    \begin{aligned}[t]
      &\forall Y\in B\,\bigl(Y\subseteq X\rightarrow Y=X\bigr)
    \end{aligned}}{%
      \FormulaRefAuto{TarskiKomplementdualMaxWitness}[thm]{1,8,9,5}}
  \proofstep{1,3,7}{%
    \begin{aligned}[t]
      &\exists X\in B\,\forall Y\in B\,
        \bigl(Y\subseteq X\rightarrow Y=X\bigr)
    \end{aligned}}{\rEI{\rAI{8,10}}}
  \proofstep{1,3}{%
    \begin{aligned}[t]
      &\exists X\in B\,\forall Y\in B\,
        \bigl(Y\subseteq X\rightarrow Y=X\bigr)
    \end{aligned}}{\rEE{6,7,11}}
  \proofstep{1,2}{%
    \begin{aligned}[t]
      &\exists X\in B\,\forall Y\in B\,
        \bigl(Y\subseteq X\rightarrow Y=X\bigr)
    \end{aligned}}{\rEE{2,3,12}}
\end{tabproof}

\FormulaThmDeltaK[Komplementduales Minimum liefert ein Maximum]{%
  \begin{aligned}[t]
    &B\subseteq \powerset(M)\dsep
    \exists U\in B^{\complement_M}\,
      \forall V\in B^{\complement_M}\,
        \bigl(V\subseteq U\rightarrow V=U\bigr)
    \vdash{}\\
    &\exists X\in B\,\forall Y\in B\,
      \bigl(X\subseteq Y\rightarrow Y=X\bigr)
  \end{aligned}
}{TarskiKomplementdualMinToMax}{%
  \DeltaDecl{Mengen}{M\dsep B\dsep U\dsep V\dsep X\dsep Y}%
}
\begin{tabproof}
  \proofstep{1}{B\subseteq \powerset(M)}{\rA}
  \proofstep{2}{%
    \begin{aligned}[t]
      &\exists U\in B^{\complement_M}\,
        \forall V\in B^{\complement_M}\,
          \bigl(V\subseteq U\rightarrow V=U\bigr)
    \end{aligned}}{\rA}
  \proofstep{3}{%
    \begin{aligned}[t]
      &U\in B^{\complement_M}\land{}\\
      &\forall V\in B^{\complement_M}\,
        \bigl(V\subseteq U\rightarrow V=U\bigr)
    \end{aligned}}{\rA}
  \proofstep{3}{U\in B^{\complement_M}}{\rAEa{3}}
  \proofstep{3}{%
    \begin{aligned}[t]
      &\forall V\in B^{\complement_M}\,
        \bigl(V\subseteq U\rightarrow V=U\bigr)
    \end{aligned}}{\rAEb{3}}
  \proofstep{3}{\exists X\in B\,\bigl(U=M\setminus X\bigr)}{%
    \FormulaRefAuto{TarskiKomplementdualChar}[thm]{4}}
  \proofstep{7}{X\in B\land U=M\setminus X}{\rA}
  \proofstep{7}{X\in B}{\rAEa{7}}
  \proofstep{7}{U=M\setminus X}{\rAEb{7}}
  \proofstep{1,3,7}{%
    \begin{aligned}[t]
      &\forall Y\in B\,\bigl(X\subseteq Y\rightarrow Y=X\bigr)
    \end{aligned}}{%
      \FormulaRefAuto{TarskiKomplementdualMinWitness}[thm]{1,8,9,5}}
  \proofstep{1,3,7}{%
    \begin{aligned}[t]
      &\exists X\in B\,\forall Y\in B\,
        \bigl(X\subseteq Y\rightarrow Y=X\bigr)
    \end{aligned}}{\rEI{\rAI{8,10}}}
  \proofstep{1,3}{%
    \begin{aligned}[t]
      &\exists X\in B\,\forall Y\in B\,
        \bigl(X\subseteq Y\rightarrow Y=X\bigr)
    \end{aligned}}{\rEE{6,7,11}}
  \proofstep{1,2}{%
    \begin{aligned}[t]
      &\exists X\in B\,\forall Y\in B\,
        \bigl(X\subseteq Y\rightarrow Y=X\bigr)
    \end{aligned}}{\rEE{2,3,12}}
\end{tabproof}

\paragraph{Äquivalenz der Prädikate}

Nach der komplementdualen Übersetzung bleiben die Prädikatbeweise kurz: Man
wendet die jeweilige Tarski-Form auf die Komplementfamilie an und übersetzt
das gefundene Extremum zurück.

\FormulaThmDeltaK[Tarski-Minimalendlichkeit liefert ein Minimum]{%
  \begin{aligned}[t]
    &\TarskiMinFinite(M)\dsep B\subseteq \powerset(M)\dsep B\neq\varnothing
    \vdash{}\\
    &\exists F\in B\,\forall G\in B\,\bigl(G\subseteq F\rightarrow G=F\bigr)
  \end{aligned}
}{TarskiMinFiniteElim}{%
  \DeltaDecl{Mengen}{M\dsep B\dsep F\dsep G}%
}
\begin{tabproof}
  \proofstep{1}{\TarskiMinFinite(M)}{\rA}
  \proofstep{2}{B\subseteq \powerset(M)}{\rA}
  \proofstep{3}{B\neq\varnothing}{\rA}
  \proofstep{1}{%
    \begin{aligned}[t]
      &\bigl(B\subseteq \powerset(M)\land B\neq\varnothing\bigr)\rightarrow{}\\
      &\exists F\in B\,\forall G\in B\,\bigl(G\subseteq F\rightarrow G=F\bigr)
    \end{aligned}}{\FormulaRefAuto{minimal Tarski-endliche Menge}[def]{1}}
  \proofstep{2,3}{B\subseteq \powerset(M)\land B\neq\varnothing}{\rAI{2,3}}
  \proofstep{1,2,3}{%
    \begin{aligned}[t]
      &\exists F\in B\,\forall G\in B\,
        \bigl(G\subseteq F\rightarrow G=F\bigr)
    \end{aligned}}{\rIE{5,4}}
\end{tabproof}

\FormulaThmDeltaK[Tarski-Endlichkeit impliziert minimale Tarski-Endlichkeit]{%
  \TarskiFinite(M)\vdash \TarskiMinFinite(M)
}{TarskiFiniteToTarskiMinFinite}{%
  \DeltaDecl{Mengen}{M\dsep B\dsep F\dsep G}%
}
\begin{tabproof}
  \proofstep{1}{\TarskiFinite(M)}{\rA}
  \proofstep{2}{B\subseteq \powerset(M)\land B\neq\varnothing}{\rA}
  \proofstep{2}{B\subseteq \powerset(M)}{\rAEa{2}}
  \proofstep{2}{B\neq\varnothing}{\rAEb{2}}
  \proofstep{2}{%
    B^{\complement_M}\subseteq \powerset(M)\land
    B^{\complement_M}\neq\varnothing}{%
    \FormulaRefAuto{TarskiKomplementdualPowerFamily}[thm]{3,4}}
  \proofstep{2}{B^{\complement_M}\subseteq \powerset(M)}{\rAEa{5}}
  \proofstep{2}{B^{\complement_M}\neq\varnothing}{\rAEb{5}}
  \proofstep{1,2}{%
    \begin{aligned}[t]
      &\exists F\in B^{\complement_M}\,
      \forall G\in B^{\complement_M}\,
      \bigl(F\subseteq G\rightarrow G=F\bigr)
    \end{aligned}}{\FormulaRefAuto{TarskiFiniteElim}[thm]{1,6,7}}
  \proofstep{1,2}{%
    \begin{aligned}[t]
      &\exists F\in B\,\forall G\in B\,
      \bigl(G\subseteq F\rightarrow G=F\bigr)
    \end{aligned}}{\FormulaRefAuto{TarskiKomplementdualMaxToMin}[thm]{3,8}}
  \proofstep{1}{%
    \begin{aligned}[t]
      &\bigl(B\subseteq \powerset(M)\land B\neq\varnothing\bigr)\rightarrow{}\\
      &\exists F\in B\,\forall G\in B\,\bigl(G\subseteq F\rightarrow G=F\bigr)
    \end{aligned}}{\rRI{2,9}}
  \proofstep{1}{\TarskiMinFinite(M)}{%
    \FormulaRefAuto{minimal Tarski-endliche Menge}[def]{10}}
\end{tabproof}

\FormulaThmDeltaK[Minimale Tarski-Endlichkeit impliziert Tarski-Endlichkeit]{%
  \TarskiMinFinite(M)\vdash \TarskiFinite(M)
}{TarskiMinFiniteToTarskiFinite}{%
  \DeltaDecl{Mengen}{M\dsep B\dsep F\dsep G}%
}
\begin{tabproof}
  \proofstep{1}{\TarskiMinFinite(M)}{\rA}
  \proofstep{2}{B\subseteq \powerset(M)\land B\neq\varnothing}{\rA}
  \proofstep{2}{B\subseteq \powerset(M)}{\rAEa{2}}
  \proofstep{2}{B\neq\varnothing}{\rAEb{2}}
  \proofstep{2}{%
    B^{\complement_M}\subseteq \powerset(M)\land
    B^{\complement_M}\neq\varnothing}{%
    \FormulaRefAuto{TarskiKomplementdualPowerFamily}[thm]{3,4}}
  \proofstep{2}{B^{\complement_M}\subseteq \powerset(M)}{\rAEa{5}}
  \proofstep{2}{B^{\complement_M}\neq\varnothing}{\rAEb{5}}
  \proofstep{1,2}{%
    \begin{aligned}[t]
      &\exists F\in B^{\complement_M}\,
      \forall G\in B^{\complement_M}\,
      \bigl(G\subseteq F\rightarrow G=F\bigr)
    \end{aligned}}{\FormulaRefAuto{TarskiMinFiniteElim}[thm]{1,6,7}}
  \proofstep{1,2}{%
    \begin{aligned}[t]
      &\exists F\in B\,\forall G\in B\,
      \bigl(F\subseteq G\rightarrow G=F\bigr)
    \end{aligned}}{\FormulaRefAuto{TarskiKomplementdualMinToMax}[thm]{3,8}}
  \proofstep{1}{%
    \begin{aligned}[t]
      &\bigl(B\subseteq \powerset(M)\land B\neq\varnothing\bigr)\rightarrow{}\\
      &\exists F\in B\,\forall G\in B\,\bigl(F\subseteq G\rightarrow G=F\bigr)
    \end{aligned}}{\rRI{2,9}}
  \proofstep{1}{\TarskiFinite(M)}{%
    \FormulaRefAuto{Tarski-endliche Menge}[def]{10}}
\end{tabproof}

\FormulaThmDeltaK[Äquivalenz der maximalen und minimalen Tarski-Endlichkeit]{%
  \TarskiFinite(M)\eqvdash \TarskiMinFinite(M)
}{TarskiMaxMinFiniteEquiv}{%
  \DeltaDecl{Mengen}{M}%
}
\begin{tabproofsplit}
  \proofpart{\(\vdash\)}
    \proofstep{1}{\TarskiFinite(M)}{\rA}
    \proofstep{1}{\TarskiMinFinite(M)}{%
      \FormulaRefAuto{TarskiFiniteToTarskiMinFinite}[thm]{1}}
  \closeproofpart

  \proofpart{\(\dashv\)}
    \proofstep{1}{\TarskiMinFinite(M)}{\rA}
    \proofstep{1}{\TarskiFinite(M)}{%
      \FormulaRefAuto{TarskiMinFiniteToTarskiFinite}[thm]{1}}
  \closeproofpart
\end{tabproofsplit}

\subsection{Zusammenfassung der Äquivalenzen}

\FormulaThmDeltaR[Äquivalenz von Endlichkeit und Dedekind-Endlichkeit]{%
  \Finite{M}\eqvdash \neg\exists H\colon \mathbb{N}\inj M
}{%
  \Finite{M}\eqvdash \neg\exists H\colon \mathbb{N}\inj M
}{%
  \DeltaDecl{Mengen}{M}%
  \DeltaDecl{Funktionen}{H}%
}
\begin{tabproofsplit}
  \proofpart{\(\vdash\)}
    \proofstep{1}{\Finite{M}}{\rA}
    \proofstep{1}{\neg\exists H\colon \mathbb{N}\inj M}{%
      \FormulaRefAuto{\Finite{M}\vdash \neg\exists H\colon \mathbb{N}\inj M}{1}}
  \closeproofpart

  \proofpart{\(\dashv\)}
    \proofstep{1}{\neg\exists H\colon \mathbb{N}\inj M}{\rA}
    \proofstep{1}{\TarskiFinite(M)}{%
      \FormulaRefAuto{\neg\exists H\colon \mathbb{N}\inj M\vdash \TarskiFinite(M)}{1},
      \FormulaRefAuto{Auswahlaxiom}[ax]}
    \proofstep{1}{\Finite{M}}{%
      \FormulaRefAuto{\TarskiFinite(M)\vdash \Finite{M}}{2}}
  \closeproofpart
\end{tabproofsplit}

\FormulaThmDeltaR[Äquivalenz von Endlichkeit und Ausschluss von Surjektionen nach \(\mathbb{N}\)]{%
  \Finite{M}\eqvdash \neg\exists H\colon M\sur\mathbb{N}
}{%
  \Finite{M}\eqvdash \neg\exists H\colon M\sur\mathbb{N}
}{%
  \DeltaDecl{Mengen}{M}%
  \DeltaDecl{Funktionen}{G\dsep H}%
}
\begin{tabproofsplit}
  \proofpart{\(\vdash\)}
    \proofstep{1}{\Finite{M}}{\rA}
    \proofstep{1}{\neg\exists G\colon \mathbb{N}\inj M}{%
      \FormulaRefAuto{\Finite{M}\vdash \neg\exists H\colon \mathbb{N}\inj M}{1}}
    \proofstep{1}{\neg\exists H\colon M\sur \mathbb{N}}{%
      \FormulaRefAuto{%
        \neg\exists F\colon \mathbb{N}\inj M
        \vdash
        \neg\exists F\colon M\sur \mathbb{N}
      }{2}}
  \closeproofpart

  \proofpart{\(\dashv\)}
    \proofstep{1}{\neg\exists H\colon M\sur\mathbb{N}}{\rA}
    \proofstep{2}{\exists G\colon \mathbb{N}\inj M}{\rA}
    \proofstep{2}{\exists H\colon M\sur\mathbb{N}}{%
      \FormulaRefAuto{%
        \exists G\colon \mathbb{N}\inj M\vdash
        \exists H\colon M\sur\mathbb{N}
      }{2}}
    \proofstep{1,2}{\bot}{\rBI{1,3}}
    \proofstep{1}{\neg\exists G\colon \mathbb{N}\inj M}{\rCI{2,4}}
    \proofstep{1}{\TarskiFinite(M)}{%
      \FormulaRefAuto{\neg\exists H\colon \mathbb{N}\inj M\vdash \TarskiFinite(M)}{5}}
    \proofstep{1}{\Finite{M}}{%
      \FormulaRefAuto{\TarskiFinite(M)\vdash \Finite{M}}{6}}
  \closeproofpart
\end{tabproofsplit}

\FormulaThmDeltaR[Äquivalenz von Endlichkeit und Bijektivität aller Selbstinjektionen]{%
  \Finite{M}\eqvdash
  \forall F\colon M\inj M\,\bigl(F\colon M\bij M\bigr)
}{%
  \Finite{M}\eqvdash
  \forall F\colon M\inj M\,\bigl(F\colon M\bij M\bigr)
}{%
  \DeltaDecl{Mengen}{M}%
  \DeltaDecl{Funktionen}{F\dsep H}%
}
\begin{tabproofsplit}
  \proofpart{\(\vdash\)}
    \proofstep{1}{\Finite{M}}{\rA}
    \proofstep{1}{\neg\exists H\colon \mathbb{N}\inj M}{%
      \FormulaRefAuto{\Finite{M}\vdash \neg\exists H\colon \mathbb{N}\inj M}{1}}
    \proofstep{1}{\neg\exists H\colon M\sur \mathbb{N}}{%
      \FormulaRefAuto{%
        \neg\exists F\colon \mathbb{N}\inj M
        \vdash
        \neg\exists F\colon M\sur \mathbb{N}
      }{2}}
    \proofstep{1}{\forall F\colon M\inj M\,\bigl(F\colon M\bij M\bigr)}{%
      \FormulaRefAuto{%
        \neg\exists H\colon M\sur\mathbb{N}\vdash
        \forall F\colon M\inj M\,\bigl(F\colon M\bij M\bigr)
      }{3}}
  \closeproofpart

  \proofpart{\(\dashv\)}
    \proofstep{1}{\forall F\colon M\inj M\,\bigl(F\colon M\bij M\bigr)}{\rA}
    \proofstep{1}{\neg\exists H\colon \mathbb{N}\inj M}{%
      \FormulaRefAuto{%
        \forall F\colon M\inj M\,\bigl(F\colon M\bij M\bigr)\vdash
        \neg\exists H\colon \mathbb{N}\inj M
      }{1}}
    \proofstep{1}{\TarskiFinite(M)}{%
      \FormulaRefAuto{\neg\exists H\colon \mathbb{N}\inj M\vdash \TarskiFinite(M)}{2}}
    \proofstep{1}{\Finite{M}}{%
      \FormulaRefAuto{\TarskiFinite(M)\vdash \Finite{M}}{3}}
  \closeproofpart
\end{tabproofsplit}

\FormulaThmDeltaR[Äquivalenz von Endlichkeit und Bijektivität aller Selbstsurjektionen]{%
  \Finite{M}\eqvdash
  \forall F\colon M\sur M\,\bigl(F\colon M\bij M\bigr)
}{%
  \Finite{M}\eqvdash
  \forall F\colon M\sur M\,\bigl(F\colon M\bij M\bigr)
}{%
  \DeltaDecl{Mengen}{M}%
  \DeltaDecl{Funktionen}{F\dsep H}%
}
\begin{tabproofsplit}
  \proofpart{\(\vdash\)}
    \proofstep{1}{\Finite{M}}{\rA}
    \proofstep{1}{\forall F\colon M\inj M\,\bigl(F\colon M\bij M\bigr)}{%
      \FormulaRefAuto{%
        \Finite{M}\eqvdash
        \forall F\colon M\inj M\,\bigl(F\colon M\bij M\bigr)
      }{1}}
    \proofstep{1}{\forall F\colon M\sur M\,\bigl(F\colon M\bij M\bigr)}{%
      \FormulaRefAuto{%
        \forall F\colon M\inj M\,\bigl(F\colon M\bij M\bigr)\vdash
        \forall F\colon M\sur M\,\bigl(F\colon M\bij M\bigr)
      }{2}}
  \closeproofpart

  \proofpart{\(\dashv\)}
    \proofstep{1}{\forall F\colon M\sur M\,\bigl(F\colon M\bij M\bigr)}{\rA}
    \proofstep{1}{\forall F\colon M\inj M\,\bigl(F\colon M\bij M\bigr)}{%
      \FormulaRefAuto{%
        \forall F\colon M\sur M\,\bigl(F\colon M\bij M\bigr)\vdash
        \forall F\colon M\inj M\,\bigl(F\colon M\bij M\bigr)
      }{1}}
    \proofstep{1}{\Finite{M}}{%
      \FormulaRefAuto{%
        \Finite{M}\eqvdash
        \forall F\colon M\inj M\,\bigl(F\colon M\bij M\bigr)
      }{2}}
  \closeproofpart
\end{tabproofsplit}

\FormulaThmDeltaR[Äquivalenz von Endlichkeit und Tarski-Endlichkeit]{%
  \Finite{M}\eqvdash \TarskiFinite(M)
}{%
  \Finite{M}\eqvdash \TarskiFinite(M)
}{%
  \DeltaDecl{Mengen}{M}%
}
\begin{tabproofsplit}
  \proofpart{\(\vdash\)}
    \proofstep{1}{\Finite{M}}{\rA}
    \proofstep{1}{\TarskiFinite(M)}{%
      \FormulaRefAuto{\Finite{M}\vdash \TarskiFinite(M)}{1}}
  \closeproofpart

  \proofpart{\(\dashv\)}
    \proofstep{1}{\TarskiFinite(M)}{\rA}
    \proofstep{1}{\Finite{M}}{%
      \FormulaRefAuto{\TarskiFinite(M)\vdash \Finite{M}}{1}}
  \closeproofpart
\end{tabproofsplit}

\FormulaThmDeltaR[Äquivalenz von Endlichkeit und minimaler Tarski-Endlichkeit]{%
  \Finite{M}\eqvdash \TarskiMinFinite(M)
}{%
  \Finite{M}\eqvdash \TarskiMinFinite(M)
}{%
  \DeltaDecl{Mengen}{M}%
}
\begin{tabproofsplit}
  \proofpart{\(\vdash\)}
    \proofstep{1}{\Finite{M}}{\rA}
    \proofstep{1}{\TarskiFinite(M)}{%
      \FormulaRefAuto{\Finite{M}\vdash \TarskiFinite(M)}{1}}
    \proofstep{1}{\TarskiMinFinite(M)}{%
      \FormulaRefAuto{TarskiFiniteToTarskiMinFinite}[thm]{2}}
  \closeproofpart

  \proofpart{\(\dashv\)}
    \proofstep{1}{\TarskiMinFinite(M)}{\rA}
    \proofstep{1}{\TarskiFinite(M)}{%
      \FormulaRefAuto{TarskiMinFiniteToTarskiFinite}[thm]{1}}
    \proofstep{1}{\Finite{M}}{%
      \FormulaRefAuto{\TarskiFinite(M)\vdash \Finite{M}}{2}}
  \closeproofpart
\end{tabproofsplit}



\section{Wohldefiniertheit der endlichen Kardinalzahl}

\subsection{Eindeutigkeit der endlichen Kardinalzahl}

\FormulaThmDeltaK[Selbstinjektionen von Peano-Anfangsabschnitten sind bijektiv]{%
  n\in\mathbb{N}\dsep F\colon \NatSeg{n}\inj \NatSeg{n}
  \vdash F\colon \NatSeg{n}\bij \NatSeg{n}
}{PeanoNatSegSelfInjectionBijective}{%
  \DeltaRow{Mengen}{n}
  \DeltaRow{Funktionen}{F\dsep G}
}
\begin{tabproof}
  \proofstep{1}{n\in\mathbb{N}}{\rA}
  \proofstep{2}{F\colon \NatSeg{n}\inj \NatSeg{n}}{\rA}
  \proofstep{1}{\Finite{\NatSeg{n}}}{%
    \FormulaRefAuto{PeanoNatSegFinite}[thm]{1}}
  \proofstep{1}{\forall G\colon \NatSeg{n}\inj \NatSeg{n}\,\bigl(G\colon \NatSeg{n}\bij \NatSeg{n}\bigr)}{%
    \FormulaRefAuto{%
      \Finite{M}\eqvdash
      \forall F\colon M\inj M\,\bigl(F\colon M\bij M\bigr)
    }{3}}
  \proofstep{1,2}{F\colon \NatSeg{n}\bij \NatSeg{n}}{\rUE{4}}
\end{tabproof}

\subsubsection{Schubfachkern}

\paragraph{Allgemeine Hilfssätze}

\FormulaThmDelta[Surjektive Selbstabbildung vermeidet kein Element]{%
  b\in B\dsep b\notin A\dsep
  G\colon B\sur B\dsep
  \forall x\in B\,G(x)\in A
  \vdash \bot
}{%
  \DeltaRow{Mengen}{A\dsep B\dsep b\dsep x}
  \DeltaRow{Funktionen}{G}
}
\begin{tabproof}
  \proofstep{1}{b\in B}{\rA}
  \proofstep{2}{b\notin A}{\rA}
  \proofstep{3}{G\colon B\sur B}{\rA}
  \proofstep{4}{\forall x\in B\,G(x)\in A}{\rA}

  \proofstep{1,3}{\exists x\in B\,G(x)=b}{%
    \FormulaRefAuto{Surjektivität}{3,1}}
  \proofstep{6}{x\in B\land G(x)=b}{\rA}
  \proofstep{6}{x\in B}{\rAEa{6}}
  \proofstep{6}{G(x)=b}{\rAEb{6}}
  \proofstep{4,6}{G(x)\in A}{\rRE{4,7}}
  \proofstep{4,6}{b\in A}{\rIE{8,9}}
  \proofstep{2,4,6}{\bot}{\rBI{10,2}}

  \proofstep{1,2,3,4}{\bot}{\rEE{5,6,11}}
\end{tabproof}

\FormulaThmDelta[Inklusionskomposition in echte Teilmenge ist nicht surjektiv]{%
  A\subseteq B\dsep b\in B\dsep b\notin A\dsep
  F\colon B\to A\dsep
  \iota_{A,B}\circ F\colon B\sur B
  \vdash \bot
}{%
  \DeltaRow{Mengen}{A\dsep B\dsep b}
  \DeltaRow{Funktionen}{F}
}
\begin{tabproof}
  \proofstep{1}{A\subseteq B}{\rA}
  \proofstep{2}{b\in B}{\rA}
  \proofstep{3}{b\notin A}{\rA}
  \proofstep{4}{F\colon B\to A}{\rA}
  \proofstep{5}{\iota_{A,B}\circ F\colon B\sur B}{\rA}
  \proofstep{1,4}{\forall x\in B\,(\iota_{A,B}\circ F)(x)\in A}{%
    \FormulaRefAuto{%
      A\subseteq B\dsep F\colon C\to A
      \vdash
      \forall x\in C\,(\iota_{A,B}\circ F)(x)\in A
    }{1,4}}
  \proofstep{1,2,3,4,5}{\bot}{%
    \FormulaRefAuto{%
      b\in B\dsep b\notin A\dsep
      G\colon B\sur B\dsep
      \forall x\in B\,G(x)\in A
      \vdash \bot
    }{2,3,5,6}}
\end{tabproof}

\paragraph{Anwendung auf natürliche Zahlen}

\FormulaThmDelta[Keine Injektion in kleinere natürliche Zahl]{%
  m\in\mathbb{N}\dsep n\in\mathbb{N}\dsep m < n
  \dsep F\colon \NatSeg{n}\inj \NatSeg{m} \vdash \bot
}{%
  \DeltaRow{Mengen}{m\dsep n}
  \DeltaRow{Funktionen}{F}
  \DeltaRow{Relationen}{<_{\mathbb{N}}}
}
\begin{tabproof}
  \proofstep{1}{m\in\mathbb{N}}{\rA}
  \proofstep{2}{n\in\mathbb{N}}{\rA}
  \proofstep{3}{m < n}{\rA}
  \proofstep{4}{F\colon \NatSeg{n}\inj \NatSeg{m}}{\rA}

  \proofstep{1,2,3}{m\in \NatSeg{n}}{%
    \FormulaRefAuto{PeanoNatSegStrictMembership}[thm]{1,2,3}}
  \proofstep{1,2,3}{\NatSeg{m}\subseteq \NatSeg{n}}{%
    \FormulaRefAuto{PeanoNatSegMembershipSubset}[thm]{1,2,5}}
  \proofstep{1}{m\notin \NatSeg{m}}{%
    \FormulaRefAuto{PeanoNatSegEndpointNotIn}[thm]{1}}
  \proofstep{1,2,3}{\iota_{\NatSeg{m},\NatSeg{n}}\colon \NatSeg{m}\inj \NatSeg{n}}{%
    \FormulaRefAuto{A\subseteq B\vdash \iota_{A,B}\colon A\inj B}{6}}
  \proofstep{1,2,3,4}{\iota_{\NatSeg{m},\NatSeg{n}}\circ F\colon \NatSeg{n}\inj \NatSeg{n}}{%
    \FormulaRefAuto{F\colon A\inj B \dsep G\colon B\inj C\vdash G\circ F\colon A\inj C}{4,8}}

  \proofstep{1,2,3,4}{\iota_{\NatSeg{m},\NatSeg{n}}\circ F\colon \NatSeg{n}\bij \NatSeg{n}}{%
    \FormulaRefAuto{PeanoNatSegSelfInjectionBijective}[thm]{2,9}}
  \proofstep{1,2,3,4}{\iota_{\NatSeg{m},\NatSeg{n}}\circ F\colon \NatSeg{n}\sur \NatSeg{n}}{%
    \FormulaRefAuto{F\colon A\bij B \vdash F\colon A\sur B}{10}}
  \proofstep{4}{F\colon \NatSeg{n}\to \NatSeg{m}}{%
    \FormulaRefAuto{Injektive Funktion}{4}}
  \proofstep{1,2,3,4}{\bot}{%
    \FormulaRefAuto{%
      A\subseteq B\dsep b\in B\dsep b\notin A\dsep
      F\colon B\to A\dsep
      \iota_{A,B}\circ F\colon B\sur B
      \vdash \bot
    }{6,5,7,12,11}}
\end{tabproof}

\subsubsection{Ausschluss gleicher Mächtigkeit}

\FormulaThmDeltaR[Keine größere natürliche Zahl ist gleichmächtig]{%
  \begin{aligned}[t]
    &m\in\mathbb{N}\dsep n\in\mathbb{N}\dsep m < n \vdash {}\\
    &\neg(\NatSeg{m} \EqCard \NatSeg{n})
  \end{aligned}
}{%
  m\in\mathbb{N}\dsep n\in\mathbb{N}\dsep m < n
  \vdash \neg(\NatSeg{m} \EqCard \NatSeg{n})
}{%
  \DeltaRow{Mengen}{m\dsep n}
  \DeltaRow{Funktionen}{F}
  \DeltaRow{Relationen}{<_{\mathbb{N}}}
}
\begin{tabproof}
  \proofstep{1}{m\in\mathbb{N}}{\rA}
  \proofstep{2}{n\in\mathbb{N}}{\rA}
  \proofstep{3}{m < n}{\rA}
  \proofstep{4}{\NatSeg{m} \EqCard \NatSeg{n}}{\rA}
  \proofstep{4}{\NatSeg{n} \EqCard \NatSeg{m}}{\FormulaRefAuto{A \EqCard B \vdash B \EqCard A}{4}}
  \proofstep{4}{\exists F\colon \NatSeg{n}\bij \NatSeg{m}}{%
    \FormulaRefAuto{A \EqCard B \coloneqq \exists F\colon A\bij B}{5}}

  \proofstep{7}{F\colon \NatSeg{n}\bij \NatSeg{m}}{\rA}
  \proofstep{7}{F\colon \NatSeg{n}\inj \NatSeg{m}}{\FormulaRefAuto{F\colon A\bij B \vdash F\colon A\inj B}{7}}
  \proofstep{1,2,3,7}{\bot}{%
    \FormulaRefAuto{m\in\mathbb{N}\dsep n\in\mathbb{N}\dsep m < n\dsep F\colon \NatSeg{n}\inj \NatSeg{m} \vdash \bot}{1,2,3,8}}

  \proofstep{1,2,3,4}{\bot}{\rEE{6,7,9}}
  \proofstep{1,2,3}{\neg(\NatSeg{m} \EqCard \NatSeg{n})}{\rCI{4,10}}
\end{tabproof}

\FormulaThmDeltaR[Gleichmächtige Peano-Anfangsabschnitte haben gleiche Länge]{%
  n\in\mathbb{N}\dsep m\in\mathbb{N}\dsep
  \NatSeg{n} \EqCard \NatSeg{m} \vdash n=m
}{%
  n\in\mathbb{N}\dsep m\in\mathbb{N}\dsep
  \NatSeg{n} \EqCard \NatSeg{m} \vdash n=m
}{%
  \DeltaRow{Mengen}{n\dsep m}
  \DeltaRow{Relationen}{<_{\mathbb{N}}}
}
\begin{tabproof}
  \proofstep{1}{n\in\mathbb{N}}{\rA}
  \proofstep{2}{m\in\mathbb{N}}{\rA}
  \proofstep{3}{\NatSeg{n} \EqCard \NatSeg{m}}{\rA}
  \proofstep{1,2}{n < m\lor m < n\lor n=m}{%
    \FormulaRefAuto{m\in\mathbb{N}\dsep n\in\mathbb{N}\vdash m<n\lor n<m\lor m=n}{1,2}}

  \proofstep{5}{n < m}{\rA}
  \proofstep{1,2,5}{\neg(\NatSeg{n} \EqCard \NatSeg{m})}{%
    \FormulaRefAuto{m\in\mathbb{N}\dsep n\in\mathbb{N}\dsep m < n \vdash \neg(\NatSeg{m} \EqCard \NatSeg{n})}{1,2,5}}
  \proofstep{1,2,3,5}{n=m}{\FormulaRefAuto{P,\,\neg P \vdash Q}{3,6}}

  \proofstep{8}{m < n}{\rA}
  \proofstep{3}{\NatSeg{m} \EqCard \NatSeg{n}}{\FormulaRefAuto{A \EqCard B \vdash B \EqCard A}{3}}
  \proofstep{2,1,8}{\neg(\NatSeg{m} \EqCard \NatSeg{n})}{%
    \FormulaRefAuto{m\in\mathbb{N}\dsep n\in\mathbb{N}\dsep m < n \vdash \neg(\NatSeg{m} \EqCard \NatSeg{n})}{2,1,8}}
  \proofstep{1,2,3,8}{n=m}{\FormulaRefAuto{P,\,\neg P \vdash Q}{9,10}}

  \proofstep{12}{n=m}{\rA}

  \proofstep{1,2,3}{n=m}{%
    \FormulaRefAuto{P_1\lor P_2\lor P_3\dsep [P_1]\vdots R\dsep [P_2]\vdots R\dsep [P_3]\vdots R\vdash R}{4,5,7,8,11,12,12}}
\end{tabproof}

\FormulaThmDeltaK[Eindeutigkeit der endlichen Kardinalzahl]{%
  \begin{aligned}[t]
    &n\in\mathbb{N}\dsep m\in\mathbb{N}\dsep{}\\
    &\NatSeg{n} \EqCard M\dsep \NatSeg{m} \EqCard M \vdash n=m
  \end{aligned}
}{FiniteCardWitnessUnique}{%
  \DeltaRow{Mengen}{M\dsep n\dsep m}
  \DeltaRow{Relationen}{<_{\mathbb{N}}}
}
\begin{tabproof}
  \proofstep{1}{n\in\mathbb{N}}{\rA}
  \proofstep{2}{m\in\mathbb{N}}{\rA}
  \proofstep{3}{\NatSeg{n} \EqCard M}{\rA}
  \proofstep{4}{\NatSeg{m} \EqCard M}{\rA}

  \proofstep{4}{M \EqCard \NatSeg{m}}{\FormulaRefAuto{A \EqCard B \vdash B \EqCard A}{4}}
  \proofstep{3,4}{\NatSeg{n} \EqCard \NatSeg{m}}{%
    \FormulaRefAuto{A \EqCard B \dsep B \EqCard C \vdash A \EqCard C}{3,5}}
  \proofstep{1,2,3,4}{n=m}{%
    \FormulaRefAuto{n\in\mathbb{N}\dsep m\in\mathbb{N}\dsep \NatSeg{n} \EqCard \NatSeg{m} \vdash n=m}{1,2,6}}
\end{tabproof}
\subsection{Definition der Kardinalzahl}

\FormulaThmDeltaK[Wohldefiniertheit der endlichen Kardinalzahl]{%
  \Finite{M}\vdash \exists! n\,\bigl(n\in\mathbb{N}\land \NatSeg{n} \EqCard M\bigr)
}{FiniteCardWellDefined}{%
  \DeltaRow{Mengen}{M\dsep n\dsep m}
}
\begin{tabproof}
  \proofstep{1}{\Finite{M}}{\rA}
  \proofstep{1}{\exists n\,\bigl(n\in\mathbb{N}\land \NatSeg{n} \EqCard M\bigr)}{%
    \FormulaRefAuto{FiniteEqCardEquiv}[thm]{1}}

  \proofstep{3}{n\in\mathbb{N}\land \NatSeg{n} \EqCard M}{\rA}
  \proofstep{3}{n\in\mathbb{N}}{\rAEa{3}}
  \proofstep{3}{\NatSeg{n} \EqCard M}{\rAEb{3}}

  \proofstep{6}{m\in\mathbb{N}\land \NatSeg{m} \EqCard M}{\rA}
  \proofstep{6}{m\in\mathbb{N}}{\rAEa{6}}
  \proofstep{6}{\NatSeg{m} \EqCard M}{\rAEb{6}}

  \proofstep{3,6}{n=m}{%
    \FormulaRefAuto{FiniteCardWitnessUnique}[thm]{4,7,5,8}}

  \proofstep{1}{\exists! n\,\bigl(n\in\mathbb{N}\land \NatSeg{n} \EqCard M\bigr)}{\UEI{2,3,6,9}}
\end{tabproof}

\FormulaDefDeltaK[Kardinalzahl einer endlichen Menge]{%
  \Finite{M}\vdash
  \tCard(M) \coloneqq
  \iota n\,\bigl(n\in\mathbb{N}\land \NatSeg{n} \EqCard M\bigr)
}{FiniteCardinality}{%
  \DeltaRow{Mengen}{M\dsep n}
  \DeltaRow{Termsymbol}{\tCard}
}

\subsection{Vergleich endlicher Mengen}

\FormulaThmDeltaK[Kardinalzahl einer endlichen Menge ist natürlich]{%
  \Finite{M}\vdash \tCard(M)\in\mathbb{N}
}{FiniteCardNatural}{%
  \DeltaRow{Mengen}{M}
}
\begin{tabproof}
  \proofstep{1}{\Finite{M}}{\rA}
  \proofstep{1}{\tCard(M)\in\mathbb{N}\land \NatSeg{\tCard(M)}\EqCard M}{%
    \FormulaRefAuto{FiniteCardinality}[def]{1}}
  \proofstep{1}{\tCard(M)\in\mathbb{N}}{\rAEa{2}}
\end{tabproof}

\FormulaThmDeltaK[Kardinalzahl einer endlichen Menge zählt die Menge]{%
  \Finite{M}\vdash \NatSeg{\tCard(M)}\EqCard M
}{FiniteCardCounts}{%
  \DeltaRow{Mengen}{M}
}
\begin{tabproof}
  \proofstep{1}{\Finite{M}}{\rA}
  \proofstep{1}{\tCard(M)\in\mathbb{N}\land \NatSeg{\tCard(M)}\EqCard M}{%
    \FormulaRefAuto{FiniteCardinality}[def]{1}}
  \proofstep{1}{\NatSeg{\tCard(M)}\EqCard M}{\rAEb{2}}
\end{tabproof}

\FormulaThmDelta[Kardinalzahlenvergleich liefert Inklusionsinjektion]{%
  \Finite{M}\dsep \Finite{N}\dsep \tCard(M)\leq \tCard(N)
  \vdash
  \iota_{\NatSeg{\tCard(M)},\NatSeg{\tCard(N)}}\colon
  \NatSeg{\tCard(M)}\inj \NatSeg{\tCard(N)}
}{%
  \DeltaRow{Mengen}{M\dsep N}
}
\begin{tabproof}
  \proofstep{1}{\Finite{M}}{\rA}
  \proofstep{2}{\Finite{N}}{\rA}
  \proofstep{3}{\tCard(M)\leq \tCard(N)}{\rA}
  \proofstep{1}{\tCard(M)\in\mathbb{N}}{%
    \FormulaRefAuto{\Finite{M}\vdash \tCard(M)\in\mathbb{N}}{1}}
  \proofstep{2}{\tCard(N)\in\mathbb{N}}{%
    \FormulaRefAuto{\Finite{M}\vdash \tCard(M)\in\mathbb{N}}{2}}
  \proofstep{1,2,3}{\NatSeg{\tCard(M)}\subseteq \NatSeg{\tCard(N)}}{%
    \FormulaRefAuto{PeanoNatSegLeqSubset}[thm]{4,5,3}}
  \proofstep{1,2,3}{\iota_{\NatSeg{\tCard(M)},\NatSeg{\tCard(N)}}\colon
    \NatSeg{\tCard(M)}\inj \NatSeg{\tCard(N)}}{%
    \FormulaRefAuto{A\subseteq B\vdash \iota_{A,B}\colon A\inj B}{6}}
\end{tabproof}

\FormulaThmDelta[Injektion aus kleinerer endlicher Kardinalzahl]{%
  \Finite{M}\dsep \Finite{N}\dsep \tCard(M)\leq \tCard(N)
  \vdash \exists F\colon M\inj N
}{%
  \DeltaRow{Mengen}{M\dsep N}
  \DeltaRow{Funktionen}{F}
}
\begin{tabproof}
  \proofstep{1}{\Finite{M}}{\rA}
  \proofstep{2}{\Finite{N}}{\rA}
  \proofstep{3}{\tCard(M)\leq \tCard(N)}{\rA}
  \proofstep{1}{\NatSeg{\tCard(M)}\EqCard M}{%
    \FormulaRefAuto{FiniteCardCounts}[thm]{1}}
  \proofstep{2}{\NatSeg{\tCard(N)}\EqCard N}{%
    \FormulaRefAuto{FiniteCardCounts}[thm]{2}}
  \proofstep{1,2,3}{\iota_{\NatSeg{\tCard(M)},\NatSeg{\tCard(N)}}\colon
    \NatSeg{\tCard(M)}\inj \NatSeg{\tCard(N)}}{%
    \FormulaRefAuto{\Finite{M}\dsep \Finite{N}\dsep \tCard(M)\leq \tCard(N)
    \vdash
    \iota_{\NatSeg{\tCard(M)},\NatSeg{\tCard(N)}}\colon
    \NatSeg{\tCard(M)}\inj \NatSeg{\tCard(N)}}{1,2,3}}
  \proofstep{1,2,3}{\exists F\colon M\inj N}{%
    \FormulaRefAuto{A \EqCard M\dsep B \EqCard N\dsep F\colon A\inj B
    \vdash \exists G\colon M\inj N}{4,5,6}}
\end{tabproof}

\FormulaThmDelta[Vergleich endlicher Mengen durch Injektionen]{%
  \Finite{M}\dsep \Finite{N}\dsep
  \neg\exists F\colon M\inj N
  \vdash \exists G\colon N\inj M
}{%
  \DeltaRow{Mengen}{M\dsep N}
  \DeltaRow{Funktionen}{F\dsep G}
}
\begin{tabproof}
  \proofstep{1}{\Finite{M}}{\rA}
  \proofstep{2}{\Finite{N}}{\rA}
  \proofstep{3}{\neg\exists F\colon M\inj N}{\rA}
  \proofstep{1}{\tCard(M)\in\mathbb{N}}{%
    \FormulaRefAuto{\Finite{M}\vdash \tCard(M)\in\mathbb{N}}{1}}
  \proofstep{2}{\tCard(N)\in\mathbb{N}}{%
    \FormulaRefAuto{\Finite{M}\vdash \tCard(M)\in\mathbb{N}}{2}}
  \proofstep{1,2}{\tCard(M)\leq \tCard(N)\lor \tCard(N)\leq \tCard(M)}{%
    \FormulaRefAuto{PeanoOrderComparison}[thm]{4,5}}

  \proofstep{7}{\tCard(M)\leq \tCard(N)}{\rA}
  \proofstep{1,2,7}{\exists F\colon M\inj N}{%
    \FormulaRefAuto{\Finite{M}\dsep \Finite{N}\dsep \tCard(M)\leq \tCard(N)\vdash \exists F\colon M\inj N}{1,2,7}}
  \proofstep{1,2,3,7}{\exists G\colon N\inj M}{%
    \FormulaRefAuto{P,\,\neg P \vdash Q}{8,3}}

  \proofstep{10}{\tCard(N)\leq \tCard(M)}{\rA}
  \proofstep{1,2,10}{\exists G\colon N\inj M}{%
    \FormulaRefAuto{\Finite{M}\dsep \Finite{N}\dsep \tCard(M)\leq \tCard(N)\vdash \exists F\colon M\inj N}{2,1,10}}

  \proofstep{1,2,3}{\exists G\colon N\inj M}{\rOE{6,7,9,10,11}}
\end{tabproof}

Die entsprechende Aussage für Surjektionen braucht Nichtleerheit auf beiden
Seiten: Aus einer Injektion \(A\to B\) erhalten wir hier eine Surjektion
\(B\to A\) nur unter \(A\neq\varnothing\).

\FormulaThmDelta[Vergleich nichtleerer endlicher Mengen durch Surjektionen]{%
  \Finite{M}\dsep \Finite{N}\dsep M\neq\varnothing\dsep N\neq\varnothing\dsep
  \neg\exists F\colon M\sur N
  \vdash \exists G\colon N\sur M
}{%
  \DeltaRow{Mengen}{M\dsep N}
  \DeltaRow{Funktionen}{F\dsep G\dsep H\dsep K}
}
\begin{tabproof}
  \proofstep{1}{\Finite{M}}{\rA}
  \proofstep{2}{\Finite{N}}{\rA}
  \proofstep{3}{M\neq\varnothing}{\rA}
  \proofstep{4}{N\neq\varnothing}{\rA}
  \proofstep{5}{\neg\exists F\colon M\sur N}{\rA}

  \proofstep{6}{\exists H\colon N\inj M}{\rA}
  \proofstep{7}{H\colon N\inj M}{\rA}
  \proofstep{4,7}{\exists F\colon M\sur N}{%
    \FormulaRefAuto{F\colon A\inj B \dsep A\neq \varnothing \vdash \exists G\colon B\sur A}{7,4}}
  \proofstep{4,5,7}{\bot}{\rBI{8,5}}
  \proofstep{4,5,6}{\bot}{\rEE{6,7,9}}
  \proofstep{4,5}{\neg\exists H\colon N\inj M}{\rCI{6,10}}

  \proofstep{1,2,4,5}{\exists K\colon M\inj N}{%
    \FormulaRefAuto{\Finite{M}\dsep \Finite{N}\dsep \neg\exists F\colon M\inj N\vdash \exists G\colon N\inj M}{2,1,11}}
  \proofstep{13}{K\colon M\inj N}{\rA}
  \proofstep{3,13}{\exists G\colon N\sur M}{%
    \FormulaRefAuto{F\colon A\inj B \dsep A\neq \varnothing \vdash \exists G\colon B\sur A}{13,3}}
  \proofstep{1,2,3,4,5}{\exists G\colon N\sur M}{\rEE{12,13,14}}
\end{tabproof}

\chapter{Von-Neumann-Charakterisierung der Endlichkeit}

In diesem Abschnitt ist die natürliche Zahl \(n\) selbst das von-Neumann-
Anfangssegment.  Die Charakterisierung wird in beiden Richtungen unmittelbar
aus den drei Endlichkeitsaxiomen und den bereits bewiesenen Sätzen über
natürliche Zahlen und Gleichmächtigkeit hergeleitet.

\FormulaThmDeltaK[Von-Neumann-Zeuge für endliche Mengen]{%
  \Finite{M}\vdash
  \exists n\in\mathbb{N}\,\bigl(n\EqCard M\bigr)%
}{B03FiniteVonNeumannEqCardForward}{%
  \DeltaRow{Mengen}{A\dsep M\dsep a\dsep n\dsep m}%
}
\begin{tabproofsplitwide}
  \proofpartwideindR[Induktionsanfang]{%
    \exists n\in\mathbb{N}\,\bigl(n\EqCard\varnothing\bigr)%
  }{%
    \exists n\in\mathbb{N}\,\bigl(n\EqCard\varnothing\bigr)%
  }[B03FiniteVonNeumannEqCardForwardBase]
    \proofstepwidestar[]{\varnothing\in\mathbb{N}}{%
      \FormulaRefAuto{\varnothing\in\mathbb{N}}}
    \proofstepwidestar[]{\varnothing\EqCard\varnothing}{%
      \FormulaRefAuto{EqCardReflexive}[thm]}
    \proofstepwidestar[]{%
      \varnothing\in\mathbb{N}\land
      \varnothing\EqCard\varnothing}{\rAI{1,2}}
    \proofstepwidestar[]{%
      \exists n\in\mathbb{N}\,\bigl(n\EqCard\varnothing\bigr)}{\rEI{3}}
  \closeproofpartwideindR

  \proofpartwideindR[Induktionsschritt]{%
    \begin{aligned}[t]
      &\Finite{A}\dsep
       \exists n\in\mathbb{N}\,\bigl(n\EqCard A\bigr)
       \dsep a\notin A\vdash{}\\
      &\exists m\in\mathbb{N}\,\bigl(m\EqCard A\cup\{a\}\bigr)
    \end{aligned}%
  }{%
    \Finite{A}\dsep
    \exists n\in\mathbb{N}\,\bigl(n\EqCard A\bigr)
    \dsep a\notin A\vdash
    \exists m\in\mathbb{N}\,\bigl(m\EqCard A\cup\{a\}\bigr)%
  }[B03FiniteVonNeumannEqCardForwardStep]
    \proofstepwidestar[1]{\Finite{A}}{\rA}
    \proofstepwidestar[2]{%
      \exists n\in\mathbb{N}\,\bigl(n\EqCard A\bigr)}{\rA}
    \proofstepwidestar[3]{a\notin A}{\rA}
    \proofstepwidestar[4]{n\in\mathbb{N}\land n\EqCard A}{\rA}
    \proofstepwidestar[4]{n\in\mathbb{N}}{\rAEa{4}}
    \proofstepwidestar[4]{n\EqCard A}{\rAEb{4}}
    \proofstepwidestar[]{n\notin n}{\FormulaRefAuto{a\not\in a}}
    \proofstepwidestar[3,4]{%
      n\cup\{n\}\EqCard A\cup\{a\}}{%
      \FormulaRefAuto{EqCardFreshAdjunction}[thm]{6,7,3}}
    \proofstepwidestar[4]{(n\cup\{n\})\in\mathbb{N}}{%
      \FormulaRefAuto{n\in\mathbb{N}\vdash
        (n\cup\{n\})\in\mathbb{N}}{5}}
    \proofstepwidestar[3,4]{%
      (n\cup\{n\})\in\mathbb{N}\land
      n\cup\{n\}\EqCard A\cup\{a\}}{\rAI{9,8}}
    \proofstepwidestar[3,4]{%
      \exists m\in\mathbb{N}\,\bigl(m\EqCard A\cup\{a\}\bigr)}{\rEI{10}}
    \proofstepwidestar[2,3]{%
      \exists m\in\mathbb{N}\,\bigl(m\EqCard A\cup\{a\}\bigr)}{%
      \rEE{2,4,11}}
  \closeproofpartwideindR

  \proofpartwide{Induktionsschluss}
    \proofstepwidestar[1]{\Finite{M}}{\rA}
    \proofstepwidestar[1]{%
      \exists n\in\mathbb{N}\,\bigl(n\EqCard M\bigr)}{%
      \FormulaRefAuto{FiniteInductionRule}[ax]{IA,IS,1}}
  \closeproofpartwide
\end{tabproofsplitwide}

\FormulaThmDeltaK[Einschränkung einer Nachfolgerbijektion]{%
  n\in\mathbb{N}\dsep F\colon\Succ(n)\bij M
  \vdash n\EqCard F[n]%
}{B03VonNeumannSuccRestrictionEqCard}{%
  \DeltaRow{Mengen}{M\dsep n}%
  \DeltaRow{Funktionen}{F\dsep G}%
}
\begin{tabproof}
  \proofstep{1}{n\in\mathbb{N}}{\rA}
  \proofstep{2}{F\colon\Succ(n)\bij M}{\rA}
  \proofstep{2}{F\colon\Succ(n)\to M}{\FormulaRefAuto{Bijektive Funktion}{2}}
  \proofstep{2}{F\colon\Succ(n)\inj M}{%
    \FormulaRefAuto{F\colon A\bij B\vdash F\colon A\inj B}{2}}
  \proofstep{1}{n\subseteq\Succ(n)}{%
    \FormulaRefAuto{n\in\mathbb{N}\vdash n\subseteq\Succ(n)}{1}}
  \proofstep{1,2}{F\restriction_n\colon n\inj M}{%
    \FormulaRefAuto{F\colon A\inj B\dsep C\subseteq A
      \vdash F\restriction_C\colon C\inj B}{4,5}}
  \proofstep{1,2}{%
    F\restriction_n\colon n\bij(F\restriction_n)[n]}{%
    \FormulaRefAuto{F\colon A\inj B\vdash F\colon A\bij F[A]}{6}}
  \proofstep{1,2}{(F\restriction_n)[n]=F[n]}{%
    \FormulaRefAuto{C\subseteq A\dsep F\colon A\to B
      \vdash (F\restriction_C)[C]=F[C]}{5,3}}
  \proofstep{1,2}{F\restriction_n\colon n\bij F[n]}{\rIE{8,7}}
  \proofstep{1,2}{\exists G\colon n\bij F[n]}{\rEI{9}}
  \proofstep{1,2}{n\EqCard F[n]}{%
    \FormulaRefAuto{A\EqCard B\coloneqq \exists F\colon A\bij B}{10}}
\end{tabproof}

\FormulaThmDeltaK[Das neue Bild liegt nicht im alten Bild]{%
  n\in\mathbb{N}\dsep F\colon\Succ(n)\bij M
  \vdash F(n)\notin F[n]%
}{B03VonNeumannSuccFreshImage}{%
  \DeltaRow{Mengen}{M\dsep n}%
  \DeltaRow{Funktionen}{F}%
}
\begin{tabproof}
  \proofstep{1}{n\in\mathbb{N}}{\rA}
  \proofstep{2}{F\colon\Succ(n)\bij M}{\rA}
  \proofstep{2}{F\colon\Succ(n)\inj M}{%
    \FormulaRefAuto{F\colon A\bij B\vdash F\colon A\inj B}{2}}
  \proofstep{1}{n\subseteq\Succ(n)}{%
    \FormulaRefAuto{n\in\mathbb{N}\vdash n\subseteq\Succ(n)}{1}}
  \proofstep{1}{n\in\powerset(\Succ(n))}{%
    \FormulaRefAuto{x\in\powerset(A)\eqvdash x\subseteq A}{4}}
  \proofstep{1}{n\in\Succ(n)}{%
    \rIE{%
      \FormulaRefAuto{a=b\vdash b=a}{%
        \FormulaRefAuto{n\in\mathbb{N}\vdash
          \Succ(n)=n\cup\{n\}}{1}},%
      \FormulaRefAuto{a\in A\cup\{a\}}}}
  \proofstep{7}{F(n)\in F[n]}{\rA}
  \proofstep{1,2,7}{n\in n}{%
    \FormulaRefAuto{%
      F\colon A\inj B\dsep Y\in\powerset(A)\dsep x\in A\dsep
      F(x)\in F[Y]\vdash x\in Y}{3,5,6,7}}
  \proofstep{}{n\notin n}{\FormulaRefAuto{a\not\in a}}
  \proofstep{1,2,7}{\bot}{\rBI{8,9}}
  \proofstep{1,2}{F(n)\notin F[n]}{\rCI{7,10}}
\end{tabproof}

\FormulaThmDeltaK[Bildzerlegung einer Nachfolgerbijektion]{%
  n\in\mathbb{N}\dsep F\colon\Succ(n)\bij M
  \vdash F[n]\cup\{F(n)\}=M%
}{B03VonNeumannSuccImageDecomposition}{%
  \DeltaRow{Mengen}{M\dsep n}%
  \DeltaRow{Funktionen}{F}%
}
\begin{tabproof}
  \proofstep{1}{n\in\mathbb{N}}{\rA}
  \proofstep{2}{F\colon\Succ(n)\bij M}{\rA}
  \proofstep{2}{F\colon\Succ(n)\to M}{\FormulaRefAuto{Bijektive Funktion}{2}}
  \proofstep{1}{\Succ(n)=n\cup\{n\}}{%
    \FormulaRefAuto{n\in\mathbb{N}\vdash\Succ(n)=n\cup\{n\}}{1}}
  \proofstep{1}{n\subseteq\Succ(n)}{%
    \FormulaRefAuto{n\in\mathbb{N}\vdash n\subseteq\Succ(n)}{1}}
  \proofstep{1}{n\in\Succ(n)}{%
    \rIE{%
      \FormulaRefAuto{a=b\vdash b=a}{4},%
      \FormulaRefAuto{a\in A\cup\{a\}}}}
  \proofstep{1}{\{n\}\subseteq\Succ(n)}{%
    \FormulaRefAuto{a\in A\vdash\{a\}\subseteq A}{6}}
  \proofstep{1,2}{F[n\cup\{n\}]=F[n]\cup F[\{n\}]}{%
    \FormulaRefAuto{%
      F\colon M\to N\dsep A\subseteq M\dsep B\subseteq M
      \vdash F[A\cup B]=F[A]\cup F[B]}{3,5,7}}
  \proofstep{1,2}{F[\{n\}]=\{F(n)\}}{%
    \FormulaRefAuto{x\in A\dsep F\colon A\to B
      \vdash F[\{x\}]=\{F(x)\}}{6,3}}
  \proofstep{1,2}{F[\Succ(n)]=F[n]\cup\{F(n)\}}{%
    \rIE{4,\rIE{9,8}}}
  \proofstep{2}{F[\Succ(n)]=M}{%
    \FormulaRefAuto{F\colon A\bij B\vdash F[A]=B}{2}}
  \proofstep{1,2}{F[n]\cup\{F(n)\}=M}{%
    \FormulaRefAuto{a=b\dsep a=c\vdash b=c}{10,11}}
\end{tabproof}

\FormulaThmDeltaK[Endlichkeit aus einem von-Neumann-Zeugen]{%
  n\in\mathbb{N}\dsep n\EqCard M\vdash \Finite{M}%
}{B03FiniteVonNeumannEqCardReverse}{%
  \DeltaRow{Mengen}{A\dsep M\dsep n}%
  \DeltaRow{Funktionen}{F}%
}
\begin{tabproofsplitwide}
  \proofpartwideindR[Induktionsanfang]{%
    \forall M\,\bigl((\varnothing\EqCard M)\rightarrow\Finite{M}\bigr)%
  }{%
    \forall M\,\bigl((\varnothing\EqCard M)\rightarrow\Finite{M}\bigr)%
  }[B03FiniteVonNeumannEqCardReverseBase]
    \proofstepwidestar[1]{\varnothing\EqCard M}{\rA}
    \proofstepwidestar[1]{\exists F\colon\varnothing\bij M}{%
      \FormulaRefAuto{A\EqCard B\coloneqq \exists F\colon A\bij B}{1}}
    \proofstepwidestar[3]{F\colon\varnothing\bij M}{\rA}
    \proofstepwidestar[3]{F\colon\varnothing\to M}{%
      \FormulaRefAuto{Bijektive Funktion}{3}}
    \proofstepwidestar[3]{F[\varnothing]=\varnothing}{%
      \FormulaRefAuto{F\colon A\to B\vdash F[\varnothing]=\varnothing}{4}}
    \proofstepwidestar[3]{F[\varnothing]=M}{%
      \FormulaRefAuto{F\colon A\bij B\vdash F[A]=B}{3}}
    \proofstepwidestar[3]{\varnothing=F[\varnothing]}{%
      \FormulaRefAuto{a=b\vdash b=a}{5}}
    \proofstepwidestar[3]{\varnothing=M}{%
      \FormulaRefAuto{a=b\dsep b=c\vdash a=c}{7,6}}
    \proofstepwidestar[]{\Finite{\varnothing}}{%
      \FormulaRefAuto{FiniteEmptyAxiom}[ax]}
    \proofstepwidestar[3]{\Finite{M}}{\rIE{8,9}}
    \proofstepwidestar[1]{\Finite{M}}{\rEE{2,3,10}}
    \proofstepwidestar[]{%
      (\varnothing\EqCard M)\rightarrow\Finite{M}}{\rRI{1,11}}
    \proofstepwidestar[]{%
      \forall M\,\bigl((\varnothing\EqCard M)\rightarrow\Finite{M}\bigr)}{%
      \rUI{12}}
  \closeproofpartwideindR

  \proofpartwideindR[Induktionsschritt]{%
    \begin{aligned}[t]
      &n\in\mathbb{N}\dsep
       \forall A\,\bigl((n\EqCard A)\rightarrow\Finite{A}\bigr)
       \vdash{}\\
      &\forall M\,\bigl((\Succ(n)\EqCard M)\rightarrow\Finite{M}\bigr)
    \end{aligned}%
  }{%
    n\in\mathbb{N}\dsep
    \forall A\,\bigl((n\EqCard A)\rightarrow\Finite{A}\bigr)
    \vdash
    \forall M\,\bigl((\Succ(n)\EqCard M)\rightarrow\Finite{M}\bigr)%
  }[B03FiniteVonNeumannEqCardReverseStep]
    \proofstepwidestar[1]{n\in\mathbb{N}}{\rA}
    \proofstepwidestar[2]{%
      \forall A\,\bigl((n\EqCard A)\rightarrow\Finite{A}\bigr)}{\rA}
    \proofstepwidestar[3]{\Succ(n)\EqCard M}{\rA}
    \proofstepwidestar[3]{\exists F\colon\Succ(n)\bij M}{%
      \FormulaRefAuto{A\EqCard B\coloneqq \exists F\colon A\bij B}{3}}
    \proofstepwidestar[5]{F\colon\Succ(n)\bij M}{\rA}
    \proofstepwidestar[1,5]{n\EqCard F[n]}{%
      \FormulaRefAuto{B03VonNeumannSuccRestrictionEqCard}[thm]{1,5}}
    \proofstepwidestar[2]{%
      (n\EqCard F[n])\rightarrow\Finite{F[n]}}{\rUE{2}}
    \proofstepwidestar[1,2,5]{\Finite{F[n]}}{\rRE{7,6}}
    \proofstepwidestar[1,5]{F(n)\notin F[n]}{%
      \FormulaRefAuto{B03VonNeumannSuccFreshImage}[thm]{1,5}}
    \proofstepwidestar[1,2,5]{\Finite{F[n]\cup\{F(n)\}}}{%
      \FormulaRefAuto{FiniteAdjunctionAxiom}[ax]{9,8}}
    \proofstepwidestar[1,5]{F[n]\cup\{F(n)\}=M}{%
      \FormulaRefAuto{B03VonNeumannSuccImageDecomposition}[thm]{1,5}}
    \proofstepwidestar[1,2,5]{\Finite{M}}{\rIE{11,10}}
    \proofstepwidestar[1,2,3]{\Finite{M}}{\rEE{4,5,12}}
    \proofstepwidestar[1,2]{%
      (\Succ(n)\EqCard M)\rightarrow\Finite{M}}{\rRI{3,13}}
    \proofstepwidestar[1,2]{%
      \forall M\,\bigl((\Succ(n)\EqCard M)\rightarrow\Finite{M}\bigr)}{%
      \rUI{14}}
  \closeproofpartwideindR

  \proofpartwide{Induktionsschluss}
    \proofstepwidestar[1]{n\in\mathbb{N}}{\rA}
    \proofstepwidestar[2]{n\EqCard M}{\rA}
    \proofstepwidestar[1]{%
      \forall A\,\bigl((n\EqCard A)\rightarrow\Finite{A}\bigr)}{%
      \FormulaRefAuto{%
        P(\varnothing)\dsep
        [x\in\mathbb{N},P(x)]\vdots P(\Succ(x))\dsep
        n\in\mathbb{N}\vdash P(n)}{IA,IS,1}}
    \proofstepwidestar[1]{%
      (n\EqCard M)\rightarrow\Finite{M}}{\rUE{3}}
    \proofstepwidestar[1,2]{\Finite{M}}{\rRE{4,2}}
  \closeproofpartwide
\end{tabproofsplitwide}

\FormulaThmDeltaK[Charakterisierung der Endlichkeit durch von-Neumann-Zahlen]{%
  \Finite{M}\eqvdash
  \exists n\in\mathbb{N}\,\bigl(n\EqCard M\bigr)%
}{FiniteVonNeumannEqCardEquiv}{%
  \DeltaRow{Mengen}{M\dsep n}%
}
\begin{tabproofsplit}
  \proofpart{\(\vdash\)}
    \proofstep{1}{\Finite{M}}{\rA}
    \proofstep{1}{\exists n\in\mathbb{N}\,\bigl(n\EqCard M\bigr)}{%
      \FormulaRefAuto{B03FiniteVonNeumannEqCardForward}[thm]{1}}
  \closeproofpart

  \proofpart{\(\dashv\)}
    \proofstep{1}{\exists n\in\mathbb{N}\,\bigl(n\EqCard M\bigr)}{\rA}
    \proofstep{2}{n\in\mathbb{N}\land n\EqCard M}{\rA}
    \proofstep{2}{n\in\mathbb{N}}{\rAEa{2}}
    \proofstep{2}{n\EqCard M}{\rAEb{2}}
    \proofstep{2}{\Finite{M}}{%
      \FormulaRefAuto{B03FiniteVonNeumannEqCardReverse}[thm]{3,4}}
    \proofstep{1}{\Finite{M}}{\rEE{1,2,5}}
  \closeproofpart
\end{tabproofsplit}

% Die frühere definitionelle Fassung wird nicht mehr registriert.  Der
% folgende alte Herleitungsblock bleibt nur als historische Arbeitsfassung
% im Quelltext und ist vollständig deaktiviert.

\end{document}
